% !TEX program = xelatex
\documentclass[11pt,a4paper]{ctexart}
% 零基础讲义公共导言区
\usepackage{amsmath,amssymb,amsthm,mathtools}
\usepackage{geometry}
\geometry{left=2.0cm,right=2.0cm,top=2.2cm,bottom=2.0cm}
\usepackage{booktabs}
\usepackage{array}
\usepackage{enumitem}
\usepackage{tcolorbox}
\tcbuselibrary{skins,breakable}
\usepackage{xcolor}
\usepackage{fancyhdr}
\usepackage{titlesec}
\usepackage{multirow}
\usepackage{longtable}

% ---------- 颜色 ----------
\definecolor{mainblue}{RGB}{25,70,130}
\definecolor{acc1}{RGB}{176,58,46}      % 定义/公式
\definecolor{acc2}{RGB}{20,105,85}      % 直觉/类比
\definecolor{acc3}{RGB}{120,80,150}     % 你可能想问
\definecolor{acc4}{RGB}{160,110,20}     % 手把手算
\definecolor{acc5}{RGB}{40,90,160}      % 考试怎么考
\definecolor{soft}{RGB}{246,247,250}

% ---------- 数学宏 ----------
\newcommand{\Hx}[1]{H\!\left(#1\right)}
\newcommand{\Hc}[1]{H_{c}\!\left(#1\right)}
\newcommand{\I}[1]{I\!\left(#1\right)}
\newcommand{\lb}{\log_{2}}
\DeclareMathOperator{\E}{\mathbb{E}}
\newcommand{\bit}{\,\text{bit}}
\newcommand{\p}[1]{p\!\left(#1\right)}

% ---------- 五种盒子 ----------
\newtcolorbox{defbox}[1]{enhanced,breakable,
  colback=soft,colframe=acc1,boxrule=0.7pt,arc=2pt,
  left=6pt,right=6pt,top=5pt,bottom=5pt,
  title={\bfseries #1},fonttitle=\color{white},colbacktitle=acc1,
  attach boxed title to top left={xshift=6pt,yshift=-2.6pt},
  boxed title style={arc=1pt,boxrule=0pt}}

\newtcolorbox{idea}[1]{enhanced,breakable,
  colback=white,colframe=acc2,boxrule=0.7pt,arc=2pt,
  left=6pt,right=6pt,top=5pt,bottom=5pt,
  title={\bfseries #1},fonttitle=\color{white},colbacktitle=acc2,
  attach boxed title to top left={xshift=6pt,yshift=-2.6pt},
  boxed title style={arc=1pt,boxrule=0pt}}

\newtcolorbox{qa}[1]{enhanced,breakable,
  colback=white,colframe=acc3,boxrule=0.7pt,arc=2pt,
  left=6pt,right=6pt,top=5pt,bottom=5pt,
  title={\bfseries #1},fonttitle=\color{white},colbacktitle=acc3,
  attach boxed title to top left={xshift=6pt,yshift=-2.6pt},
  boxed title style={arc=1pt,boxrule=0pt}}

\newtcolorbox{calc}[1]{enhanced,breakable,
  colback=white,colframe=acc4,boxrule=0.7pt,arc=2pt,
  left=6pt,right=6pt,top=5pt,bottom=5pt,
  title={\bfseries #1},fonttitle=\color{white},colbacktitle=acc4,
  attach boxed title to top left={xshift=6pt,yshift=-2.6pt},
  boxed title style={arc=1pt,boxrule=0pt}}

\newtcolorbox{exam}[1]{enhanced,breakable,
  colback=white,colframe=acc5,boxrule=0.7pt,arc=2pt,
  left=6pt,right=6pt,top=5pt,bottom=5pt,
  title={\bfseries #1},fonttitle=\color{white},colbacktitle=acc5,
  attach boxed title to top left={xshift=6pt,yshift=-2.6pt},
  boxed title style={arc=1pt,boxrule=0pt}}

% 兼容旧环境名
\newenvironment{thmbox}[1]{\begin{defbox}{#1}}{\end{defbox}}
\newenvironment{keypoint}[1]{\begin{idea}{#1}}{\end{idea}}
\newenvironment{warning}[1]{\begin{qa}{#1}}{\end{qa}}
\newenvironment{example}[1]{\begin{calc}{#1}}{\end{calc}}
\newenvironment{problem}[1]{\begin{defbox}{#1}}{\end{defbox}}

% ---------- 标题格式 ----------
\titleformat{\section}{\Large\bfseries\color{acc1}}{第\,\thesection\,讲}{0.6em}{}
\titleformat{\subsection}{\large\bfseries\color{mainblue}}{\thesubsection}{0.5em}{}
\titleformat{\subsubsection}{\normalsize\bfseries\color{acc2}}{\thesubsubsection}{0.4em}{}

% ---------- 页眉页脚 ----------
\pagestyle{fancy}
\fancyhf{}
\fancyhead[L]{\small 信息论基础\ \textbar\ 零基础讲义}
\fancyhead[R]{\small\leftmark}
\fancyfoot[C]{\small\thepage}
\renewcommand{\headrulewidth}{0.4pt}

\allowdisplaybreaks
\setlength{\parskip}{2.5pt}
\linespread{1.08}

% ---------- 特殊符号兜底（避免缺字）----------
\usepackage{newunicodechar}
\newunicodechar{✓}{\ensuremath{\checkmark}}
\newunicodechar{✗}{\ensuremath{\times}}
\xeCJKDeclareCharClass{CJK}{"2460 -> "24FF}   % ① ② ③ 等圈号交给中文字体
\newunicodechar{→}{\ensuremath{\rightarrow}}
\newunicodechar{≥}{\ensuremath{\ge}}
\newunicodechar{≤}{\ensuremath{\le}}
\newunicodechar{≠}{\ensuremath{\ne}}
\newunicodechar{×}{\ensuremath{\times}}
\newunicodechar{≈}{\ensuremath{\approx}}
\newunicodechar{∞}{\ensuremath{\infty}}
\newunicodechar{∈}{\ensuremath{\in}}
\newunicodechar{⊕}{\ensuremath{\oplus}}
\newunicodechar{∑}{\ensuremath{\sum}}
\newunicodechar{∫}{\ensuremath{\int}}
\newunicodechar{√}{\ensuremath{\surd}}
\newunicodechar{σ}{\ensuremath{\sigma}}
\newunicodechar{π}{\ensuremath{\pi}}
\newunicodechar{λ}{\ensuremath{\lambda}}
\newunicodechar{μ}{\ensuremath{\mu}}
\newunicodechar{Δ}{\ensuremath{\Delta}}
\newunicodechar{α}{\ensuremath{\alpha}}
\newunicodechar{β}{\ensuremath{\beta}}
\newunicodechar{ε}{\ensuremath{\varepsilon}}


\begin{document}

\begin{titlepage}
\centering
\vspace*{2.0cm}
{\Huge\bfseries\color{acc1} 信息论基础\ \ 零基础讲义}\\[1.0em]
{\Large 一节课没听也能看懂的版本}\\[2.5em]
\rule{0.7\textwidth}{0.8pt}\\[1.5em]
{\large 从“概率 $p(x)$ 是什么”讲到“汉明码怎么构造”}\\[0.5em]
{\large 每个符号都翻译成人话\ \textbullet\ 每个公式都配手算例子}\\[0.5em]
{\large 每讲末尾附「考试怎么考」与自测题}\\[1.5em]
\rule{0.7\textwidth}{0.8pt}\\[3.5em]
\begin{minipage}{0.86\textwidth}
\begin{tcolorbox}[colback=soft,colframe=acc1,boxrule=0.7pt,arc=2pt]
\textbf{怎么用这本讲义（重要）}\\[4pt]
\begin{enumerate}[leftmargin=1.5em,itemsep=2pt,topsep=2pt]
  \item \textbf{务必按顺序读}第1--4讲，它们只讲三件事：$p(x)$ 是什么、$I(x)$ 是什么、$H(X)$ 是什么。这三件事不透，后面全是天书。
  \item 读的时候\textbf{手里拿支笔}，看到例题就遮住答案自己算一遍——信息论不背公式是靠算熟的。
  \item 每讲末尾的橙色框「\textbf{考试怎么考}」是得分清单，考前只看这些框也能捞回大量分数。
  \item 紫色框「\textbf{你可能想问}」是我预先替你问的问题，别跳过，那都是你真正会卡住的地方。
\end{enumerate}
\end{tcolorbox}
\end{minipage}
\vfill
{\large 配套材料：《作业1-超详细版》（第2章作业逐步讲解）}\\[3cm]
\end{titlepage}

\tableofcontents
\newpage

\section{先补地基：这门课到底在算什么}

\begin{idea}{一句话说明这门课在干什么}
信息论就是在回答一个问题：\textbf{“一条消息值多少信息？”}而且它要给这个“多少”一个能算的数字，单位叫 bit。
整本书所有公式，都是围绕“怎么算这个数字、及这个数字有什么用”展开的。
\end{idea}

\begin{qa}{我完全没听过课，现在最该先弄懂什么？}
只要弄懂下面这条“主线”，剩下都是它的变形：
\[
\underbrace{p(x)}_{\text{这件事发生的概率}}
\ \longrightarrow\
\underbrace{I(x)=-\lb p(x)}_{\text{这件事发生告诉了你多少信息}}
\ \longrightarrow\
\underbrace{\Hx{X}=\sum_x p(x)I(x)}_{\text{平均每件事告诉你的信息}}
\]
说得再白一点：
\begin{enumerate}[leftmargin=1.6em,itemsep=1pt,topsep=2pt]
  \item $p(x)$：这件事有\textbf{多可能}发生（0 到 1 之间的一个数）。
  \item $I(x)$：这件事\textbf{一旦告诉你}，你获得了多少信息（越好猜的信息越少，越意外的信息越多）。
  \item $\Hx{X}$：把上面每条信息的“收获”按它发生的可能性\textbf{加权平均}，就是“平均每次能听到多少信息”。
\end{enumerate}
本讲先把 $p(x)$ 这一层的地基打好；下一讲再说 $I(x)$，再下一讲说 $\Hx{X}$。
\end{qa}

\subsection{概率 $p(x)$：可能性的度量}

\subsubsection*{随机试验与样本点}

先看几个熟悉的例子：

\begin{center}
\renewcommand{\arraystretch}{1.35}
\begin{tabular}{l|l|l}
\toprule
试验 & 所有可能结果（样本点） & 每个结果的概率\\
\midrule
掷一颗骰子 & $1,2,3,4,5,6$ & 各 $1/6$\\
抛一枚公平硬币 & 正面、反面 & 各 $1/2$\\
明天天气 & 晴、阴、雨 & 例如 $0.5,\ 0.25,\ 0.25$\\
从 10 个数字里随机发一个 & $0,1,2,\dots,9$ & 各 $1/10$\\
\bottomrule
\end{tabular}
\end{center}

\textbf{概率 $p(x)$ 就是“结果 $x$ 出现的可能性”}，满足两条铁律：

\begin{defbox}{概率的两条铁律（做题时用来检查自己有没有算错）}
\begin{enumerate}[leftmargin=1.6em,itemsep=1pt,topsep=2pt]
  \item \textbf{非负}：$0\le p(x)\le 1$，即每个概率都不小于 0、不大于 1；
  \item \textbf{归一}：所有可能结果的概率\textbf{加起来必须正好等于 1}，即
        $\displaystyle\sum_{x}p(x)=1$。
\end{enumerate}
例：$p(1)=p(2)=\cdots=p(6)=\tfrac16$，六个 $\tfrac16$ 加起来是 $1$，合格。
如果你算出一组概率加起来是 $0.9$ 或 $1.1$，那一定算错了——这是考场上最快发现错误的方法。
\end{defbox}

\subsubsection*{随机变量 $X$：把结果贴个数字标签}

\textbf{随机变量只是一个“把结果翻译成数字”的标签}，不是真的要解方程。例如：
\begin{itemize}[leftmargin=1.6em,itemsep=1pt,topsep=2pt]
  \item 掷骰子，$X$ = 朝上的点数，则 $X$ 取值 $1,2,\dots,6$；
  \item 抛硬币，定义 $X=0$ 表示“反面”、$X=1$ 表示“正面”，则 $X$ 取值 $0,1$；
  \item 掷两颗骰子，$X$ = 点数之和，则 $X$ 取值 $2,3,\dots,12$（注意不是每个值等可能！）。
\end{itemize}
大写的 $X$ 表示“这个还没发生的随机结果”，小写的 $x$ 表示“其中某一个具体取值”。所以
$p(x)=P(X=x)$ 就是“$X$ 恰好取到 $x$ 的概率”。

\begin{calc}{手把手算：两颗骰子之和 $Y$ 的分布表}
两颗骰子独立，各自 $1\sim6$ 等可能，共 $6\times 6=36$ 种等可能情况。数一数“和为 $y$”有几种情况（记为 $c_y$）：
\begin{center}
\renewcommand{\arraystretch}{1.2}
\begin{tabular}{c|ccccccccccc}
\toprule
和 $y$ & 2 & 3 & 4 & 5 & 6 & 7 & 8 & 9 & 10 & 11 & 12\\
\midrule
情况数 $c_y$ & 1 & 2 & 3 & 4 & 5 & 6 & 5 & 4 & 3 & 2 & 1\\
概率 $p(y)=c_y/36$ & $\frac1{36}$ & $\frac2{36}$ & $\frac3{36}$ & $\frac4{36}$ & $\frac5{36}$ & $\frac6{36}$ & $\frac5{36}$ & $\frac4{36}$ & $\frac3{36}$ & $\frac2{36}$ & $\frac1{36}$\\
\bottomrule
\end{tabular}
\end{center}
检查归一：$1+2+3+4+5+6+5+4+3+2+1=36$，而 $36\times\frac{1}{36}=1$，合格。

\emph{读法示范}：$p(7)=\frac{6}{36}=\frac16\approx0.167$ 是最大的一个（因为凑成 7 的方式最多：$1+6,2+5,3+4,4+3,5+2,6+1$），而 $p(2)=\frac{1}{36}$ 是最小的。这说明\textbf{“和是 7”比“和是 2”好猜得多}——这一点下一讲会变成具体的 bit 数。
\end{calc}

\subsection{联合分布、边缘分布、条件概率}

这是本章最容易混的三样东西，但区别其实就是“看表格的角度不同”。

\subsubsection*{联合分布 $p(x,y)$}

把两个随机变量放在一起：$p(x,y)=P(X=x\ \text{且}\ Y=y)$，意思是“$X$ 取 $x$ \textbf{同时} $Y$ 取 $y$ 的概率”。

\begin{calc}{一张联合分布表长什么样（这张表后面作业里要用）}
\begin{center}
\renewcommand{\arraystretch}{1.3}
\begin{tabular}{c|cc}
\toprule
 & $Y=0$ & $Y=1$\\
\midrule
$X=0$ & $1/3$ & $1/3$\\
$X=1$ & $1/3$ & $0$\\
\bottomrule
\end{tabular}
\end{center}
读法：$p(0,0)=\frac13$、$p(0,1)=\frac13$、$p(1,0)=\frac13$、$p(1,1)=0$。
四格加起来 $\frac13+\frac13+\frac13+0=1$，合格。

注意 $p(1,1)=0$ 是很有信息量的一句话：它表示“\textbf{不可能同时出现 $X=1$ 和 $Y=1$}”，即一旦 $X=1$，$Y$ 就必定是 0。这种“确定性关联”是后面算条件熵的关键。
\end{calc}

\subsubsection*{边缘分布 $p(x)$：把表的一行（或一列）加起来}

只知道 $X$ 的分布、不管 $Y$ 时，就把同一行的数加起来：
\[
p(x)=\sum_{y}p(x,y).
\]
上面那张表：$p(X=0)=\frac13+\frac13=\frac23$，$p(X=1)=\frac13+0=\frac13$；同理 $p(Y=0)=\frac23$，$p(Y=1)=\frac13$。

\begin{idea}{“边缘”这个词怎么理解}
把所有行的数加起来，写在表格最右边一列——就像写在纸的\textbf{边缘}上，所以叫边缘分布。中文里也叫“边际分布”。
\end{idea}

\subsubsection*{条件概率 $p(x|y)$：已经知道 $Y$ 之后的概率}

$p(x|y)$ 读作“\textbf{在 $Y=y$ 的条件下，$X=x$ 的概率}”。算法只有一个：
\[
p(x|y)=\frac{p(x,y)}{p(y)} .
\]
\begin{calc}{手把手算：用上面那张表求条件概率}
已知 $Y=0$：$p(Y=0)=\frac23$，于是
\[
p(X=0|Y=0)=\frac{p(0,0)}{p(Y=0)}=\frac{1/3}{2/3}=\frac12,\qquad
p(X=1|Y=0)=\frac{p(1,0)}{p(Y=0)}=\frac{1/3}{2/3}=\frac12 .
\]
也就是“知道 $Y=0$ 之后，$X$ 是 0 还是 1 各占一半，完全猜不准”。

已知 $Y=1$：$p(Y=1)=\frac13$，于是
\[
p(X=0|Y=1)=\frac{p(0,1)}{p(Y=1)}=\frac{1/3}{1/3}=1,\qquad
p(X=1|Y=1)=\frac{p(1,1)}{p(Y=1)}=\frac{0}{1/3}=0 .
\]
也就是“知道 $Y=1$ 之后，$X$ 必定是 0，一点不用猜了”。

\emph{体会一下}：同样是问 $X$，知道 $Y=0$ 时你还得猜（50\% 对 50\%），知道 $Y=1$ 时你完全确定。\textbf{“还要不要猜、要猜多少”就是下一讲要量化的东西。}
\end{calc}

\begin{defbox}{独立性：这门课反复用到的“最好情况”}
若对一切 $x,y$ 都有
\[
p(x,y)=p(x)\,p(y),
\]
就称 $X$ 与 $Y$ \textbf{相互独立}，意思就是“$Y$ 的信息对猜 $X$ 毫无帮助”。
两个等价说法（后面会证明，可以先记着）：
\begin{itemize}[leftmargin=1.6em,itemsep=1pt,topsep=1pt]
  \item 独立 $\iff$ $p(x|y)=p(x)$（知道 $Y$ 后 $X$ 的概率一点没变）；
  \item 独立 $\iff$ $\I{X;Y}=0$（互信息为零，第5讲讲）。
\end{itemize}
上面那张表\textbf{不独立}：因为 $p(0,0)=\frac13$ 而 $p(X=0)p(Y=0)=\frac23\times\frac23=\frac49\ne\frac13$。
\end{defbox}

\subsection{期望 $\E[\cdot]$：其实就是“加权平均”}

这是全书出现频率最高的运算，但它只是“按概率加权求平均”：

\begin{defbox}{期望的定义}
若 $X$ 取值为 $x_1,\dots,x_n$，概率为 $p(x_i)$，则对任意函数 $f$：
\[
\E\bigl[f(X)\bigr]=\sum_{i=1}^{n}p(x_i)\,f(x_i).
\]
\textbf{人话}：把每个可能值算出来的 $f(x_i)$，按它出现的概率加权，加起来。
\end{defbox}

\begin{calc}{例子：掷骰子的平均点数}
\[
\E[X]=\sum_{x=1}^{6}\frac16\cdot x=\frac16(1+2+3+4+5+6)=\frac{21}{6}=3.5 .
\]
注意 $3.5$ 根本不是骰子的任何一个面——\textbf{平均值可以不等于任何可能取值}，它只是“长期来看的平均”。这一点很重要，因为熵也是这种“平均”，也常常不是某个具体取值的信息量。
\end{calc}

\begin{qa}{为什么到处都是“加权平均”？}
因为信息论关心的是\textbf{长期的、平均的表现}：一条消息传一次可能很意外也可能很平淡，但传一千万次，平均每次花多少 bit 才是工程上关心的数字（决定要买多贵的线路、多大的硬盘）。所以本书几乎所有核心量都是期望。
\end{qa}

\subsection{对数 $\lb$：这本课唯一的“新数学”}

全书约定：$\lb$ 表示\textbf{以 2 为底的对数}，即 $\lb n$ 表示“2 的几次方等于 $n$”。

\begin{center}
\renewcommand{\arraystretch}{1.25}
\begin{tabular}{c|ccccccccc}
\toprule
$n$ & 1 & 2 & 4 & 8 & 16 & 3 & 5 & 6 & 10\\
\midrule
$\lb n$ & 0 & 1 & 2 & 3 & 4 & 1.585 & 2.322 & 2.585 & 3.322\\
\bottomrule
\end{tabular}
\end{center}

这些数字怎么来的：$2^1=2$ 所以 $\lb2=1$；$2^2=4$ 所以 $\lb4=2$；$2^3=8$ 所以 $\lb8=3$。
不是 2 的整数次方的就不整齐（比如 $\lb 3\approx1.585$ 表示 $2^{1.585}\approx3$），考试时直接写近似值即可。

\begin{defbox}{三个必须记住的对数性质}
\[
\lb(xy)=\lb x+\lb y,\qquad
\lb\frac{x}{y}=\lb x-\lb y,\qquad
\lb(x^{k})=k\lb x .
\]
含意：\textbf{对数把乘法变成加法}。这正是为什么信息量要用对数来定义——两条独立消息的“总信息量”应该等于两条各自信息量之和（加法），而对数正好把“概率相乘”变成“信息量相加”。
\end{defbox}

\begin{qa}{为什么一定要用 log，用别的函数不行吗？}
这正是教材 2.3 节“熵函数的唯一性”在证明的事：只要你承认“信息量应满足几条常识要求”（概率越小信息越大、两条独立消息的信息量应相加），那么\textbf{能用的函数只有对数这一类}，别的函数都通不过检验。所以 log 不是随便选的，是被“逼”出来的——第2讲会给出这个推导。
\end{qa}

\subsection{本讲符号总表（贴在桌上用）}

\begin{center}
\renewcommand{\arraystretch}{1.45}
\begin{tabular}{l|l|l}
\toprule
符号 & 读作 & 意思\\
\midrule
$X$ & “随机变量 X” & 一个还没确定的随机结果\\
$x$ & “某个取值 x” & $X$ 的某一个具体取值\\
$p(x)$ & “x 的概率” & $X=x$ 的可能性，$0\le p(x)\le1$\\
$p(x,y)$ & “x 和 y 的联合概率” & $X=x$ 且 $Y=y$ \textbf{同时}发生的概率\\
$p(y)=\sum_x p(x,y)$ & “边缘分布” & 只看 $Y$、不管 $X$ 时的分布\\
$p(x|y)$ & “给定 y 时 x 的条件概率” & 已知 $Y=y$ 后，$X=x$ 的概率 $=\frac{p(x,y)}{p(y)}$\\
$\E[f(X)]$ & “f(X) 的期望” & 按概率加权的平均值\\
$\lb n$ & “log 以 2 为底 n” & $2^{\lb n}=n$\\
$H(X)$ & “X 的熵” & 平均信息量 / 平均不确定性（第3讲）\\
$H(X|Y)$ & “给定 Y 后 X 的条件熵” & 知道 $Y$ 之后，$X$ 还剩下多少不确定性（第4讲）\\
$I(X;Y)$ & “X 与 Y 的互信息” & 知道 $Y$ 帮你在 $X$ 上省下多少不确定性（第5讲）\\
\bottomrule
\end{tabular}
\end{center}

\begin{exam}{第1讲：考试怎么用}
这一讲本身不会单独出题，但它是后面所有计算的工具。考场上最常用的两招：
\begin{enumerate}[leftmargin=1.6em,itemsep=1pt,topsep=2pt]
  \item \textbf{拿到题先画分布表}。给联合分布的题，第一件事就是把 $2\times2$、$3\times3$ 的表填出来，并把行和列的和加在边缘上——后面不管要求什么熵，都从这张表出发。
  \item \textbf{随时用“加起来等于 1”查错}。如果某一步算出的条件概率加起来不等于 1，立刻回头找错，不要往下算。
\end{enumerate}
\end{exam}

\begin{qa}{自测：你真的会了吗（有答案）}
\textbf{问}：设 $p(X=0,Y=0)=0.1$，$p(X=0,Y=1)=0.4$，$p(X=1,Y=0)=0.3$，$p(X=1,Y=1)=0.2$。求 $p(X=0)$、$p(Y=1)$、$p(X=0|Y=1)$，并判断 $X,Y$ 是否独立。

\textbf{答}：$p(X=0)=0.1+0.4=0.5$；$p(Y=1)=0.4+0.2=0.6$；
$p(X=0|Y=1)=\dfrac{p(0,1)}{p(Y=1)}=\dfrac{0.4}{0.6}=\dfrac23$。
检验独立性：$p(X=0)p(Y=1)=0.5\times0.6=0.3$，而 $p(0,1)=0.4\ne0.3$，故\textbf{不独立}。
\end{qa}

\clearpage
\section{自信息 $I(x)$：一条消息值多少 bit}

\begin{idea}{先做一道猜数字的题，你就懂 $I(x)$ 了}
我心里想了一个 $1$ 到 $8$ 之间的整数，你来猜。你每次只能问一个是非题（“是大于 4 吗？”之类），我如实回答。

最聪明的问法是“对半砍”：先问“大于 4 吗？”——无论我怎么答，剩下的可能性都从 8 个减半成 4 个；再砍一刀变 2 个；再砍一刀变 1 个。所以\textbf{最多 3 次一定能猜中}。

现在换个说法：我一开始心想的那个数，在你眼里有 $8$ 种可能，而我一告诉你答案，可能性从 $8$ 变成 $1$。你猜了 $3$ 次，用了 $3$ 个“是/不是”的回答。于是我们很自然地这样规定：
\[
\text{把可能性砍掉一半，就记 }\ 1\ \text{个信息量（1 个二进制单位）},\qquad \text{砍 3 刀} \Rightarrow 3\ \text{个单位}.
\]
因为 $8=2^3$，所以“$8$ 种可能里告诉你哪一个”正好值 $3$ 个单位，而 $3=\lb 8$。
把这件事翻译成公式：\textbf{如果这件事的先验概率是 $p$，那么它发生所提供的信息量是 $-\lb p$}。
\end{idea}

\subsection{自信息的定义}

\begin{defbox}{定义（自信息 / 信息量）}
设某件事 $x$ 发生的概率为 $p(x)$，则它\textbf{发生（被你知道）时所提供的信息量}，称为 $x$ 的\textbf{自信息}，记作
\[
\boxed{\ I(x)=-\lb p(x)=\lb\frac{1}{p(x)}\ }.
\]
单位是 \textbf{bit}（比特）。
\end{defbox}

\begin{idea}{这个式子的三处“人话”翻译}
\begin{enumerate}[leftmargin=1.6em,itemsep=1pt,topsep=2pt]
  \item \textbf{为什么是 $\frac{1}{p}$ 取对数}：$p(x)$ 越小（越意外），$\frac{1}{p(x)}$ 越大，信息量越大——越意外的消息越值钱。概率为 $1$ 的事（“太阳明天会升起来”）谈不上任何信息。
  \item \textbf{为什么外面有个负号}：因为 $p\le1$，所以 $\lb p\le0$，前面加个负号才能让信息量变成\textbf{非负}的数。所以写 $-\lb p$ 和写 $\lb\frac1p$ 完全是一回事，后者更不容易算错。
  \item \textbf{为什么是 log}：这样才具有“\textbf{可加性}”——两条互不相关消息同时告诉你的信息量，等于各自信息量之和（见下面性质 3）。
\end{enumerate}
\end{idea}

\subsection{手把手算几个最基本的自信息}

\begin{calc}{例 1：抛公平硬币}
$p(\text{正面})=\frac12$，于是
\[
I(\text{正面})=-\lb\frac12=\lb 2=1\bit .
\]
\textbf{这就是“1 bit”这个词的由来}：抛一次公平硬币、告诉你结果，你得到的正好是 1 bit 信息。比如硬盘上一个 bit 存的就是“是/否”。
\end{calc}

\begin{calc}{例 2：掷一颗骰子}
$p(\text{出现 }6)=\frac16$，于是
\[
I=-\lb\frac16=\lb 6\approx2.585\bit .
\]
为什么不是整数？因为 $6$ 不是 2 的整数次方，没法用整数个“是/不是”问题问完。$2.585$ 的意思是“平均需要 2.585 个二进制位才能确定结果”。
同理，从 $8$ 匹马中告诉你哪匹赢（等概率）：$-\lb\frac18=\lb8=3\bit$，正好 3 个问题；从 $64$ 个方块中猜中一个：$-\lb\frac1{64}=6\bit$。
\end{calc}

\begin{calc}{例 3：不公平的硬币（这里藏着本课的核心思想）}
设 $p(\text{正})=0.9$，$p(\text{反})=0.1$，则
\[
I(\text{正})=-\lb 0.9\approx0.152\bit,\qquad
I(\text{反})=-\lb 0.1\approx3.322\bit .
\]
\textbf{体会}：如果硬币严重偏向正面，“正面”这个结果几乎不意外，只值 $0.152$ bit；而一旦出现“反面”，那是很惊人的消息，值 $3.322$ bit。
可见\textbf{“意外的消息更值钱”}。那么平均每次抛硬币能收到多少信息？你不能只取一个数——这就是下一讲的“熵”：把 $0.152$ 和 $3.322$ 按 $0.9$ 和 $0.1$ 加权平均。
\end{calc}

\begin{calc}{例 4：天气预报（考试常用数据）}
设明天天气：晴 $p=0.5$、阴 $p=0.25$、雨 $p=0.25$。各事件的自信息为
\begin{align*}
I(\text{晴})&=-\lb0.5=1\bit,\\
I(\text{阴})&=-\lb0.25=-\lb(2^{-2})=2\bit,\\
I(\text{雨})&=-\lb0.25=2\bit .
\end{align*}
（这里用了一个小技巧：$0.25=\frac14=2^{-2}$，所以 $-\lb0.25=2$，不用查表。）
\end{calc}

\subsection{自信息的三个性质（对应“为什么用 log”的常识要求）}

\begin{defbox}{性质 1：单调递减}
概率越大，信息量越小：
\[
p(x_1)<p(x_2)\ \Longrightarrow\ I(x_1)>I(x_2).
\]
人话：越罕见的事，说出来越有信息。
\end{defbox}

\begin{defbox}{性质 2：非负性}
\[
p(x)\le 1\ \Longrightarrow\ I(x)=-\lb p(x)\ge 0,
\]
且 $p(x)=1$ 时 $I(x)=0$。
人话：“早知道了一定会发生的事”不提供任何信息。
\end{defbox}

\begin{defbox}{性质 3：可加性（这条决定了只能用 log）}
若事件 $x$ 与 $y$ 相互独立（$p(x,y)=p(x)p(y)$），则
\[
I(xy)=-\lb p(x,y)=-\lb\bigl[p(x)p(y)\bigr]=-\lb p(x)-\lb p(y)=I(x)+I(y).
\]
验证一下：抛两次公平硬币，两个结果的信息量应为 $1+1=2$ bit；
直接算：$p(\text{正正})=\frac12\cdot\frac12=\frac14$，$I=-\lb\frac14=2\bit$。$\ $对上！
人话：\textbf{两条互不相干的消息，总信息量等于各自之和}——这几乎是任何“信息量”定义都必须满足的要求，而只有对数能做到（若定义 $I=1/p$，则 $\frac{1}{1/4}=4\ne2+2$，不成立）。
\end{defbox}

\subsection{更一般的三种“自信息”}

除了“一个事件的自信息”，书上还定义了另外两种，写法很像，要分清：

\begin{defbox}{联合自信息、条件自信息}
\[
I(x,y)=-\lb p(x,y),\qquad
I(x|y)=-\lb p(x|y),\qquad
I(x;y)=\lb\frac{p(x,y)}{p(x)p(y)} .
\]
\begin{itemize}[leftmargin=1.6em,itemsep=1pt,topsep=1pt]
  \item $I(x,y)$：\textbf{“$x$ 和 $y$ 同时发生”}这个联合事件的信息量；
  \item $I(x|y)$：\textbf{在已知 $y$ 的条件下}，再告诉你“$x$ 发生了”还值多少信息（因为已知 $y$ 后 $p(x|y)$ 可能变大，所以这个值通常比 $I(x)$ 小）；
  \item $I(x;y)$：互信息，衡量“$y$ 帮你在 $x$ 上省了多少”——它\textbf{可正可负}（第5讲讲）。
\end{itemize}
\end{defbox}

\begin{calc}{手把手：用上一讲那张联合分布表算三种自信息}
表：$p(0,0)=\frac13$，$p(0,1)=\frac13$，$p(1,0)=\frac13$，$p(1,1)=0$；边缘 $p(X=0)=\frac23$，$p(X=1)=\frac13$，$p(Y=0)=\frac23$，$p(Y=1)=\frac13$。

（1）$I(X=1)=-\lb\frac13=\lb3\approx1.585$ bit（$\frac13$ 比 $\frac12$ 小，所以比 1 bit 略大）。
（2）联合自信息 $I(X=0,Y=0)=-\lb\frac13=\lb3\approx1.585$ bit。
（3）条件自信息：已知 $Y=0$ 后 $X=0$ 的条件概率是 $\frac12$（上一讲算过），故
\[
I(X=0|Y=0)=-\lb\frac12=1\bit .
\]
对比一下：$I(X=0)=\lb\frac32\approx0.585$ bit，而 $I(X=0|Y=0)=1$ bit。
\textbf{什么？知道 $Y=0$ 之后反而更意外了？}是的，这种情况完全可能：因为 $Y=0$ 这个条件把 $p(X=0)$ 从 $\frac23$ 压到了 $\frac12$。所以\textbf{“条件不增信息量”并不是单条事件成立的}，它只对“平均以后”的熵成立——这是初学者最容易错的地方，第4讲会专门讲。
\end{calc}

\begin{exam}{第2讲：考试怎么用}
\begin{enumerate}[leftmargin=1.6em,itemsep=1pt,topsep=2pt]
  \item \textbf{直接给概率求信息量}：$I=-\lb p$，套公式，注意单位写 bit。答题时写“$I(x)=-\lb 0.25=2$ bit”这样的完整一行即可得分。
  \item \textbf{反着问“需要多少位”}：若信息量是 $I$ bit，则大约需要 $I$ 个二进制位表示（有小数时向上取整）。考试里常问“$n$ 元等概消息需要多少 bit”$\Rightarrow \lb n$。
  \item \textbf{判断题“概率越小信息量越大”}：正确；\textbf{“概率为 0 的事件信息量为无穷大”}：这是数学上的极限（现实中不会发生），一般不出这种题，若问就答“$I$ 随 $p\to0$ 而趋于无穷”。
  \item \textbf{千万别把 $I(x)$ 和 $H(X)$ 混为一谈}：$I(x)$ 是“某一条消息”的信息量（随具体取值变化，有大小不一定相等）；$H(X)$ 是“平均每条”的信息量（一个数）。下一讲就是讲 $H$。
\end{enumerate}
\end{exam}

\begin{qa}{自测：你真的会了吗（有答案）}
\textbf{问}：某信源发出四种符号，概率分别为 $1/2,\ 1/4,\ 1/8,\ 1/8$。求每个符号的自信息，以及（不用算平均）哪个符号最有信息量。

\textbf{答}：$I_1=-\lb\frac12=1$ bit；$I_2=-\lb\frac14=2$ bit；$I_3=I_4=-\lb\frac18=3$ bit。
最有信息量的是概率最小、也就是最罕见的两个符号（各 3 bit）。
顺便说：如果你要设计编码，直觉上也应该把最短的码字给最常出现的符号——这正是第8讲 Huffman 编码的全部思想。
\end{qa}

\clearpage
\section{熵 $H(X)$：平均每条消息值多少 bit}

\begin{idea}{一句话：熵就是“自信息的加权平均”}
上一讲算出了每条消息各自值多少 bit，但每条都不一样（$0.152$ 和 $3.322$）。工程上我们要一个\textbf{平均指标}：长期来看，平均每发一条消息，你收到多少 bit？

答案就是把各自的信息量按它们出现的概率加权平均：
\[
\Hx{X}=p(x_1)I(x_1)+p(x_2)I(x_2)+\cdots=\sum_x p(x)\cdot\bigl(-\lb p(x)\bigr).
\]
这个“平均自信息”就叫\textbf{熵}，书上记作 $H(X)$（$H$ 源自物理学里的“熵”，表示混乱/不确定程度）。
\end{idea}

\subsection{熵的定义}

\begin{defbox}{定义 2.1.3（信源熵 / 信息熵 / 香农熵）}
设离散随机变量 $X$ 的取值集合为 $\{x_1,x_2,\dots,x_n\}$，概率为 $p(x_i)$，则 $X$ 的\textbf{熵}定义为
\[
\boxed{\ \Hx{X}=\E\bigl[I(x)\bigr]=\sum_{i=1}^{n}p(x_i)\bigl(-\lb p(x_i)\bigr)
=-\sum_{i=1}^{n}p(x_i)\lb p(x_i)\ }
\]
单位 bit（对数底为 2 时）；若底取 $e$ 则单位是 nat，底取 10 则单位是 Hart，但\textbf{本课程统一用 $\lb$，单位一律 bit}。
\end{defbox}

\begin{qa}{这条公式到底在算什么？三句话}
\begin{enumerate}[leftmargin=1.6em,itemsep=1pt,topsep=2pt]
  \item $-\lb p(x_i)$ 是“第 $i$ 种情况\textbf{一旦发生}值多少 bit”；
  \item 再乘上 $p(x_i)$ 表示“这种情况\textbf{出现的可能性有多大}，就按多大份额算”；
  \item 全部加起来 = “平均每种情况值多少 bit”，也就是\textbf{平均不确定性}。
\end{enumerate}
所以同一个 $H(X)$ 有两个等价的说法：\textbf{“平均信息量”}（告诉你结果后你平均收到多少 bit）与\textbf{“平均不确定性”}（告诉你之前你对结果有多糊涂）。这两者是同一件事的两面：不确定性越大，一旦告诉你，你得到的信息就越多。
\end{qa}

\subsection{手把手算熵（三个必须会算的例子）}

\begin{calc}{例 1：掷一颗骰子（等概率情形）}
六种结果各 $\frac16$，于是
\begin{align*}
\Hx{X}&=-\sum_{x=1}^{6}\frac16\lb\frac16
=-6\cdot\frac16\cdot\lb\frac16\\
&=-\lb\frac16=\lb 6\approx2.585\bit .
\end{align*}
（这里用了一个省力技巧：$6$ 项完全一样，直接乘 $6$。注意 $-\lb\frac16$ 就是“六种等可能”的不确定性，与上一讲单条消息的自信息相同——因为等概率时每条消息信息量一样，平均值当然不变。）

\textbf{结论记牢}：等概 $n$ 元信源的熵是 $\lb n$。特别地，抛公平硬币 $\Hx{X}=\lb2=1$ bit。
\end{calc}

\begin{calc}{例 2：不公平硬币（二元信源，全书最重要的一族）}
设 $p(\text{正})=p$，$p(\text{反})=1-p$。定义
\[
H(p)=-p\lb p-(1-p)\lb(1-p),
\]
这个函数叫\textbf{二元熵函数}。代几个数进去：
\begin{center}
\renewcommand{\arraystretch}{1.25}
\begin{tabular}{c|ccccccc}
\toprule
$p$ & 0（或1） & 0.01 & 0.1 & 0.25 & 0.5 & 0.75 & 0.9\\
\midrule
$H(p)$ & 0 & 0.081 & 0.469 & 0.811 & 1 & 0.811 & 0.469\\
\bottomrule
\end{tabular}
\end{center}
\textbf{怎么算的（以 $p=0.1$ 为例）}：
\begin{align*}
H(0.1)&=-0.1\times\lb0.1-0.9\times\lb0.9\\
&=-0.1\times(-3.322)-0.9\times(-0.152)\\
&=0.332+0.137=0.469\bit .
\end{align*}
\textbf{读出三条信息}：
\begin{itemize}[leftmargin=1.6em,itemsep=1pt,topsep=1pt]
  \item $p=\frac12$ 时最大，等于 $1$ bit —— 这时最猜不准；
  \item $p=0$ 或 $p=1$ 时为零 —— 结果完全确定（永远正面或永远反面），没有任何不确定性；
  \item 函数关于 $p=\frac12$ 对称：$H(0.1)=H(0.9)$（“偏 0.1”和“偏 0.9”一样糊涂）。
\end{itemize}
\end{calc}

\begin{calc}{例 3：天气预报（考试常用）}
晴 $0.5$、阴 $0.25$、雨 $0.25$：
\begin{align*}
\Hx{X}&=-0.5\lb0.5-0.25\lb0.25-0.25\lb0.25\\
&=0.5\times1+0.25\times2+0.25\times2\\
&=0.5+0.5+0.5=1.5\bit .
\end{align*}
\textbf{算这类题的标准写法}：先列表（概率、$-\lb p$、乘积），再求和：
\begin{center}
\renewcommand{\arraystretch}{1.25}
\begin{tabular}{c|c|c|c}
\toprule
天气 & $p_i$ & $-\lb p_i$ & $p_i(-\lb p_i)$\\
\midrule
晴 & 0.5 & 1 & 0.5\\
阴 & 0.25 & 2 & 0.5\\
雨 & 0.25 & 2 & 0.5\\
\midrule
合计 & 1 & — & $\Hx{X}=1.5$\\
\bottomrule
\end{tabular}
\end{center}
这一列“$p_i(-\lb p_i)$”加起来就是熵。\textbf{考试时永远这么列表}，既清楚又不容易算错。
\end{calc}

\begin{defbox}{约定：$0\lb 0=0$}
如果某个 $p(x_i)=0$（这种情况根本不可能发生），那么它对应的项 $0\cdot\lb 0$ 规定为 $0$，即“不可能发生的事不贡献任何不确定性”。
\textbf{人话}：算熵时可以把概率为 0 的项直接划掉不写。
（严格说这是极限 $\lim_{p\to0^+}p\lb p=0$，你在做作业时直接写“含 0 的项略去”就行。）
\end{defbox}

\subsection{熵的三重含义（同一个公式，三种说法）}

\begin{idea}{同一件事的三种说法，考试都可能问}
\begin{enumerate}[leftmargin=1.6em,itemsep=2pt,topsep=2pt]
  \item \textbf{不确定性的度量}：收到结果之前，你对它有多糊涂。$H$ 越大越糊涂。
  \item \textbf{平均信息量}：收到结果之后，你平均收获多少 bit。$H$ 越大，这条消息平均越“值钱”。
  \item \textbf{无损压缩的下限}：平均每个符号\textbf{至少要}用 $H(X)$ 个二进制位来表示，做不到更短。这就是第8讲信源编码里 $\bar L\ge H(X)$ 的来源。
\end{enumerate}
第三条最实用：它给了“压缩的极限”。比如一段中文每个字熵是 8 bit，那你无论怎么压缩，都不可能平均用少于 8 bit 表示一个字。
\end{idea}

\subsection{熵的性质（考试判断/填空的高频点）}

\begin{defbox}{性质 1：非负性}
\[
\Hx{X}\ge 0,
\]
且 $\Hx{X}=0$ 当且仅当 $X$ 的取值完全确定（某个 $p(x_i)=1$，其余全为 0）。
理由：每一项 $p_i(-\lb p_i)$ 都 $\ge0$。
\end{defbox}

\begin{defbox}{性质 2：极值性（离散熵最大值，必背）}
设 $X$ 有 $n$ 个可能取值，则
\[
0\le \Hx{X}\le \lb n,
\]
且\textbf{右端取等号当且仅当所有取值等可能}（$p(x_i)=\frac1n$）。
\textbf{人话}：在所有“$n$ 种结果”的分布里，\textbf{等概率是最不确定的}。这一点直觉上也对：如果你知道某个结果特别容易出现，那你就没那么糊涂了。
例：骰子 $\Hx{X}=\lb6\approx2.585$ 是六种收手里的最大值；换成概率 $(0.5,0.1,0.1,0.1,0.1,0.1)$ 后熵只有约 $2.16$ bit。
\end{defbox}

\begin{defbox}{性质 3：对称性}
熵只看\textbf{概率的集合}，与“哪个概率对应哪个符号”无关：
\[
\Hx{p_1,p_2,\dots,p_n}=\Hx{p_2,p_1,\dots,p_n}\ \ \text{等等任意重排}.
\]
人话：把“晴 0.5、雨 0.5”改成“晴 0.5、雪 0.5”，熵不变——熵不关心符号叫什么，只关心概率长什么样。
\end{defbox}

\begin{defbox}{性质 4：可加性（分组原则，做难题要用）}
若把 $n$ 个符号分个组，先按组算一次熵、再在组内算一次熵，两次相加等于原来的熵：
\[
\Hx{p_1,\dots,p_n}=H(\text{组概率})+\sum_{\text{组}}p(\text{组})\cdot H(\text{组内条件概率}).
\]
\textbf{人话}：不确定性可以“先分大类、再分小类”分批算，加起来一样。
\textbf{用处}：遇到“概率里有一大堆 0.01 的小概率符号”这类题，可以先归并成两组，能省很多计算。
\end{defbox}

\begin{defbox}{性质 5：上凸性}
$H(p)$ 是概率向量 $\mathbf{p}=(p_1,\dots,p_n)$ 的\textbf{上凸函数}；$-H$ 是下凸函数。
\textbf{人话}：把两个分布“掺”在一起（取凸组合），熵不会比原来两个的平均值更小——\textbf{混合只会让事情更不确定}。
（这也是“等概时熵最大”的另一种说法。）
\end{defbox}

\subsection{为什么偏偏是“熵”这个函数？——唯一性定理一句话版}

教材 2.3 节证明了：如果要求一个“不确定性度量”函数 $H(p_1,\dots,p_n)$ 满足三条常识：
\begin{enumerate}[leftmargin=1.6em,itemsep=1pt,topsep=2pt]
  \item \textbf{连续性}：概率稍微变一点，熵也只能稍微变一点（不能跳）；
  \item \textbf{单调性}：等概情况下，可选结果越多熵越大（$\lb n$ 随 $n$ 增大）；
  \item \textbf{分组可加性}：像上面性质 4 那样可以分批算；
\end{enumerate}
那么\textbf{唯一可能的函数形式就是} $H=-\sum p_i\lb p_i$（相差一个常数倍数）。

\begin{idea}{这一条对你的实际意义}
考试里如果问“为什么要用这个公式”“熵函数为什么唯一”，标准答法就是：
\textbf{“因为它满足连续性、单调性、分组可加性三条公理，而这三条公理唯一确定了 $H=-\sum p\lb p$ 这一形式。”}——把三条公理的名字写出来就是分。
\end{idea}

\subsection{初学者最常见的四个错}

\begin{qa}{四个必查的坑}
\begin{enumerate}[leftmargin=1.6em,itemsep=1pt,topsep=2pt]
  \item \textbf{忘了检查概率和为 1}。比如给 $(0.5,0.3,0.3)$（和为 1.1），要先归一化成 $(0.4545,0.2727,0.2727)$ 这类东西，或题目本身有误——先查这个，能救回很多分。
  \item \textbf{把 $\lb$ 当成自然对数}。课本一律 $\lb$；如果你按 $\ln$ 算了，结果要再除以 $\ln2\approx0.693$。例如 $\ln6=1.792$，而 $\lb6=2.585$（$1.792/0.693=2.585$）。
  \item \textbf{把 $H(X)$ 当成某一个符号的信息量}。$H$ 是平均值；单个符号的信息量是 $I(x)=-\lb p(x)$，两者只有在等概时才相等。
  \item \textbf{漏掉概率为 0 的项还好，但漏掉“小概率项”就是致命错误}。比如 $p=(0.99,0.01)$，那个 $0.01$ 看起来可以忽略，实际上贡献了 $0.01\times6.64=0.066$ bit，占全部熵 $0.081$ 的八成以上。
\end{enumerate}
\end{qa}

\begin{exam}{第3讲：考试怎么用}
\begin{enumerate}[leftmargin=1.6em,itemsep=1pt,topsep=2pt]
  \item \textbf{基本计算题}：给概率分布求 $H$。答题格式：列“$p_i$ / $-\lb p_i$ / 乘积”三列表格，最后求和，写单位 bit。
  \item \textbf{问最大熵}：$n$ 个符号的最大熵是 $\lb n$，在等概时达到；二元时最大是 1 bit（$p=\frac12$）。
  \item \textbf{判断/改错题}：如“等概信源的熵最大”“离散熵不能为负”“概率为 0 的项可以略去”——这三条都对。
  \item \textbf{与编码挂钩}：问“这段消息至少需要多少 bit”——答 $\Hx{X}$ 再乘符号数；问“最多能压缩到什么程度”——答 $H(X)$ 就是下限。
\end{enumerate}
\end{exam}

\begin{qa}{自测：三问（有答案）}
（1）$X$ 等可能取 $1,2,3,4$，$\Hx{X}=$？\\
（2）$X$ 取 $0.9,0.1$ 两个值，$\Hx{X}\approx$？\\
（3）$X$ 的概率分布为 $(0.6,0.2,0.2)$，$\Hx{X}=$？

\textbf{答}：
（1）$\lb4=2$ bit。
（2）$H(0.1)=0.469$ bit（查上面的表）。
（3）$-0.6\lb0.6-2\times0.2\lb0.2$，其中 $-\lb0.6=0.737$、$-\lb0.2=2.322$，故
$0.6\times0.737+0.2\times2.322+0.2\times2.322=0.442+0.464+0.464=1.371$ bit。
\end{qa}

\clearpage
\section{联合熵与条件熵：两个变量一起看}

\begin{idea}{这一讲要回答三个问题}
\begin{enumerate}[leftmargin=1.6em,itemsep=2pt,topsep=2pt]
  \item 如果一次同时告诉你 $X$ 和 $Y$ 两个值，这份“打包消息”值多少 bit？——\textbf{联合熵 $H(X,Y)$}
  \item 如果已经告诉了你 $Y$，再告诉你 $X$ 还值多少 bit？——\textbf{条件熵 $H(X|Y)$}
  \item 这两者与单个的 $H(X)$、$H(Y)$ 之间是什么关系？——\textbf{链式法则}
\end{enumerate}
只要把这三个问题算清楚，第2章的大半分数就到手了（作业 2.1、2.7 全是这一讲的内容）。
\end{idea}

\subsection{联合熵 $H(X,Y)$}

\begin{defbox}{定义}
\begin{align*}
\Hx{X,Y}&=\E\bigl[-\lb p(X,Y)\bigr]
=-\sum_{x}\sum_{y}p(x,y)\lb p(x,y)\\
&=\sum_{x}\sum_{y}p(x,y)\lb\frac{1}{p(x,y)} .
\end{align*}
\textbf{人话}：“把 $X$ 和 $Y$ \textbf{同时}告诉你”这份打包消息平均值多少 bit，也就是两个变量的\textbf{总不确定性}。
\end{defbox}

\begin{calc}{手把手算联合熵（就用这张 2×2 表，务必跟着算一遍）}
联合分布表：
\begin{center}
\renewcommand{\arraystretch}{1.3}
\begin{tabular}{c|cc|c}
\toprule
 & $Y=0$ & $Y=1$ & 行和（$X$ 的边缘）\\
\midrule
$X=0$ & $1/3$ & $1/3$ & $2/3$\\
$X=1$ & $1/3$ & $0$ & $1/3$\\
\midrule
列和（$Y$ 的边缘） & $2/3$ & $1/3$ & $1$\\
\bottomrule
\end{tabular}
\end{center}

\textbf{第一步：列出所有非零格子的 $p$ 与 $-\lb p$。}（$p(1,1)=0$，按约定 $0\lb0=0$，直接略去。）
\begin{center}
\renewcommand{\arraystretch}{1.3}
\begin{tabular}{c|c|c|c}
\toprule
格子 & $p(x,y)$ & $-\lb p(x,y)$ & 乘积\\
\midrule
$(0,0)$ & $1/3$ & $\lb3\approx1.585$ & $0.528$\\
$(0,1)$ & $1/3$ & $1.585$ & $0.528$\\
$(1,0)$ & $1/3$ & $1.585$ & $0.528$\\
\midrule
合计 & & & $\Hx{X,Y}\approx1.585$\\
\bottomrule
\end{tabular}
\end{center}

\textbf{第二步：加起来。}
\[
\Hx{X,Y}=3\times\frac13\times\lb3=\lb3\approx1.585\bit .
\]
（技巧：三个格子概率都是 $\frac13$，所以直接乘 3。更通用地说，若三个非零格概率均等，“打包不确定性”就是 $\lb3$。）
\end{calc}

\begin{qa}{为什么这里 $H(X,Y)=\lb3$ 这么“整”？}
因为联合分布里只有 3 个等可能的组合，等于“从 3 个组合里等可能挑一个”，按第3讲结论，熵就是 $\lb3$。
\textbf{反过来说}：只要你在表上看到“$k$ 个非零且相等的格子”，联合熵立刻就是 $\lb k$，不用一个个算。这个小技巧在考场上能省一两分钟。
\end{qa}

\subsection{条件熵 $H(X|Y)$：知道 $Y$ 之后还剩多少不确定性}

\begin{defbox}{定义（分两步看最清楚）}
\textbf{第一步}：给定某个具体的 $Y=y$，此时 $X$ 的条件分布是 $p(x|y)$，它的熵叫\textbf{条件熵（给定 $y$）}：
\[
\Hx{X|Y=y}=-\sum_{x}p(x|y)\lb p(x|y).
\]
\textbf{第二步}：$y$ 本身是随机的，各种 $y$ 出现的概率不同，所以加权平均：
\[
\boxed{\ \Hx{X|Y}=\sum_{y}p(y)\,\Hx{X|Y=y}\ }.
\]
\end{defbox}

\begin{calc}{手把手算条件熵（继续用那张表）}
\textbf{（1）先算 $Y=0$ 的情形。}已知 $Y=0$（概率 $\frac23$），条件分布上一讲算过：
\[
p(X=0|Y=0)=\frac{p(0,0)}{p(Y=0)}=\frac{1/3}{2/3}=\frac12,\qquad
p(X=1|Y=0)=\frac{1/3}{2/3}=\frac12 .
\]
这是“一半一半”，所以
\[
\Hx{X|Y=0}=-\frac12\lb\frac12-\frac12\lb\frac12=\frac12+\frac12=1\bit .
\]

\textbf{（2）再算 $Y=1$ 的情形。}已知 $Y=1$（概率 $\frac13$）：
\[
p(X=0|Y=1)=\frac{p(0,1)}{p(Y=1)}=\frac{1/3}{1/3}=1,\qquad
p(X=1|Y=1)=\frac{0}{1/3}=0 .
\]
完全确定，所以
\[
\Hx{X|Y=1}=-1\lb1=0\bit .
\]

\textbf{（3）加权平均}（这一步千万别漏，很多同学直接写 $1$ 就错了）：
\[
\Hx{X|Y}=p(Y=0)\cdot1+p(Y=1)\cdot0
=\frac23\times1+\frac13\times0=\frac23\approx0.667\bit .
\]

\textbf{（4）对照检查。}原来 $\Hx{X}=-\frac23\lb\frac23-\frac13\lb\frac13=0.918$ bit；
知道 $Y$ 之后剩下 $0.667$ bit。所以知道 $Y$ 帮你省下了
\[
0.918-0.667=0.251\bit
\]
的不确定性——这个差值就是第5讲的“互信息”。\textbf{记住这个“省下来的差额”的说法，它是整本书的灵魂。}
\end{calc}

\begin{qa}{为什么“条件熵”要加权平均，而不能直接把各个 $\Hx{X|Y=y}$ 加起来？}
因为不同的 $y$ 出现的频率不一样。如果 $Y=1$ 这种情况极其罕见（概率 0.001），那它对“长期平均”的影响就微乎其微。所以必须按 $p(y)$ 加权。\textbf{一句话：条件熵 = 各种情况下的剩余不确定性，按该情况出现的概率加权平均。}
\end{qa}

\subsection{链式法则：拆开算的两个方向}

\begin{defbox}{链式法则（本章第一公式，必背）}
\[
\boxed{\ \Hx{X,Y}=\Hx{X}+\Hx{Y|X}=\Hx{Y}+\Hx{X|Y}\ }
\]
\textbf{人话翻译（第一式）}：要把“$X$ 和 $Y$ 同时告诉你”的不确定性拆开算，可以这样分两步——
\textbf{先告诉 $X$}（这份消息值 $\Hx{X}$），\textbf{再告诉 $Y$}（这时你已经知道 $X$ 了，所以 $Y$ 只值 $\Hx{Y|X}$）。两份加起来就是打包的总不确定性。
\end{defbox}

\begin{calc}{用刚才的数验证链式法则（考场上的验算手段）}
右边：$\Hx{X}+\Hx{Y|X}=0.918+0.667=1.585$；左边：$\Hx{X,Y}=1.585$。对上。

顺带一个变形（后面做作业常用）：
\[
\Hx{X|Y}=\Hx{X,Y}-\Hx{Y},\qquad \Hx{Y|X}=\Hx{X,Y}-\Hx{X}.
\]
代进去：$\Hx{X|Y}=1.585-0.918=0.667$ ✓，$\Hx{Y|X}=1.585-0.918=0.667$ ✓。
\textbf{注意}：这次两个值恰好相等，是因为这张表的结构特殊（行和与列和一模一样）。\textbf{一般情形下两者不相等}，这是作业 2.1 最重要的考点。
\end{calc}

\subsection{两条最重要的不等式}

\begin{defbox}{不等式 1：条件不增熵（知道得越多，越不糊涂）}
\[
\Hx{X|Y}\le\Hx{X},\qquad \Hx{Y|X}\le\Hx{Y}.
\]
\textbf{人话}：多知道一个信息，只会让你对 $X$ 更确定，绝不会更糊涂。
\textbf{等号成立条件}：$X$ 与 $Y$ 独立（知道了 $Y$ 等于白知道）。
\end{defbox}

\begin{qa}{疑点：上一讲我算出来“知道 $Y=0$ 之后 $X$ 的信息量反而变大了”，是不是矛盾？}
不矛盾，而且这是极好的问题。区别在于“\textbf{一条事件}”和“\textbf{平均}”：
\begin{itemize}[leftmargin=1.6em,itemsep=1pt,topsep=1pt]
  \item 单条事件的条件自信息 $I(x|y)=-\lb p(x|y)$ 可能比 $I(x)$ 大（当 $p(x|y)<p(x)$ 时）；
  \item 但\textbf{平均以后}（即条件熵 $\Hx{X|Y}$）一定不大于 $\Hx{X}$。
\end{itemize}
所以在卷子上写“因为 $p(x|y)$ 可能小于 $p(x)$，所以条件熵更大”是\textbf{错误}的；正确的说法永远是“条件熵不超过原熵”。
\end{qa}

\begin{defbox}{不等式 2：次可加性（作业要求证明的那一条）}
\[
\boxed{\ \Hx{X,Y}\le\Hx{X}+\Hx{Y}\ }
\]
\textbf{人话}：两个变量“打包”的不确定性，不会超过“分别算再相加”。为什么？因为分开算时，你把“$Y$ 和 $X$ 可能重复的部分”算了两遍（$\Hx{Y|X}\le\Hx{Y}$ 正是这个意思）。
\textbf{等号成立条件}：$X$ 与 $Y$ 相互独立。
\textbf{一句话证明}：
\[
\Hx{X,Y}=\Hx{X}+\Hx{Y|X}\ \overset{\Hx{Y|X}\le\Hx{Y}}{\le}\ \Hx{X}+\Hx{Y}.
\]
（作业里要求完整证明，我会在《作业一详解》里给出三种证法。）
\end{defbox}

\begin{calc}{用具体数字体会“打包比分开省”}
本题：$\Hx{X,Y}=1.585$，而 $\Hx{X}+\Hx{Y}=0.918+0.918=1.836$。
确实 $1.585<1.836$，省下了 $0.251$ bit —— 正好是刚才算出的“知道 $Y$ 省下的不确定性”！
所以有一个非常重要的恒等式（第5讲再正式出现）：
\[
\Hx{X}+\Hx{Y}-\Hx{X,Y}=\I{X;Y}
\quad(\text{互信息}).
\]
\end{calc}

\subsection{把作业 2.7 的“新变量”先演练一遍}

作业 2.7 会给你上面那张表，再定义 $Z=X\oplus Y$（$\oplus$ 是\textbf{模 2 加}：两个相同为 0，不同为 1），让你求 $\Hx{X|Z}$ 之类。这里先演示怎么求 $Z$ 的分布——老师说这一步错得最多。

\begin{calc}{手把手：求 $Z=X\oplus Y$ 的分布}
逐格算 $z=x\oplus y$：
\begin{center}
\renewcommand{\arraystretch}{1.3}
\begin{tabular}{c|c|c|c}
\toprule
$(x,y)$ & 概率 & $z=x\oplus y$ & 说明\\
\midrule
$(0,0)$ & $1/3$ & $0\oplus0=0$ & 相同 $\Rightarrow0$\\
$(0,1)$ & $1/3$ & $0\oplus1=1$ & 不同 $\Rightarrow1$\\
$(1,0)$ & $1/3$ & $1\oplus0=1$ & 不同 $\Rightarrow1$\\
$(1,1)$ & $0$ & $1\oplus1=0$ & 这一格概率为 0，不影响\\
\bottomrule
\end{tabular}
\end{center}
把 $z$ 相同的概率并起来：
\[
p(Z=0)=\frac13+0=\frac13,\qquad p(Z=1)=\frac13+\frac13=\frac23 .
\]
（检查：$\frac13+\frac23=1$ ✓。）

于是 $\Hx{Z}=-\frac13\lb\frac13-\frac23\lb\frac23=0.528+0.390=0.918$ bit。

\textbf{接下去求 $\Hx{X|Z}$ 的套路}：分别看 $Z=0$、$Z=1$ 时 $X$ 的条件分布：
$Z=0$ 只可能来自 $(0,0)$（概率 $\frac13$），此时 $X$ 必为 0，故 $\Hx{X|Z=0}=0$；
$Z=1$ 来自 $(0,1)$ 与 $(1,0)$ 各 $\frac13$，条件分布为 $(\frac12,\frac12)$，故 $\Hx{X|Z=1}=1$ bit。
加权：
\[
\Hx{X|Z}=\frac13\times0+\frac23\times1=\frac23\bit .
\]
\textbf{注意别踩的坑}：有人直接把 $\Hx{X|Z}$ 写成 $1$（忘了加权）或写成 $0$（只看了 $Z=0$ 的情形）。\textbf{一定要“分别算 + 加权”。}
\end{calc}

\begin{exam}{第4讲：考试怎么用}
\begin{enumerate}[leftmargin=1.6em,itemsep=1pt,topsep=2pt]
  \item \textbf{给联合分布表求各种熵}的标准流程（照这个顺序写，稳）：
        “① 补全表格边缘 → ② 算 $\Hx{X},\Hx{Y}$（各自一维）→ ③ 算 $\Hx{X,Y}$（对各非零格）→ ④ 用链式法则得条件熵 → ⑤ 用 $\Hx{X}+\Hx{Y}-\Hx{X,Y}$ 得互信息”。
  \item \textbf{求 $\Hx{X|Z}$ 这类“新变量”题}：诀窍是“先求 $Z$ 的分布 → 再对每个 $Z$ 值求条件分布 → 最后加权”。老师说这一步最容易错，一定要把三小步都写在卷面上。
  \item \textbf{常见陷阱}：
        （i）$\Hx{X|Y}\ne\Hx{Y|X}$；
        （ii）条件熵要加权；
        （iii）概率为 0 的格子按 $0\lb0=0$ 略去，不要写“无穷大”；
        （iv）$Z$ 是 $X,Y$ 的确定性函数时，$\Hx{Z|X,Y}=0$、$\Hx{X,Y,Z}=\Hx{X,Y}$（这条能秒杀很多题）。
\end{enumerate}
\end{exam}

\begin{qa}{自测：（有答案）}
\textbf{问}：设 $p(0,0)=0.25$，$p(0,1)=0.25$，$p(1,0)=0.25$，$p(1,1)=0.25$（就是“独立且各半”）。求 $\Hx{X},\Hx{Y},\Hx{X,Y},\Hx{X|Y}$，并说明为什么 $\Hx{X,Y}=\Hx{X}+\Hx{Y}$。

\textbf{答}：边缘都是 $(0.5,0.5)$，故 $\Hx{X}=\Hx{Y}=1$ bit；联合四个格各 $0.25$，故 $\Hx{X,Y}=\lb4=2$ bit；
$\Hx{X|Y}=\Hx{X,Y}-\Hx{Y}=2-1=1$ bit。
这里 $X,Y$ 独立，故 $\Hx{X|Y}=\Hx{X}$，于是 $\Hx{X,Y}=1+1=2$，等号成立——正是次可加性的等号条件。
\end{qa}

\clearpage
\section{互信息：知道一件事，能帮你猜中另一件事多少}

% ==================================================================
\subsection{先说人话：你每天早上都在用互信息}

先不讲公式，讲个你天天遇到的事。

早上 7 点，你站在宿舍门口，脑子里有两个问号：
\begin{enumerate}[leftmargin=1.8em,itemsep=1pt,topsep=2pt]
  \item \textbf{今天下雨吗？}（记为 $X$）
  \item \textbf{我要不要带伞？}（记为 $Y$）
\end{enumerate}
你对第一个问号是\emph{不确定}的：也许是 $30\%$ 会下雨，也许是 $90\%$。
这种“心里没底”的程度，信息论里有个专门的度量，叫\textbf{不确定性}
（后面你会看到它就是熵 $H(X)$）。

现在楼下大爷喊了一句：“今儿下雨！”（这是一个消息）。
这一句话，让你对“带不带伞”的犹豫几乎消失了——原来你还在纠结，
现在你基本确定要带伞。

学长要强调的是这句话的\textbf{价值}：大爷这句话给你带来的东西，
不是“下雨”这两个字，而是那句更抽象的话——
\begin{center}
\fbox{\parbox{0.86\textwidth}{\centering
\textbf{你关于“带不带伞”的不确定性，减少了多少}\\[2pt]
\footnotesize 减少的这部分不确定性，就是这句话的\textbf{信息量}。
}}
\end{center}
\textbf{互信息（mutual information）衡量的就是这个“减少了多少”。}
写成公式的样子，就是本讲的主角 $\I{X;Y}$，读作“$X$ 和 $Y$ 的互信息”，
意思是“\textbf{知道 $Y$ 之后，我对 $X$ 的不确定性减少了多少}（单位 bit）”。

\begin{idea}{本讲唯一的“一根筋”思路}
整讲你只需要记住一句话，剩下的全是这句话的变体：

\medskip
\textbf{信息量 $=$ 不确定性的减少量。}

\medskip
\begin{itemize}[leftmargin=1.4em,itemsep=1pt,topsep=1pt]
  \item 一件\textbf{本来就没准}的事（$p=0.5$），告诉你了，信息量\textbf{大}；
  \item 一件\textbf{本来就很准}的事（$p=0.99$），再说一遍，信息量\textbf{小}；
  \item 知道 $Y$ 之后 $X$ 的不确定性降低得越多，$\I{X;Y}$ 就越大。
\end{itemize}
\end{idea}

\subsubsection{先把这四个符号认全（不认全后面全是天书）}

\begin{center}
\small
\begin{tabular}{@{}lll@{}}
\toprule
符号 & 怎么读 & 大白话 \\
\midrule
$p(x)$ & “$p$ 括号 $x$” & 事件“$X$ 取值为 $x$”发生的\textbf{概率}，一个 $0\sim1$ 的数 \\
$I(x)$ & “$x$ 的自信息” & 事件 $x$ 真的发生了这件事，本身\textbf{含有多少信息量} \\
 & & $I(x)=-\lb p(x)$，单位 bit \\
$H(X)$ & “$X$ 的熵” & $X$ \textbf{整体的平均不确定性}，$H(X)=\sum_x p(x)I(x)$ \\
$H(X|Y)$ & “在 $Y$ 的条件下 $X$ 的条件熵” & \textbf{知道 $Y$ 之后}，$X$ 还剩多少不确定性 \\
$\I{X;Y}$ & “$X$ 和 $Y$ 的互信息” & 知道 $Y$ 之后，$X$ 的不确定性\textbf{减少了多少} \\
\bottomrule
\end{tabular}
\end{center}

还有一条本讲要反复用的“换算器”（考试不给你，要背）：
\[
\lb 2=1,\qquad \lb 3\approx1.58496\approx1.585,\qquad
\lb 5\approx2.32193\approx2.322,\qquad \lb 10\approx3.32193\approx3.322 .
\]
用法举例：$\lb 0.4=\lb\frac{2}{5}=\lb 2-\lb 5=1-2.32193=-1.32193$，
$\lb 0.3=\lb\frac{3}{10}=1.58496-3.32193=-1.73697$。
看到小数就用“拆成分数、再用上面三个数”的办法，别去按计算器蒙。

另外约定一件事：本讲用小写 $x,y$ 表示\textbf{一个具体的取值}（一个事件），
用大写 $X,Y$ 表示\textbf{整个随机变量}（所有取值一起）。这个大小写之别，
正是本讲前半段（事件互信息）和后半段（平均互信息）的区别。

% ==================================================================
\subsection{事件的互信息：一件事能“泄露”多少关于另一件事的信息}

\subsubsection{定义，以及分子分母到底是什么}

\begin{defbox}{定义（事件之间的互信息）}
设 $X,Y$ 是两个离散随机变量，$p(x,y)$ 读作“$x$ 和 $y$ \textbf{同时}发生的概率”
（联合概率），$p(x)$、$p(y)$ 分别是“$x$ 单独发生的概率”和“$y$ 单独发生的概率”
（边缘概率）。事件 $y$ 对事件 $x$ 的\textbf{互信息}定义为
\[
\I{x;y}=\lb\frac{p(x,y)}{p(x)\,p(y)} .
\]
底取 $2$ 时单位是 \textbf{bit}（本课程一律用底 $2$）。
\end{defbox}

\begin{idea}{三个零件的分工}
把 $\I{x;y}=\lb\dfrac{p(x,y)}{p(x)p(y)}$ 拆开看，它其实在问一句话：
\textbf{“这两件事一起发生的频率，比‘各管各的’应该有的频率高多少倍？”}
\begin{itemize}[leftmargin=1.4em,itemsep=1pt,topsep=1pt]
  \item \textbf{分母 $p(x)p(y)$}：如果 $X$ 和 $Y$ \emph{完全没关系}
        （这种情况叫\textbf{统计独立}），那么同时发生的概率就应该是
        “$p(x)$ 乘 $p(y)$”。所以分母是\textbf{“没关系时本该有的概率”}，
        也就是你的\textbf{基准线}（瞎猜的命中率）。
  \item \textbf{分子 $p(x,y)$}：这两件事\textbf{实际}同时发生的概率。
  \item \textbf{比值}：实际 $\div$ 应该 $=1$ 表示“跟没关系一模一样”；
        $>1$ 表示“一起出现得异常频繁，它俩有牵连（正相关）”；
        $<1$ 表示“它俩老是互相躲着（负相关）”。
\end{itemize}
最后套上 $\lb$，是把“倍数”翻译成“bit”，这个下文马上解释。
\end{idea}

\subsubsection{为什么一定要套对数？（这里用了一个技巧）}

因为\textbf{对数能把乘法变成加法}，而“独立”正好是乘法。

具体地：如果 $x$ 和 $y$ 独立，则 $p(x,y)=p(x)p(y)$，于是联合自信息
$I(x,y)=-\lb p(x,y)=-\lb p(x)-\lb p(y)=I(x)+I(y)$——
\textbf{两件互不相干的事，信息量正好相加}（“明天太阳升起”和“明天下雨”
这两条消息合起来的信息量，就该等于两条相加）。取对数就是为了保住这个
“独立 $\Rightarrow$ 相加”的性质；不取对数，加法就凑不出来了。

再看比特的含义：比值 $=2$（实际发生的机会是瞎猜的 $2$ 倍）
$\Rightarrow\lb 2=1$ bit，正好等于“一次公平抛硬币”的信息量；
比值 $=4\Rightarrow\lb 4=2$ bit；比值 $=\frac14\Rightarrow\lb\frac14=-2$ bit。
所以你读这个数可以这样读：\textbf{“$1$ bit $=$ 一次公平抛硬币的把握”}。

\subsubsection{第一个数值例子：天气与带伞}

\begin{calc}{例 1：把天气表的四格全算出来}
学长给了一张真实点的 $2\times2$ 表（横向是带不带伞，纵向是下不下雨，
格子里的数是概率）。先看\textbf{边缘}（每行、每列加起来）：

\begin{center}
\small
\begin{tabular}{@{}lccc@{}}
\toprule
 & 带伞 $y=u$ & 不带伞 $y=\bar u$ & 边缘 $p(x)$ \\
\midrule
下雨 $x=r$ & $0.4$ & $0.1$ & $0.5$ \\
不下雨 $x=\bar r$ & $0.2$ & $0.3$ & $0.5$ \\
边缘 $p(y)$ & $0.6$ & $0.4$ & $1$ \\
\bottomrule
\end{tabular}
\end{center}

\noindent\textbf{四格逐个算}（分子是联合概率，分母是两个边缘概率之积）：
\[
\begin{aligned}
\text{第 1 格：}\ \I{r;u}&=\lb\frac{p(r,u)}{p(r)p(u)}=\lb\frac{0.4}{0.5\times0.6}
=\lb\frac{0.4}{0.3}=\lb\frac{4}{3}=\lb 4-\lb 3\\
&=2-1.58496=0.41504\approx0.415\bit ,\\
\text{第 2 格：}\ \I{r;\bar u}&=\lb\frac{p(r,\bar u)}{p(r)p(\bar u)}
=\lb\frac{0.1}{0.5\times0.4}=\lb\frac{0.1}{0.2}=\lb\frac12=-1\bit ,\\
\text{第 3 格：}\ \I{\bar r;u}&=\lb\frac{p(\bar r,u)}{p(\bar r)p(u)}
=\lb\frac{0.2}{0.5\times0.6}=\lb\frac{0.2}{0.3}=\lb\frac23\\
&=\lb 2-\lb 3=1-1.58496=-0.58496\approx-0.585\bit ,\\
\text{第 4 格：}\ \I{\bar r;\bar u}&=\lb\frac{p(\bar r,\bar u)}{p(\bar r)p(\bar u)}
=\lb\frac{0.3}{0.5\times0.4}=\lb\frac{0.3}{0.2}=\lb\frac32\\
&=\lb 3-\lb 2=1.58496-1=0.58496\approx0.585\bit .
\end{aligned}
\]
怎么读这四个数：
\begin{itemize}[leftmargin=1.4em,itemsep=1pt,topsep=2pt]
  \item 第 1 格 $+0.415$：下雨时带伞的机会是“不管天气瞎猜”的 $\frac43$ 倍——
        大爷那句“下雨”确实把你往“带伞”这边推了一把。
  \item 第 2 格 $-1$：\textbf{负数！}“下雨”这条消息让“不带伞”\textbf{更难}发生了
        （$0.4$ 降到 $0.2$），所以是 $-1$ bit 的“反向帮助”。负号不是算错。
  \item 第 3、4 格是一对反号：不下雨时“带伞”变难（$-0.585$）、
        “不带伞”变容易（$+0.585$），符合常识。
\end{itemize}
小结：$+0.415,\,-1,\,-0.585,\,+0.585$——有正有负，说明“带伞”这件事
有时帮我们猜天气、有时反而误导（不带伞的人里也有相当多是懒，
跟天气没关系）。
\end{calc}

\begin{calc}{例 2：掷两颗骰子——完全独立，答案就是 $0$}
$X$ 是第一颗骰子的点数，$Y$ 是第二颗。取 $x=6,y=6$：
\[
p(x,y)=\frac1{36},\quad p(x)=\frac16,\quad p(y)=\frac16
\ \Longrightarrow\
\I{x;y}=\lb\frac{1/36}{(1/6)\times(1/6)},
\]
\[
\Longrightarrow\ \I{x;y}=\lb\frac{1/36}{1/36}=\lb 1=0\bit .
\]
两颗骰子各管各的，“第二颗是几点”对“第一颗是几点”一分钱信息也不给，
这就是 $0$ 的含义。
\end{calc}

\begin{calc}{例 3：掷两颗骰子——有牵连时是正数}
练一个“看着没关系、其实强关联”的：设 $X$ 是\textbf{第一颗骰子的点数}
（取值 $1,\dots,6$），$Y$ 是\textbf{两颗骰子的点数之和}（取值 $2,\dots,12$）。
现在看两个事件：$x$ 表示“第一颗是 $6$ 点”，$y$ 表示“和是 $12$”。
注意“和是 $12$”只有一种可能：$6+6$。于是
\[
p(x=6)=\frac16,\qquad p(y=12)=\frac1{36},\qquad p(x=6,\,y=12)=\frac1{36},
\]
\[
\I{x=6;\,y=12}=\lb\frac{1/36}{(1/6)\times(1/36)}
=\lb\frac{1/36}{1/216}=\lb 6\approx2.585\bit .
\]
人话：一旦知道“和是 $12$”，第一颗必然是 $6$——这条消息几乎把答案
直接端到你手里，所以是很大的正数 $2.585$ bit。
顺带看一眼上界：事件 $x=6$ 自己的信息量是 $I(x=6)=-\lb\frac16=\lb 6\approx2.585$ bit，
这里 $\I{x;y}=2.585$ 正好顶到 $I(x)$，说明“全知道了”。
（与此对比，例 1 里概率高达 $0.4$ 的格子只值 $0.415$ bit，
概率只有 $0.0278$ 的这一格却值 $2.585$ bit——\textbf{概率大小和关联强度
根本不是一回事}，这一点问答部分还会说。）
\end{calc}

\subsubsection{和自信息的关系：三种等价的写法}

把 $p(x,y)=p(x|y)p(y)$ 代进定义，再把对数拆开（$\lb$ 把乘积变加法，
这又是本讲那个“技巧”），立刻得到
\[
\I{x;y}=\lb\frac{p(x|y)}{p(x)}
=\underbrace{-\lb p(x)}_{I(x)}-\underbrace{\bigl(-\lb p(x|y)\bigr)}_{I(x|y)}
=I(x)-I(x|y),
\]
其中 $p(x|y)$ 读作“\textbf{在 $y$ 已经发生的条件下} $x$ 发生的概率”，
$I(x|y)=-\lb p(x|y)$ 叫\textbf{条件自信息}（知道 $y$ 以后 $x$ 还剩多少信息量）。
同理
\[
\I{x;y}=I(y)-I(y|x)=I(x)+I(y)-I(x,y).
\]
人话翻译：\textbf{互信息 $=$ 原来需要的“惊讶度”，减去知道 $y$ 之后剩下的“惊讶度”。}
这已经是下一节“平均互信息”三条公式的雏形了——它们是同一件事
在“事件层面”和“随机变量层面”的两次出现。

% ==================================================================
\subsection{互信息可正、可负、可为零；以及它最可爱的性质：对称}

\begin{defbox}{符号法则（一定要记住这一条）}
\begin{center}
\small
\begin{tabular}{@{}lll@{}}
\toprule
情形 & 数学条件 & 人话 \\
\midrule
正 & $p(x|y)>p(x)$ & 知道 $y$ 之后，$x$ 更可能发生了（帮上忙） \\
零 & $p(x|y)=p(x)$，即 $p(x,y)=p(x)p(y)$ & 知不知道 $y$，对猜 $x$ 毫无影响 \\
负 & $p(x|y)<p(x)$ & 知道 $y$ 之后，$x$ 反而更不可能了（帮倒忙） \\
\bottomrule
\end{tabular}
\end{center}
注意“零”那一行的\textbf{充要}关系：对这一对事件而言，
$\I{x;y}=0\iff x$ 与 $y$ \textbf{独立}。为什么？因为 $\lb$ 只在自变量为 $1$ 时取 $0$。
\end{defbox}

\begin{idea}{为什么会“帮倒忙”（负数不是bug）}
例 1 里的 $\I{r;\bar u}=-1$ bit 可以这么理解：你已经知道“今天下雨”了，
按常识下雨更可能带伞，那么“不带伞”的可能性就从 $0.4$ 掉到 $0.2$。
一条消息让某件事\emph{更不可能}，它就在这件事上提供了负的信息量
——就像你听到“今天有暴雨”，会更确定“今天不适合晾被子”。

而“平均互信息”永远非负（后面“基本性质”一节证明）：\textbf{负的格子会被正的格子抵消}
（例 1 里 $-1$ 被 $+0.415$、$+0.585$ 抵掉了）。
所以考试里如果算出 $\I{X;Y}<0$，一定是算错了。
\end{idea}

\begin{defbox}{对称性}
\[
\I{x;y}=\lb\frac{p(x,y)}{p(x)p(y)}
=\lb\frac{p(y,x)}{p(y)p(x)}=\I{y;x}.
\]
人话：“知道下雨对猜带伞的帮助”和“知道带伞对猜下雨的帮助”是\textbf{同一个数}。
\end{defbox}

\begin{idea}{对称性为什么值得单独说？因为它最容易在考试里被写反}
互信息对称，但\textbf{熵不对称}：$H(X|Y)\ne H(Y|X)$，一般是不等的！
例 1 里 $H(X|Y)=0.87549$ 而 $H(Y|X)=0.84644$，两者差 $0.029$。
记住这一对“长得像但不是一回事”的量：
\begin{itemize}[leftmargin=1.4em,itemsep=1pt,topsep=1pt]
  \item $\I{X;Y}$：两个圆\textbf{重叠}的面积，对称；
  \item $H(X|Y)$：$X$ 那个圆里\textbf{没被 $Y$ 盖住}的面积，不对称。
\end{itemize}
\end{idea}

还有两个上界，考试常拿来做“不可能性”判断：因为 $p(x|y)\le1$，所以
\[
\I{x;y}=I(x)-I(x|y)\le I(x),\qquad
\I{x;y}\le I(y).
\]
人话：一件事带给你的关于 $x$ 的信息，不可能超过“$x$ 自己发生”所含的信息量
（例 3 里 $\I{x;y}=2.585=I(x)$ 就是“全知道”的临界点）。
这两条上界，就是后面“基本性质”一节里极值性 $\I{X;Y}\le\min\{\Hx{X},\Hx{Y}\}$ 的事件级原型。

% ==================================================================
\subsection{平均互信息：把每一格的账按“发生频率”加权}

事件互信息只管\textbf{一个格子}。但我们通常要问的是整体：
“变量 $X$ 和变量 $Y$ 到底有多相关？”于是把所有格子的账按
\textbf{它实际发生的频率}加权平均。这个“加权平均”就是数学期望 $\E[\cdot]$。

\begin{defbox}{定义（平均互信息，教材定义 3.1.3）}
\[
\I{X;Y}=\E\bigl[\I{x;y}\bigr]
=\sum_{x\in\mathcal X}\sum_{y\in\mathcal Y} p(x,y)\,\lb\frac{p(x,y)}{p(x)p(y)} .
\]
人话：\textbf{对每一格算一句“帮了多少”，再按那一格实际出现的概率
$p(x,y)$ 加权平均}。权重是 $p(x,y)$ 而不是 $p(x)p(y)$——
因为我们算的是“现实中真的发生过的那些情况”的平均感受。
\end{defbox}

\begin{qa}{为什么权重是 $p(x,y)$，不是 $p(x)p(y)$？}
因为期望的定义是“取值 $\times$ 该取值出现的概率”。$\I{x;y}$ 这个量
出现在“$x,y$ 同时发生”这个场合，而“$x,y$ 同时发生”的概率正是 $p(x,y)$。
用 $p(x)p(y)$ 当权重是在算一个\textbf{虚构}的平均：它按“假设两件事已经独立”
来分配频次，而不是按现实里的频次，于是“$X$ 常见还是少见”这个信息会被一起抹掉，
算出来就不再是“知道 $Y$ 到底帮我多少”了。
一句话：\textbf{对数里面已经用 $p(x,y)$ 算账，外面的权重也得用同一个频率，
内外一致才叫平均互信息}。
\end{qa}

还有一种“信道视角”的写法，把 $p(x,y)=p(x)p(y|x)$ 代进去：
\[
\I{X;Y}=\sum_{x}p(x)\sum_{y}p(y|x)\,\lb\frac{p(y|x)}{p(y)}
=\sum_{x}p(x)\sum_{y}p(y|x)\lb\frac{p(y|x)}{\sum_{x'}p(x')p(y|x')},
\]
其中 $p(y|x)$ 读作“\textbf{输入是 $x$ 时，输出变成 $y$ 的概率}”（后面章节的
“转移概率矩阵”就是它）。课本（式 3.1.16）之所以要写这一版，是因为它说明了：
\textbf{平均互信息完全由“输入分布 $p(x)$ $+$ 信道 $p(y|x)$”决定}——
这就是第 4、5 章信道容量讨论的出发点。

% ==================================================================
\subsection{平均互信息的三种等价写法（三条都要背，考试任选一条）}

\begin{defbox}{定理（互信息与其他熵的关系，教材定理 3.2.1、3.2.2、3.2.3）}
\[
\I{X;Y}=\Hx{X}-\Hx{X|Y}=\Hx{Y}-\Hx{Y|X}=\Hx{X}+\Hx{Y}-\Hx{X,Y}.
\]
\end{defbox}

三条公式的人话翻译（这段是本讲最值钱的部分，请逐字读）：

\begin{itemize}[leftmargin=1.6em,itemsep=3pt,topsep=2pt]
  \item \textbf{第一条 $\I{X;Y}=\Hx{X}-\Hx{X|Y}$：}
        \textbf{互信息 $=$ 原来的不确定性 $-$ 知道 $Y$ 之后剩下的不确定性。}
        这是最本色的一条，就是本讲开头那句话的公式版：
        “不确定性减少了多少”。也正因为 $\Hx{X|Y}\ge0$（剩下的不确定性
        不可能是负的），立刻得到 $\I{X;Y}\le\Hx{X}$。
  \item \textbf{第二条 $\I{X;Y}=\Hx{Y}-\Hx{Y|X}$：}
        \textbf{反过来问也一样}：知道 $X$ 之后，$Y$ 的不确定性减少量
        和上一条\textbf{完全相等}。这就是平均值层面的对称性
        （事件级对称 $\I{x;y}=\I{y;x}$ 加权平均后自然还对称）。
  \item \textbf{第三条 $\I{X;Y}=\Hx{X}+\Hx{Y}-\Hx{X,Y}$：}
        \textbf{两个各自的不确定性加起来，减去“合起来看”的不确定性。}
        为什么？如果 $X,Y$ 完全没关系，合起来看的不确定性正好是两者相加
        （$H(X,Y)=H(X)+H(Y)$），多出来的部分就是 $0$；两者越相关，
        “合起来看”越是省事（重复的信息只算一次），$H(X,Y)$ 就越小，
        差额 $\I{X;Y}$ 就越大。\textbf{所以 $\I{X;Y}$ 就是“重复计算的那部分不确定性”。}
\end{itemize}

\begin{idea}{韦恩图记忆法（教材图 3.2 就是它）}
画两个圆：左边圆面积 $\Hx{X}$，右边圆面积 $\Hx{Y}$，两圆并起来
$\Hx{X,Y}$，中间重叠部分是 $\I{X;Y}$。于是
\begin{itemize}[leftmargin=1.4em,itemsep=1pt,topsep=1pt]
  \item 左边“月牙” $=$ $\Hx{X|Y}$：$X$ 独有的不确定性（$Y$ 没帮你解决的那部分）；
  \item 右边“月牙” $=$ $\Hx{Y|X}$；
  \item 重叠 $=$ $\I{X;Y}$：双方共有的、重复的那部分不确定性；
  \item 并集 $=$ $\Hx{X,Y}$；
  \item 两圆不相交 $\iff$ 独立 $\iff$ $\I{X;Y}=0$。
\end{itemize}
考试时在草稿纸角上画这两个圈，$7$ 条形如“$\Hx{X|Y}=\Hx{X,Y}-\Hx{Y}$”
的关系式全都自动想得起来。
\end{idea}

\begin{calc}{顺手验证“等价”不是口号：用定义推第一条}
学长不想让你背，想让你看它为什么成立。从定义
$\I{X;Y}=\sum_{x,y}p(x,y)\lb\frac{p(x,y)}{p(x)p(y)}$ 出发，
把对数里的分式拆成 $p(x|y)/p(x)$（\textbf{技巧：}因为 $p(x,y)=p(x|y)p(y)$，
分式上下同除 $p(y)$ 即可）：
\[
\I{X;Y}=\sum_{x,y}p(x,y)\lb p(x|y)-\sum_{x,y}p(x,y)\lb p(x).
\]
第一项 $=\sum_{y}p(y)\sum_{x}p(x|y)\lb p(x|y)=-\Hx{X|Y}$；
第二项 $=\sum_{x}\Bigl(\sum_y p(x,y)\Bigr)\lb p(x)=\sum_x p(x)\lb p(x)=-\Hx{X}$。
两式相减即 $\I{X;Y}=\Hx{X}-\Hx{X|Y}$。$\square$
（看到没——\textbf{每一次“把互信息换成一堆 $H$ 的加减”都是这个套路}：
先把 $\lb$ 里的分式拆成“联合 $\div$ 边缘”，再让一组求和“收”成边缘分布。）
\end{calc}

% ==================================================================
\subsection{手把手：把天气表从头算到尾，三条公式对答案}

\begin{calc}{完整计算（$2\times2$ 表，算术不跳步）}
\textbf{已知}（就是例 1 那张表）：
\[
p(r,u)=0.4,\quad p(r,\bar u)=0.1,\quad p(\bar r,u)=0.2,\quad p(\bar r,\bar u)=0.3 .
\]
换算器（就是本讲开头要你背的那三个数）：$\lb 3=1.58496$，$\lb 5=2.32193$，$\lb 10=3.32193$。
由它们立马得到：$\lb 0.4=-1.32193$，$\lb 0.6=-0.73697$，
$\lb 0.1=-3.32193$，$\lb 0.2=-2.32193$，$\lb 0.3=-1.73697$。

\medskip
\noindent\textbf{步骤 1：先算边缘分布}（行、列分别相加）
\[
p(r)=0.4+0.1=0.5,\quad p(\bar r)=0.2+0.3=0.5,\quad
p(u)=0.4+0.2=0.6,\quad p(\bar u)=0.1+0.3=0.4 .
\]
$X$ 是公平的（一半一半），$Y$ 偏一点（$60\%$ 带伞）。

\medskip
\noindent\textbf{步骤 2：算三个熵}（都用 $H=-\sum p\lb p$）
\[
\begin{aligned}
\Hx{X}&=-0.5\lb 0.5-0.5\lb 0.5=0.5\times1+0.5\times1=1.00000\bit ,\\
\Hx{Y}&=-0.6\lb 0.6-0.4\lb 0.4
=0.6\times0.73697+0.4\times1.32193=0.97095\bit ,\\
\Hx{X,Y}&=-\bigl[\,0.4\lb 0.4+0.1\lb 0.1+0.2\lb 0.2+0.3\lb 0.3\,\bigr]\\
&=0.4\times1.32193+0.1\times3.32193+0.2\times2.32193+0.3\times1.73697\\
&=0.52877+0.33219+0.46439+0.52109=1.84644\bit .
\end{aligned}
\]

\medskip
\noindent\textbf{步骤 3：用三条公式分别算 $\I{X;Y}$，看结果一不一样}
\[
\begin{aligned}
\text{式一：}\ &\Hx{X|Y}=\Hx{X,Y}-\Hx{Y}=1.84644-0.97095=0.87549,\\
&\I{X;Y}=\Hx{X}-\Hx{X|Y}=1.00000-0.87549=0.12451\bit ;\\
\text{式二：}\ &\Hx{Y|X}=\Hx{X,Y}-\Hx{X}=1.84644-1.00000=0.84644,\\
&\I{X;Y}=\Hx{Y}-\Hx{Y|X}=0.97095-0.84644=0.12451\bit ;\\
\text{式三：}\ &\I{X;Y}=\Hx{X}+\Hx{Y}-\Hx{X,Y}
=1.00000+0.97095-1.84644=0.12451\bit .
\end{aligned}
\]
三条全给 $0.12451$，一bit不差。

\medskip
\noindent\textbf{步骤 4：用定义（四格加权）再算一遍，第四次对答案}
\[
\begin{aligned}
\I{X;Y}={}&0.4\times\lb\frac{0.4}{0.3}+0.1\times\lb\frac{0.1}{0.2}
+0.2\times\lb\frac{0.2}{0.3}+0.3\times\lb\frac{0.3}{0.2}\\
={}&0.4\times0.41504+0.1\times(-1)+0.2\times(-0.58496)+0.3\times0.58496\\
={}&0.16602-0.10000-0.11699+0.17549=0.12452 .
\end{aligned}
\]
（最后一格的 $0.00001$ 差是四舍五入攒的，不是公式错。
考试里一般保留 $3\sim4$ 位有效数字即可。）

\medskip
\noindent\textbf{步骤 5：这个 $0.1245$ bit 说明什么？}
$\Hx{X}=1$ bit 是“完全猜不着”，$\Hx{X|Y}=0.8755$ bit 只降了
$0.1245$ bit（约 $12\%$）。人话：\textbf{光知道“我带不带伞”，
对你猜“下不下雨”的帮助很有限}——因为带伞这事还掺了“懒”
“教室有伞”等与天气无关的因素。对应地，
$\I{X;Y}=0.12451<0.97095=\Hx{Y}<1=\Hx{X}$，两个上界都满足。

\noindent\textbf{顺手做个“极值”对照}（这就是下一小节要说的性质）：
\begin{itemize}[leftmargin=1.4em,itemsep=1pt,topsep=1pt]
  \item \textbf{完全同步的表} $p(r,u)=0.6,\ p(\bar r,\bar u)=0.4$，另两格为 $0$
        （下雨一定带伞、不下雨一定不带）：两边都是 $\Hx{0.6}=0.97095$ bit，
        $\Hx{X|Y}=0$，故 $\I{X;Y}=0.97095\bit$，正好顶到上界。
  \item \textbf{独立表} $p(r,u)=0.3,\ p(r,\bar u)=0.2,\ p(\bar r,u)=0.3,\
        p(\bar r,\bar u)=0.2$：每格都等于两边缘之积，故 $\I{X;Y}=\lb1=0\bit$。
\end{itemize}
\textbf{注意独立表和本节的表边缘分布完全一样}（$X$ 都 $0.5/0.5$，$Y$ 都 $0.6/0.4$），
互信息却一个是 $0$、一个是 $0.12451$——\textbf{光看边缘定不出互信息，
关键看“怎么搭配”}。
\end{calc}

\begin{calc}{第二组数据：二元对称信道（BSC）——考试最爱的小例子}
$X$ 是一枚公平硬币（$0/1$ 各 $\frac12$）。$Y$ 是对 $X$ 的观测，
但每次有 $0.1$ 的概率看错（专业名称：错误转移概率 $p=0.1$）。
于是 $p(y|x)$ 两行都是 $(0.9,\,0.1)$ 的两个排列，且
\[
p(y=0)=0.5\times0.9+0.5\times0.1=0.5,
\]
即 $Y$ 也是公平硬币，$\Hx{Y}=1$ bit；而（两行条件分布相同，直接加权）
\[
\Hx{Y|X}=0.5\bigl[-\!0.9\lb0.9-0.1\lb0.1\bigr]\times2
=-0.9\lb 0.9-0.1\lb 0.1 .
\]
用换算器：$\lb 0.9=\lb\frac{9}{10}=\lb 9-\lb 10$，而 $\lb 9=2\lb 3=3.16993$，
故 $\lb 0.9=3.16993-3.32193=-0.15200$；又 $\lb 0.1=-3.32193$。所以
\[
\Hx{Y|X}=0.9\times0.15200+0.1\times3.32193=0.13680+0.33219=0.46899\bit ,
\]
\[
\I{X;Y}=\Hx{Y}-\Hx{Y|X}=1-0.46899=0.53101\bit .
\]
人话：$1$ bit 的不确定性，信道只送到 $0.531$ bit，剩下 $0.469$ bit
就是“噪声吃掉的部分”。这条式子后面第 5 章讲信道容量时天天用。
（顺便记住 $\Hx{0.1}\approx0.469$ 这个数，考试出现频率极高。）
\end{calc}

% ==================================================================
\subsection{互信息的四条基本性质（每条一句人话）}

\begin{defbox}{性质清单}
设 $X,Y$ 为离散随机变量，则
\begin{enumerate}[leftmargin=1.8em,itemsep=2pt,topsep=2pt]
  \item \textbf{非负性：}$\I{X;Y}\ge0$，且等号成立 $\iff$ $X,Y$ 统计独立。
        \emph{人话：知道点东西总比什么都不知道好；平均值上不可能“帮倒忙”。}
  \item \textbf{极值性：}$\I{X;Y}\le\min\{\Hx{X},\Hx{Y}\}$。
        \emph{人话：从 $Y$ 里能榨出的关于 $X$ 的信息，不可能超过 $X$ 本身的不确定性，
        也不可能超过“$Y$ 这个容器里总共装的东西”。}
  \item \textbf{独立为零：}$X,Y$ 独立 $\iff\I{X;Y}=0$（与性质 1 的等号条件同一件事）。
        \emph{人话：两件互不相干的事，互相提供不了任何信息。}
  \item \textbf{自信息：}$\I{X;X}=\Hx{X}$。
        \emph{人话：自己当然完全知道自己，不确定性一次全消。}
\end{enumerate}
\end{defbox}

\subsubsection{非负性的证明（只用 $\ln u\le u-1$ 这一个技巧）}

\begin{calc}{证明 $\I{X;Y}\ge0$，并拿到取等条件}
约定：$p(x,y)=0$ 的格子那项算 $0$（因为 $\lim_{t\to0}t\lb t=0$），
所以只需处理 $p(x,y)>0$ 的格子。\textbf{技巧：用一个初等不等式}
\[
\ln u\le u-1\quad(u>0)\ \Longleftrightarrow\ -\lb u\ \ge\ \frac{1-u}{\ln 2}
\quad(\text{因为 } \lb u=\tfrac{\ln u}{\ln 2},\ \ln 2\approx0.6931).
\]
把 $u=\dfrac{p(x)p(y)}{p(x,y)}$ 代进去：
\[
-\lb\frac{p(x)p(y)}{p(x,y)}\ \ge\ \frac{1}{\ln 2}\left(1-\frac{p(x)p(y)}{p(x,y)}\right).
\]
两边乘 $p(x,y)$（正数，不等号方向不变）并对所有格子求和：
\[
\I{X;Y}=\sum_{x,y}p(x,y)\lb\frac{p(x,y)}{p(x)p(y)}
\ \ge\ \frac{1}{\ln 2}\left[\sum_{x,y}p(x,y)-\sum_{x,y}p(x)p(y)\right].
\]
括号里两个和都是 $1$（分别是“全表加起来 $=1$”和“边缘乘起来分块求和 $=1$”），
所以右边 $=\dfrac{1}{\ln 2}(1-1)=0$，即 $\I{X;Y}\ge0$。

\textbf{取等条件}：$\ln u=u-1$ 只在 $u=1$ 时成立，即每个非零格子都满足
$p(x,y)=p(x)p(y)$——这正是\textbf{独立}。反过来独立时每格比值都是 $1$，
每项都是 $\lb1=0$。证完。$\square$
\end{calc}

\subsubsection{极值性与另外两条}

\begin{itemize}[leftmargin=1.6em,itemsep=3pt,topsep=2pt]
  \item \textbf{$\I{X;Y}\le\Hx{X}$}：由 $\I{X;Y}=\Hx{X}-\Hx{X|Y}$，
        而条件熵 $\Hx{X|Y}\ge0$（\emph{它是一堆“$-\lb p$”的加权平均，每项都非负}），
        所以自动成立；取等 $\iff\Hx{X|Y}=0$，即知道 $Y$ 就等于知道 $X$。
  \item \textbf{$\I{X;Y}\le\Hx{Y}$}：把上一条里的 $X,Y$ 换个名字（因为对称）。
        两条合起来就是\textbf{极值性}：
        \[
        \I{X;Y}\le\min\{\Hx{X},\Hx{Y}\}.
        \]
  \item \textbf{$\I{X;X}=\Hx{X}-\Hx{X|X}=\Hx{X}-0=\Hx{X}$}：
        用第一条件公式：已知 $X$ 时“$X$ 还是什么”已经确定了，条件概率
        $p(x|x)=1$，故 $-\lb1=0$，也就是“剩下的不确定性 $=0$”。数值感受：BSC 例子里若不出错
        （$p=0$），则 $\I{X;Y}=\Hx{Y}=\Hx{X}=1$ bit，与 $\I{X;X}$ 一样大。
\end{itemize}

\begin{exam}{两条性质的高频用法}
\begin{enumerate}[leftmargin=1.8em,itemsep=2pt,topsep=2pt]
  \item 题目问“$X,Y$ 是否独立”$\Rightarrow$ 只看一个数：$\I{X;Y}$ 是否为 $0$
        （等价地 $p(x,y)$ 是否处处 $=p(x)p(y)$）。
  \item 题目要你证明某个量“$\le\Hx{X}$”$\Rightarrow$ 把那个量写成
        $\I{X;Y}$ 或 $\Hx{X}-\Hx{X|Y}$，再用条件熵非负。
  \item 让你“证明 $\I{X;Y}\ge0$”就是上面那段不等式，或者干脆用
        $\I{X;Y}=\Hx{X}-\Hx{X|Y}\ge0$（两种写法阅卷都认）。
\end{enumerate}
\end{exam}

% ==================================================================
\subsection{条件互信息与链式法则：把“知道”分成两步}

现实里的“知道”常常是一步一步来的：先知道 $Y$，再知道 $Z$。
条件互信息问的是：\textbf{在已经知道 $Y$ 的前提下，$Z$ 还能额外告诉你多少关于 $X$ 的事？}

\begin{defbox}{定义（条件互信息，教材式 3.3.3）}
\[
\I{X;Y|Z}=\sum_{z}p(z)\sum_{x,y}p(x,y|z)\,\lb\frac{p(x,y|z)}{p(x|z)\,p(y|z)} .
\]
人话：\textbf{把 $Z$ 固定成某一个值 $z$，在这一层里算一遍普通的互信息，
再按 $p(z)$ 加权平均}。所以条件互信息不是新东西，就是“分组算互信息，
然后组间加权平均”。三条等价写法照抄上一节（把每个 $H$ 都加上“$|Z$”）：
\[
\I{X;Y|Z}=\Hx{X|Z}-\Hx{X|YZ}=\Hx{Y|Z}-\Hx{Y|XZ}
=\Hx{X|Z}+\Hx{Y|Z}-\Hx{X,Y|Z}.
\]
同样有非负性 $\I{X;Y|Z}\ge0$，和对称性 $\I{X;Y|Z}=\I{Y;X|Z}$。
\end{defbox}

\begin{idea}{“分两步知道”的直觉}
$X$ 是一台机器的故障原因，$Y$ 是仪表读数，$Z$ 是维修手册。
\begin{itemize}[leftmargin=1.4em,itemsep=1pt,topsep=1pt]
  \item $\I{X;Y}$：光看仪表，能推断多少故障原因；
  \item $\I{X;Z|Y}$：\textbf{已经看过仪表之后}再翻手册，还能额外推断出多少；
  \item $\I{X;YZ}$：仪表 $+$ 手册一起用，总共能推断多少。
\end{itemize}
“一共的收获”$=$“仪表给的”$+$“看过仪表之后再翻手册额外给的”——
这句话本身就是加法的定义，不用证明。
把这句大白话写成公式，就是链式法则（下面这个式子）。
\end{idea}

\begin{defbox}{链式法则（教材式 3.3.9，作业 3.3 考的就是它）}
\[
\I{X;YZ}=\I{X;Y}+\I{X;Z|Y}\ =\ \I{X;Z}+\I{X;Y|Z}.
\]
人话：\textbf{（$Y$, $Z$ 一起提供的关于 $X$ 的信息）
$=$（$Y$ 提供的）$+$（已知 $Y$ 后 $Z$ 额外提供的）}。
两个式子都对，因为“先知道 $Y$ 还是先知道 $Z$”只是顺序不同，总量一样。
\end{defbox}

\begin{calc}{证明链式法则（只有三行，但每行都要会对）}
记 $\I{X;YZ}$ 就是“把 $(Y,Z)$ 当成一个整体”再用公式：
\[
\I{X;YZ}=\Hx{X}-\Hx{X|YZ}
=\underbrace{\bigl[\Hx{X}-\Hx{X|Y}\bigr]}_{\I{X;Y}}
+\underbrace{\bigl[\Hx{X|Y}-\Hx{X|YZ}\bigr]}_{\I{X;Z|Y}} .
\]
第一个方括号按定义就是 $\I{X;Y}$；第二个方括号正是“已知 $Y$ 之后，
再把 $Z$ 弄明白所减少的不确定性”，也就是 $\I{X;Z|Y}$（见上面条件互信息
的第一条等价写法）。$\square$

\textbf{注意那个“技巧”}：这里没有算任何概率，只是加一项减一项
（$-\Hx{X|Y}$ 和 $+\Hx{X|Y}$），然后按定义认领。证明熵关系的题都是这一招：
\textbf{加一项减一项，凑成定义的样子}。
\end{calc}

\begin{calc}{两个数值例子（一个“白给”，一个“惊人”）}
\textbf{例 A（冗余：第二份只是重复）}：$X$ 是公平硬币，$Y=X$，$Z=X$（两份副本）。
$\I{X;Y}=\Hx{X}=1$ bit；但知道 $Y$ 之后 $Z$ 一点新东西都没有，
所以 $\I{X;Z|Y}=0$；链式法则给出 $\I{X;YZ}=1+0=1$ bit。
人话：\textbf{同一个消息听两遍，第二遍不值钱。}

\textbf{例 B（协同：$1=0+1$，这就是作业 3.1 的标准答案）}：
$X,Y$ 是两枚\textbf{独立}的公平硬币，令 $Z=X\oplus Y$（异或：相同为 $0$，不同为 $1$）。
\begin{itemize}[leftmargin=1.4em,itemsep=1pt,topsep=1pt]
  \item 先算 $\I{X;Y}$：因为独立，$p(x,y)=p(x)p(y)$ 处处成立，所以 $\I{X;Y}=0$ bit。
  \item 再算 $\I{X;Y|Z}$：固定 $z=0$，此时 $(x,y)$ 只能是 $(0,0)$ 或 $(1,1)$，
        各占一半；$p(x|z=0)=\tfrac12$，$p(y|z=0)=\tfrac12$，
        $p(x,y|z=0)=\tfrac12$。于是这一组里
        $\I{x;y|z=0}=\lb\frac{1/2}{(1/2)(1/2)}=\lb2=1$ bit。
        固定 $z=1$ 同理也是 $1$ bit。加权平均：$\I{X;Y|Z}=1$ bit。
  \item 验证链式法则：$\I{X;YZ}=\I{X;Y}+\I{X;Z|Y}=0+\I{X;Z|Y}$。
        而由于 $X=Y\oplus Z$，知道 $(Y,Z)$ 就等于知道 $X$，
        所以 $\I{X;YZ}=\Hx{X}=1$ bit，从而 $\I{X;Z|Y}=1$ bit。
        于是 $1=0+1$，两边相等。
\end{itemize}
人话：\textbf{两枚硬币本来毫无关系（互信息 $0$），但一旦告诉你“它俩的异或”，
它们就锁死了（互信息 $1$）}。这说明“$X$ 与 $Y$ 独立”和
“已知 $Z$ 时 $X$ 与 $Y$ 独立”是两回事——这是下一节那个反例的伏笔。
\end{calc}

% ==================================================================
\subsection{数据处理：信息处理只会丢信息，不会凭空造信息}

\subsubsection{马尔可夫链是什么意思}

$X\to Y\to Z$ 构成\textbf{马尔可夫链}，意思是：
\textbf{$Z$ 只通过 $Y$ 和过去相连}，用条件概率写就是
\[
p(z|x,y)=p(z|y)\quad\text{对所有 }x,y,z ,
\]
读作：“已经知道 $Y$ 之后，$X$ 的取值对 $Z$ 没有任何额外影响”。
生活版：$X$ 是你写的作业，$Y$ 是拍照后的照片，$Z$ 是学长的转述。
照片坏了以后，学长只能照着照片说——\textbf{他说的只取决于照片，
不取决于你原件写了什么}（原件他没见过）。

\begin{defbox}{数据处理定理（教材习题 3.10，PPT 讲评重点）}
若 $X\to Y\to Z$ 是马尔可夫链，则
\[
\I{X;Z}\le\I{X;Y}.
\]
人话：\textbf{信息处理（拍照、压缩、转发、算平均）只会让“关于 $X$ 的信息”
变少或不变，绝不可能变多}——处理不创造信息，只损失信息。
\end{defbox}

\begin{calc}{三步证明（这题考试的写法就这样）}
\textbf{第 1 步}（链式法则，两种顺序都写出来）：
\[
\I{X;YZ}=\I{X;Y}+\I{X;Z|Y}=\I{X;Z}+\I{X;Y|Z}.
\]
\textbf{第 2 步}（用马尔可夫条件，把一项干掉）：
因为 $p(z|x,y)=p(z|y)$，所以“已知 $Y$ 时，$X$ 与 $Z$ 独立”，
从而 $\I{X;Z|Y}=0$。于是第 1 步的左边两个式子化为
\[
\I{X;Y}=\I{X;YZ}=\I{X;Z}+\I{X;Y|Z}.
\]
\textbf{第 3 步}（用条件互信息非负性收尾）：
$\I{X;Y|Z}\ge0$，故 $\I{X;Y}\ge\I{X;Z}$。$\square$

\textbf{记忆口诀}：\textbf{“链式法则 $+$ 马尔可夫杀一项 $+$ 非负性收尾”}，
三道工序，全章证明题的通用模板。
\end{calc}

\begin{calc}{数值感受：两次拍照，信息一路变少（BSC 级联）}
接着用前面“二元对称信道”那个例子：$X$ 公平硬币 $\to$ 拍照（错 $0.1$）得 $Y$ $\to$ 再拍一次
（错 $0.1$）得 $Z$。
\begin{itemize}[leftmargin=1.4em,itemsep=1pt,topsep=1pt]
  \item 第一步：$\I{X;Y}=1-0.46899=0.53101$ bit（前面算过）。
  \item 两次拍照合起来等价于“一次错多少”？做一次“翻错几次”的分类：
        只看错一次（第一次错且第二次对，或第一次对且第二次错）
        $=0.1\times0.9+0.9\times0.1=0.18$；两次都对 $=0.9\times0.9=0.81$；
        两次都错 $=0.1\times0.1=0.01$。从头到尾看错一次的概率就是 $0.18$，
        所以 $X\to Z$ 等价于一个错误率 $0.18$ 的 BSC。
  \item 于是 $\Hx{Z|X}=-[0.18\lb0.18+0.82\lb0.82]$。拆 $\lb$ 的\textbf{技巧：
        把小数写成“分母含 $50$ 的分数”}：
        $\lb0.18=\lb\frac{9}{50}=\lb9-\lb50=3.16993-5.64386=-2.47393$
        （其中 $\lb 9=2\lb 3$，$\lb 50=\lb2+2\lb5$），
        同理 $\lb0.82=\lb\frac{41}{50}=5.35755-5.64386=-0.28631$。代入：
        \[
        \Hx{Z|X}=0.18\times2.47393+0.82\times0.28631
        =0.44531+0.23477=0.68008\bit ,
        \]
        所以 $\I{X;Z}=1-0.68008=0.31992\bit$。
  \item 结论：$0.31992<0.53101$，确实变少了。
        （顺便得到一个规律：只要原来的错误率小于 $\frac12$，再独立地错一次
        总比原来错得更多，所以级联一次比一次糟——第二次纯属浪费胶卷。）
\end{itemize}
\end{calc}

\begin{qa}{“信息处理只会减少信息”——会不会有例外？}
\textbf{会——只要 $X\to Y\to Z$ 不是马尔可夫链，这个不等式就可能不成立。}
用的时候两条前提缺一不可：（1）$X\to Y\to Z$ 必须真的是条链；
（2）比的是“$\I{X;Z}$ 与 $\I{X;Y}$”这种	extbf{同一方向}的量，
不是拿条件互信息去比。
最著名的反例就是前面链式法则里的\textbf{例 B（异或）}：$X,Y$ 独立，
$\I{X;Y}=0$；但加上 $Z=X\oplus Y$ 之后，$\I{X;Z|Y}=1$ bit。
也就是说
\[
\I{X;Y}=0\ \ <\ \ \I{X;Z|Y}=1 .
\]
人话：\textbf{“条件化”可以让互信息变大（$0\to1$），也可以让它变小}；
但\textbf{顺着马尔可夫链往前传}（$X\to Y\to Z$ 这种方向）只会变小。
这两个方向千万别混——这是本章最容易丢分的一句话。
\end{qa}

% ==================================================================
\begin{exam}{考试怎么考（第 3 章只考这四类题）}
\textbf{题型 1：给联合分布表，求 $\I{X;Y}$。}（最高频，几乎必有一道）

\emph{答题就写这 5 行，一行都别省：}
\begin{enumerate}[leftmargin=1.8em,itemsep=1pt,topsep=1pt]
  \item 抄表并算出边缘分布 $p(x),p(y)$（并把两者加起来各验一次 $=1$）；
  \item 写出 $\Hx{X}$、$\Hx{Y}$、$\Hx{X,Y}$ 的数值（或写成 $-\sum p\lb p$ 带数字）；
  \item 写出所用公式，如 $\I{X;Y}=\Hx{X}+\Hx{Y}-\Hx{X,Y}$；
  \item 代入数字，写出结果与单位（\textbf{bit}）；
  \item 校验：$0\le\I{X;Y}\le\min\{\Hx{X},\Hx{Y}\}$，若超出范围立刻回头查表。
\end{enumerate}

\textbf{题型 2：证明熵与互信息的关系式}（如 $\Hx{X,Y,Z}=\Hx{X}+\Hx{Y|X}+\Hx{Z|XY}$、
$\I{X;YZ}=\I{X;Y}+\I{X;Z|Y}$ 之类）。\emph{套路固定}：先把所有 $\I{\cdot}$ 用
$\Hx{a}-\Hx{a|b}$ 或 $\Hx{a}+\Hx{b}-\Hx{a,b}$ 展开成 $H$ 的加减，
\textbf{加一项减一项凑定义}，让中间项相互抵消，最后只剩非负量。
\textbf{关键：每一步都要注明依据}（“由链式法则”“由条件熵非负”），
阅卷按步给分。

\textbf{题型 3：判断独立 / 求极值 / 判断某不等式成立。}
独立 $\iff\I{X;Y}=0\iff p(x,y)=p(x)p(y)$ 处处成立；
$\I{X;Y}$ 的最大值 $\min\{\Hx{X},\Hx{Y}\}$ 在“一方完全决定另一方”时取到。

\textbf{题型 4：数据处理定理 $X\to Y\to Z\Rightarrow\I{X;Z}\le\I{X;Y}$。}
照抄前面数据处理那三道工序，第 2 步一定要\textbf{写出马尔可夫条件的表达式}
$p(z|xy)=p(z|y)$，再说明“故 $\I{X;Z|Y}=0$”。

\medskip
\noindent\textbf{易错点（每一条都有人栽过）}
\begin{itemize}[leftmargin=1.6em,itemsep=1pt,topsep=1pt]
  \item $\Hx{X|Y}\ne\Hx{Y|X}$（一般不等，例 1 里是 $0.875$ 对 $0.846$），
        但 $\I{X;Y}=\I{Y;X}$（严格相等）。\textbf{对称的只有互信息，不是条件熵。}
  \item 别把 $\I{X;Y}$ 和 $\Hx{X|Y}$ 搞混：前者是两圆的\textbf{重叠}，
        后者是左圆\textbf{露出来的月牙}。方向也相反：
        $\I{X;Y}$ 越大，$\Hx{X|Y}$ 越小（$\I{X;Y}=\Hx{X}-\Hx{X|Y}$）。
  \item $\I{x;y}$（小写，单个事件）可正可负；$\I{X;Y}$（大写，平均）
        \textbf{恒非负}。算出负数必是算错，别顺着往下写。
  \item 分子写联合概率、分母写两个边缘概率的\textbf{乘积}；
        别写成 $\lb\frac{p(x|y)}{p(y)}$ 这种把分母漏掉一项的形式。
  \item 权重用 $p(x,y)$；忘了权重（直接把四个事件互信息加起来）是最常见的错。
  \item 单位是 bit（底 $2$）；题目若用 $\ln$，单位是 nat，别混。
  \item 单看某一格 $\I{x;y}$ 可能是负数，这不矛盾（例 1 里就有 $-1$ 和 $-0.585$）；
        但平均值 $\I{X;Y}$ 一定 $\ge0$，两类量别混。
  \item 要证独立，必须 $p(x,y)=p(x)p(y)$ \textbf{每一格}都成立（或直接算出
        $\I{X;Y}=0$）；只在某一格上相等说明不了什么。
\end{itemize}
\end{exam}

% ==================================================================
\begin{qa}{你可能想问}
\textbf{问 1：互信息和熵到底什么关系？可以互相算出来吗？}

答：\textbf{不能互相算，它们是“两个圆”的关系}。
$\Hx{X}$、$\Hx{Y}$ 是两个圆的面积，$\I{X;Y}$ 是重叠面积。
知道两个圆和一个重叠，还差一个“并集”才能确定全部，所以单靠 $\Hx{X}$ 和
$\Hx{Y}$ 定不出 $\I{X;Y}$：前面手算过的天气表和“独立表”$(0.3,0.2,0.3,0.2)$ 的
边缘分布\textbf{完全一样}（都是 $0.5/0.5$ 和 $0.6/0.4$），
但一个是 $0.12451$ bit，一个是 $0$ bit。
记住三句：$\I{X;X}=\Hx{X}$（完全相关时重叠 $=$ 全圆）、
独立时 $\I{X;Y}=0$（不重叠）、其余介于两者之间。
真要靠 $H$ 算 $I$，必须\textbf{再多知道一个} $H(X,Y)$（或 $H(X|Y)$、$H(Y|X)$ 之一）。

\textbf{问 2：为什么互信息不会超过 $\Hx{X}$？会不会出现 $\I{X;Y}>\Hx{X}$？}

答：看第一条公式 $\I{X;Y}=\Hx{X}-\Hx{X|Y}$。$\Hx{X|Y}$ 是一堆 $-\lb p(x|y)$
的加权平均，每个 $p(x|y)\le1$，所以 $-\lb p(x|y)\ge0$，平均当然 $\ge0$。
\textbf{“剩下的不确定性”不可能是负的}——你不可能“知道得比全部还多”。
极端情形 $\Hx{X|Y}=0$（知道 $Y$ 就等于知道 $X$）时取等 $\I{X;Y}=\Hx{X}$，
如“下雨一定带伞、不下雨一定不带”。同理 $\I{X;Y}\le\Hx{Y}$，
所以上限是 $\min\{\Hx{X},\Hx{Y}\}$。

\textbf{问 3：为什么不直接用 $p(x,y)$ 的大小判断相关性，非要绕一个比值加 $\lb$？}

答：因为 $p(x,y)$ 的数值\textbf{同时}被“$X$ 有多常见”和“$Y$ 有多常见”压着，
横向比较会误判。看本讲两个例子：
例 1 里 $p(r,u)=0.4$（挺大一个数），对应的互信息只有 $0.415$ bit；
例 3 里 $p(x=6,y=12)=\frac1{36}=0.0278$（小得可怜），对应的互信息却是
$2.585$ bit。\textbf{要是看概率大小，就会得出“天气和带伞比骰子更相关”
这种荒唐结论。}
比值 $\dfrac{p(x,y)}{p(x)p(y)}$ 把“各自多常见”除掉了（分母就是“没关系时
本该有的概率”），于是所有关联都被折算到同一把尺子上：
比值 $=1$ 就是“毫无关联”这条零线，$>1$ 正相关，$<1$ 负相关。
取 $\lb$ 则是为了让“独立 $\Rightarrow$ 相加”：独立时比值恒为 $1$，
$\lb1=0$，信息量正好可以一节一节往上加（前面解释对数时讲过）。
一句话：\textbf{比值负责“除掉常见程度的干扰”，$\lb$ 负责“把倍数翻译成可加的 bit”。}
\end{qa}

\clearpage
\section{连续型随机变量的熵：为什么它可以是个负数}

前面几讲算的熵都是\textbf{离散}随机变量的：掷骰子、抛硬币。这一讲换成\textbf{连续}随机变量：
身高、温度、电压、噪声。你会碰到一件让初学者发愣的事——\textbf{连续随机变量的熵可以是负数}：
$X$ 在 $[-1,1]$ 上均匀分布时，$\Hc{X^2}\approx-0.443\bit$。信息量怎么会是负的？
本讲把负号的来历讲透：先讲清"连续随机变量"（第一、二节），再给连续熵的定义与来历（第三节），
再算两个必背结果（第四节），然后手把手算一道作业原型题（第五节），最后是最大熵定理、
熵功率与答题模板（第六至九节）。零基础没关系，会四则运算就跟得上。

\subsection{第一节\quad 连续随机变量：为什么"取到某一个点"的概率是 0}

\subsubsection{先看离散，再看连续}

抛一枚硬币，记 $X=0$ 为反面、$X=1$ 为正面，$\p{X=0}=\p{X=1}=0.5$。记号 $\p{\cdot}$ 就是"概率函数"：
给它一个取值，它告诉你这个取值出现的可能性。离散情形有三个特征，下面对照的就是这三条：
(1) 取值\textbf{一个一个数得出来}；(2) 每个取值都有一份\textbf{大于 $0$} 的概率 $\p{x_i}>0$；
(3) 所有取值概率加起来恰好是 $1$，即 $\sum_i\p{x_i}=1$——概率就是"份额"。
连续情形的样子完全不同：取\textbf{某地成年男性的身高} $X$（厘米），问 $X$ 恰好等于
$170.0000000\cdots$ 厘米的概率是多少？答案是 $\mathbf{0}$，不是"很小"，是\textbf{严格等于零}。
因为 $170$ 后面可以跟无限多位小数，$170.0001$、$170.000100001$……有无限多种可能，
而总概率只有 $1$ 这一份，无限多种可能去分一份，每种拿到的只能是 $0$。

\begin{defbox}{核心区别：连续时只有"区间"才有概率}
离散时"$X$ 等于某个值"有概率 $\p{x_i}>0$；连续时"$X$ 恰好等于某个具体值"的概率\textbf{永远是 0}，
即 $\p{X=x}=0$。有概率的只有"$X$ \textbf{落在某区间里}"：
\[
P(a<X\le b)=\int_a^b f(x)\,dx ,
\]
这个 $f(x)$ 叫\textbf{概率密度函数}（简称\textbf{密度}）。
所以"身高在 169 到 171 厘米之间概率为 0.35"是有意义的话；
"恰好等于 170 的概率是 0.02"则毫无意义——任何单点概率都是 $0$。
\end{defbox}

\subsubsection{$f(x)$ 不是概率，而是"概率的密度"}

这是全讲最容易误解的地方：$f(x)$ 名字带"概率"，但它\textbf{本身不是概率}。先看一个请记一辈子的类比。

\begin{idea}{一根铁棍的"线密度"}
一根铁棍线密度 $2$ 千克/米。问：铁棍上\textbf{某一个几何点}有多少质量？答案是 $0$——
点没有长度，再密的棍子一个点也没有质量；但问"\textbf{从 3 米到 4 米}这段多重？"
答案是 $2\times1=2$ 千克。规律是\textbf{质量 $=$ 密度 $\times$ 长度}，
而且\textbf{密度可以大于 $1$}（$2$ 千克/米就大于 $1$），这不影响任何事。
翻译成概率语言：$f(x)$ 就是"概率的密度"，\textbf{概率 $=$ 密度 $\times$ 区间长度}
（区间很短、密度近似不变时）：某一点的"概率"是 $0$，但一段区间的概率可以很大。
\end{idea}

\begin{calc}{$f(x)>1$ 完全正常}
密度函数必须满足两个条件：
\[
\text{(1)}\ f(x)\ge0\ \text{对所有 }x ;
\qquad\qquad
\text{(2)}\ \int_{-\infty}^{+\infty}f(x)\,dx=1 .
\]
条件 (2) 的大白话是"曲线 $y=f(x)$ 下方的\textbf{总面积等于 $1$}"（总面积 $=$ 总概率）。
\textbf{注意条件里没有 $f(x)\le1$！}$f$ 可以等于 $4$、$100$——它是密度不是概率，不受 $1$ 限制。
例如 $X\sim U(0,\tfrac14)$，区间内密度是常数，由总面积 $=1$ 定出
$f(x)=\dfrac{1}{\tfrac14-0}=4$（$0<x\le\tfrac14$），其他处为 $0$。
$f(x)=4$ 看着吓人，但从不拿它当概率，只用它乘长度：
$P\bigl(0<X\le\tfrac18\bigr)=4\times\tfrac18=0.5$，是个合格的概率（总面积
$\int_0^{1/4}4\,dx=1$ 也正常）。\textbf{记住}：$f(x)$ 的单位是"每单位长度的概率"，
所以可以大于 $1$；真正的概率是\textbf{面积}，面积永远不超过 $1$。
\end{calc}

两种最常用的密度，后面反复用：
\[
X\sim U(a,b):\quad f(x)=\frac{1}{b-a}\ (a\le x\le b),\qquad
X\sim N(m,\sigma^2):\quad f(x)=\frac{1}{\sqrt{2\pi}\,\sigma}\exp\!\left(-\frac{(x-m)^2}{2\sigma^2}\right).
\]
均匀分布的意思是"落进长度相等的区间，概率也相等"；正态分布是中间高两边低的钟形曲线，
$m$ 是均值（钟形中心），$\sigma^2$ 是方差（胖瘦，越大越分散），$\exp(t)$ 就是 $\mathrm{e}^t$
（$\mathrm{e}\approx2.71828$）。两条曲线总面积都是 $1$，但正态的积分写不出简单表达式，
考试只需记住"总面积 $=1$、方差 $=\sigma^2$"这两条。

\subsection{第二节\quad 积分这件事，你只要会这么多}

连续熵里全是积分号，先把门槛抹平。\textbf{图像一}：$\int_a^b f(x)dx$ 就是"曲线 $y=f(x)$、
横轴、$x=a$、$x=b$ 围出来的面积"（横轴上方算正），$f$ 是概率密度时它等于 $P(a<X\le b)$。
\textbf{图像二}：把 $[a,b]$ 切成无数宽 $\Delta x$ 的细条，每条近似是"高 $f(x_i)$、宽 $\Delta x$"
的矩形、面积 $f(x_i)\Delta x$，全部加起来再让 $\Delta x\to0$ 就是积分。
\textbf{记住：积分就是求和的连续版}——第三节全靠这句话。

\begin{defbox}{本讲要用的标准积分（含主力积分的来历）}
\[
\int_a^b c\,dx=c\,(b-a)\quad(c\ \text{常数}),
\qquad
\int_0^1y^{\alpha}dy=\frac{1}{\alpha+1}\quad(\alpha>-1),
\qquad
\int_0^1\frac{dy}{\sqrt y}=2 ,
\]
\[
\boxed{\ \int_0^1 y^{a-1}\ln y\,dy=-\frac{1}{a^{\,2}}\ }\qquad(a>0).
\]
最后这条是本讲的\textbf{主力积分}，第五节全靠它。来历：用求导公式
$(y^{a}\ln y)'=a\,y^{a-1}\ln y+y^{a-1}$（"前导后不导 $+$ 前不导后导"）从 $0$ 积到 $1$ 得
$\bigl[y^{a}\ln y\bigr]_{0}^{1}=\int_0^1a\,y^{a-1}\ln y\,dy+\int_0^1y^{a-1}dy$；
左边 $=0$（$y=1$ 时 $\ln1=0$；$y\to0^+$ 时 $y^{a}\ln y\to0$，因为 $a>0$ 时幂函数压得住对数），
右边第二个积分 $=\dfrac1a$，于是 $0=a\int_0^1y^{a-1}\ln y\,dy+\dfrac1a$，即主力积分成立。
\textbf{本讲只用 $a=\tfrac12$}：$\int_0^1y^{-1/2}\ln y\,dy=-\frac{1}{(1/2)^2}=-4$。
注意这是对数 $\ln$（底 $\mathrm{e}$），而熵的定义用 $\lb$（底 $2$），两者差一个 $\ln2$，第五节会救命。
\end{defbox}

还有一个第三节要用的概念：\textbf{期望}，即"加权平均"，权重是密度：
$\E\bigl[g(X)\bigr]=\int_{-\infty}^{+\infty}g(x)f(x)\,dx$。平均身高是 $\E[X]$，
"离均值平方的平均"就是方差 $\E[(X-m)^2]=\sigma^2$。这些积分不必会算，只要知道它是加权平均。

\subsection{第三节\quad 连续熵（微分熵）的定义}

\subsubsection{定义，以及它和离散熵的血缘}

\begin{defbox}{定义（连续熵 / 微分熵）}
设连续型随机变量 $X$ 的密度为 $f(x)$，定义
\[
\Hc{X}=-\int_{-\infty}^{+\infty}f(x)\lb f(x)\,dx=\E\bigl[-\lb f(X)\bigr].
\]
\textbf{人话}：对每个 $x$ 先算 $-\lb f(x)$（"这个位置有多稀罕"），再用密度当权重做加权平均，
即"平均的稀罕程度"。\textbf{它是离散熵的"积分版"}：把 $\Hx{X}=-\sum_i\p{x_i}\lb\p{x_i}$
里的"求和 $\sum$"换成"积分 $\int$"、"概率 $\p{x_i}$"换成"密度 $f(x)dx$"即可，
所以也常记作 $H(X)$、$h(X)$，或叫\textbf{微分熵}。对照看：
\end{defbox}

\begin{center}
\small
\begin{tabular}{@{}l c c@{}}
\toprule
 & 离散 & 连续 \\
\midrule
描述工具 & 概率 $\p{x_i}$ & 密度 $f(x)$ \\
熵的定义 & $\Hx{X}=-\sum_i\p{x_i}\lb\p{x_i}$ & $\Hc{X}=-\displaystyle\int f(x)\lb f(x)dx$ \\
能否为负 & 不能（恒 $\ge0$） & \textbf{能} \\
\bottomrule
\end{tabular}
\end{center}

最后一行先给了答案，下面解释"为什么"——这是本章最常被问的一句话。

\subsubsection{它从哪来：切成小格子，再让格子无限窄}

用"积分是无限细的求和"这个图像，一步步推。这是理解全讲的钥匙，请慢慢读。

\begin{calc}{核心推导：切成宽 $\Delta x$ 的小格子}
\textbf{第 1 步（切格子）}：把取值范围切成一根根宽 $\Delta x$ 的小格子（像把棍子切成薄片），
第 $i$ 格是 $[a+(i-1)\Delta x,\ a+i\Delta x]$，取格内一点 $x_i$ 当代表。
\textbf{第 2 步（写概率）}：$X$ 落进第 $i$ 格的概率 $=$ 那根细条的面积 $\approx f(x_i)\Delta x$。
\textbf{第 3 步（当成离散问题算熵）}："落在哪个格子"是个离散随机变量，取值概率
$p_i=f(x_i)\Delta x$，套离散熵公式（第二步用 $\lb(uv)=\lb u+\lb v$）：
\begin{align*}
H_\Delta
=-\sum_{i}f(x_i)\Delta x\Bigl[\lb f(x_i)+\lb\Delta x\Bigr]
=-\sum_{i}f(x_i)\Delta x\,\lb f(x_i)\ -\ \lb\Delta x\sum_{i}f(x_i)\Delta x .
\end{align*}
\textbf{第 4 步（认亲戚）}：第一项是"密度 $\times$ $\lb$密度"再求和，$\Delta x\to0$ 时正是积分
$-\int f\lb f\,dx=\Hc{X}$，是个\textbf{有限值}；第二项里 $\sum_if(x_i)\Delta x$ 是所有细条面积之和
$=$ 曲线下总面积 $=1$。于是 $H_\Delta\approx\Hc{X}-\lb\Delta x$。
\end{calc}

关键就在 $-\lb\Delta x$。$\Delta x$ 是小于 $1$ 的正数（比如 $0.001$），故 $\lb\Delta x<0$、
$-\lb\Delta x>0$，且 $\Delta x$ 越小它越大：$H_\Delta=\Hc{X}+(-\lb\Delta x)$，
而 $-\lb\Delta x\to+\infty$（$\Delta x\to0$）。

\begin{defbox}{结论：真熵是 $+\infty$，只留有限的那一块；以及三条必记推论}
照离散熵的定义老实取极限，连续随机变量的\textbf{熵是无穷大}——这讲得通：要精确写下一个连续量
需要无限多比特，即"连续随机变量的潜在信息量无穷大"。Shannon 的处理很干脆：
\textbf{把趋于无穷的那一项扔掉，只留有限的那一项}，叫\textbf{微分熵}：
$\Hc{X}=\lim_{\Delta x\to0}\bigl(H_\Delta+\lb\Delta x\bigr)$。
所以必须清楚：$\Hc{X}$ \textbf{不是}"真正的无穷大熵"，而是"真熵减去无穷大常数
$-\lb\Delta x$ 后剩下的有限部分"。由此有三条必记推论：
\\[1pt]
\textbf{(1) 可以为负。} 因为 $\Hc{X}=H_\Delta+\lb\Delta x$，而 $\lb\Delta x\to-\infty$ 可以任意小，
剩下这块当然可能为负，例如 $X\sim U(0,\tfrac12)$ 时 $\Hc{X}=\lb\tfrac12=-1\bit<0$。
更直接的理由：单项 $-\lb f(x)$ 只在 $f(x)\le1$ 时 $\ge0$，而密度\textbf{可以大于 $1$}，
此时 $\lb f(x)>0$、单项为负（第五节的 $X^2$ 即是）。所以负的微分熵\textbf{不代表}"比确定还确定"，
只说"相对同单位下密度为 $1$ 的基准分布，这个变量更集中"。\\[1pt]
\textbf{(2) 只有"相对大小"有意义。} 那个无穷大常数 $-\\lb\Delta x$ \textbf{只跟格子宽度有关、
跟分布无关}，比较两个变量时自动抵消：$H_\Delta(X)-H_\Delta(Y)=\Hc{X}-\Hc{Y}$（同一 $\Delta x$）。
所以连续熵只能比"谁更不确定"，不能像离散熵那样说"这个信源值 3 比特"。
\textbf{数值还依赖单位}：同一个 $X$，用"厘米"记账与用"米"记账差 $\lb100$。\\[1pt]
\textbf{(3) 非负性没了，别的性质还在}（链式法则、条件熵不等式都保留，只是 $-\lb f$ 不再保底非负）。
\end{defbox}

\subsubsection{顺手记一个公式：$\Hc{g(X)}$（第五节要用）}

\begin{defbox}{变换公式：$\Hc{g(X)}=\Hc{X}+\E\bigl[\lb|g'(X)|\bigr]$}
设 $y=g(x)$ 是\textbf{一一映射}（严格单调：一直增或一直减），$g'(x)$ 是它的导数（斜率）。则
\[
\Hc{g(X)}=\Hc{X}+\E\bigl[\lb\bigl|g'(X)\bigr|\bigr].
\]
\textbf{人话}：变换后的熵 $=$ 原来的熵 $+$ "斜率绝对值的对数"的平均。斜率高，$x$ 的微小变化
被放大成 $y$ 的大变化，分布更"摊开"、密度更低、熵更大，加号就是这个意思。
\textbf{来历}：一一映射下 $f_Y(y)=\dfrac{f_X(x)}{|g'(x)|}$（分布被拉伸、密度变矮），代入定义
并换变量（$dy=|g'(x)|dx$）：$-\int f_X(x)\lb\dfrac{f_X(x)}{|g'(x)|}dx
=-\int f_X\lb f_X\,dx+\int f_X\lb|g'|\,dx=\Hc{X}+\E[\lb|g'(X)|]$。
\textbf{特例}：$g(x)=ax$ 给出 $\Hc{aX}=\Hc{X}+\lb|a|$，取 $a=2$ 就是"摊宽一倍、熵加 $1$ bit"，
与直觉一致；它也说明\textbf{连续熵对一一变换不保持不变}（离散熵是保持的）。
\end{defbox}

\subsection{第四节\quad 两个必背结果：均匀分布与正态分布}

\subsubsection{(a) 均匀分布：$\Hc{X}=\lb(b-a)$}

\begin{calc}{推导：$X\sim U(a,b)$}
\textbf{第 1 步（写密度）}：$f(x)=\dfrac{1}{b-a}$（$a\le x\le b$），区间外 $f=0$。
\textbf{第 2 步（代定义）}：区间外概率质量为 $0$、对积分无贡献（即离散里"$0\lb0=0$"的连续版），
只需在 $[a,b]$ 上积：
\[
\Hc{X}=-\int_{-\infty}^{+\infty}f\lb f\,dx=-\int_a^b\frac{1}{b-a}\,\lb\frac{1}{b-a}\,dx .
\]
\textbf{第 3 步（化简对数）}：$\lb\dfrac{1}{b-a}=-\lb(b-a)$，且被积表达式里\textbf{没有 $x$}、
整个是\textbf{常数}，用 $\int_a^bc\,dx=c(b-a)$：
\[
\Hc{X}=-\int_a^b\frac{1}{b-a}\Bigl[-\lb(b-a)\Bigr]dx
=\frac{\lb(b-a)}{b-a}\int_a^b dx=\frac{\lb(b-a)}{b-a}\cdot(b-a)=\lb(b-a).
\]
即 $\boxed{\ X\sim U(a,b)\Rightarrow\Hc{X}=\lb(b-a)\ }$。
\end{calc}

\textbf{怎么记}：\textbf{均匀分布的微分熵 $=$ 区间长度的对数}（长度翻倍，熵 $+1$ bit）。
三个立刻能用的数：$U(0,1)$ 时 $\Hc{X}=\lb1=0\bit$，恰在"零"这条分界线上；
$U(0,\tfrac12)$ 时 $\Hc{X}=\lb\tfrac12=-1\bit$，\textbf{负数}！区间短于 $1$ 个单位就为负；
$U(-1,1)$ 时 $\Hc{X}=\lb2=1\bit$，正是作业 2.10 第一问。

\subsubsection{(b) 正态分布：$\Hc{X}=\dfrac12\lb(2\pi\mathrm{e}\sigma^2)$}

推导只需两件事：\textbf{总面积 $=1$}、\textbf{方差 $=\sigma^2$}，两个积分都不必动手算。

\begin{calc}{推导：$X\sim N(m,\sigma^2)$，并代个数验证}
\textbf{第 1 步（写密度）}：
$f(x)=\dfrac{1}{\sqrt{2\pi}\,\sigma}\exp\!\left(-\dfrac{(x-m)^2}{2\sigma^2}\right)$。
\textbf{第 2 步（取以 $2$ 为底的对数）}：用 $\lb(uv)=\lb u+\lb v$、$\lb\mathrm{e}^{t}=t\lb\mathrm{e}$，
\[
\lb f(x)=-\lb\bigl(\sqrt{2\pi}\,\sigma\bigr)-\frac{(x-m)^2}{2\sigma^2}\,\lb\mathrm{e}
\qquad(\text{常数部分}+\text{含 }x\text{ 的部分}).
\]
\textbf{第 3、4 步（取期望，代入方差）}：$\E$ 是线性的，常数提出去；再由 $m$ 是均值得
$\E[(X-m)^2]=\sigma^2$：
\begin{align*}
\Hc{X}=\E\bigl[-\lb f(X)\bigr]
&=\lb\bigl(\sqrt{2\pi}\,\sigma\bigr)+\frac{\lb\mathrm{e}}{2\sigma^2}\E\bigl[(X-m)^2\bigr]\\
&=\lb\bigl(\sqrt{2\pi}\,\sigma\bigr)+\frac{\lb\mathrm{e}}{2\sigma^2}\cdot\sigma^2
=\lb\bigl(\sqrt{2\pi}\,\sigma\bigr)+\frac12\lb\mathrm{e}.
\end{align*}
\textbf{第 5 步（并成一个对数）}：用 $\lb\sqrt u=\tfrac12\lb u$、$\lb u+\lb v=\lb(uv)$，
\[
\Hc{X}=\frac12\lb\bigl(2\pi\sigma^2\bigr)+\frac12\lb\mathrm{e}
=\frac12\Bigl[\lb(2\pi\sigma^2)+\lb\mathrm{e}\Bigr]=\frac12\lb\bigl(2\pi\mathrm{e}\,\sigma^2\bigr).
\]
即 $\boxed{\ X\sim N(m,\sigma^2)\Rightarrow\Hc{X}=\frac12\lb(2\pi\mathrm{e}\sigma^2)\ }$。
\textbf{代个数}（考试爱让你算数值），取 $\sigma=1$：
\begin{align*}
\Hc{X}&=\frac12\lb(2\pi\mathrm{e}\cdot1)=\frac12\bigl(\lb2+\lb\pi+\lb\mathrm{e}\bigr)
=\frac12\bigl(1+1.6515+1.4427\bigr)\\
&=\frac12\times4.0942\approx2.047\bit ,
\end{align*}
其中 $\lb2=1$，$\lb\pi=\log_23.14159\approx1.6515$，
$\lb\mathrm{e}=\log_22.71828\approx1.4427$（即 $1/\ln2\approx1.4427$，很常用，建议背下）。
\end{calc}

\textbf{人话}：正态分布的熵\textbf{只跟方差有关、跟均值无关}——把钟形曲线平移（改 $m$）形状没变，
不确定度当然不变；压扁拉宽（改 $\sigma^2$）才变。两个结果一起记：$U(a,b)\Rightarrow\lb(b-a)$
（$b-a<1$ 时为负）；$N(m,\sigma^2)\Rightarrow\frac12\lb(2\pi\mathrm{e}\sigma^2)$
（$\sigma^2\lesssim0.0585$ 时为负）。可见"负熵"并不稀奇：分布足够集中（区间够短、方差够小）时，
微分熵就是负的。

\subsection{第五节\quad 手把手：$X\sim U(-1,1)$，求 $\Hc{X^2}$ 与 $\Hc{X^3}$}

这就是作业 2.10 的原型：\textbf{已知 $X$ 的分布，求 $g(X)$ 的熵}。
老师说"$\log$ 函数积分易求错"，所以下面每一步算术都写出来，一个数不跳。
第一问套第四节 (a)：$\Hc{X}=\lb(b-a)=\lb\bigl(1-(-1)\bigr)=\lb2=1\bit$
（用定义验也对：$f=\tfrac12$，$\Hc{X}=-\int_{-1}^1\tfrac12\lb\tfrac12\,dx=-\bigl(-\lb2\bigr)=1$）。

\subsubsection{第二问：$\Hc{X^2}$——关键是先求 $X^2$ 的密度}

\textbf{先排除一个坑：$X^2$ 不是均匀分布。} 有人想当然写"$f(x^2)=\tfrac12$，所以熵是 $1$"，这是错的。
$x$ 在 $[-1,1]$ 上均匀，但平方之后靠近 $0$ 的一段 $[-\epsilon,\epsilon]$ 被压成很小的 $[0,\epsilon^2]$，
靠近 $\pm1$ 的地方被摊开——所以 $X^2$ 在 $0$ 附近\textbf{特别密}。必须老实求密度，用\textbf{分布函数法}。

\begin{calc}{第 1、2 步：先求分布函数 $F_Y(y)$，求导得密度 $f_Y(y)$}
\textbf{第 1 步}：$X\in[-1,1]\Rightarrow Y=X^2\in[0,1]$，故只需看 $0\le y\le1$：
\begin{align*}
F_Y(y)&=P\bigl(X^2\le y\bigr)=P\Bigl(-\sqrt{y}\le X\le\sqrt{y}\Bigr)
\qquad(\text{开根号，}\pm\text{ 都留})\\
&=\int_{-\sqrt{y}}^{\ \sqrt{y}}\frac12\,dx=\frac12\Bigl[\sqrt y-(-\sqrt y)\Bigr]
=\frac12\cdot2\sqrt y=\sqrt y .
\end{align*}
读法：$X$ 落在 $[-\sqrt y,\sqrt y]$ 里，区间长 $2\sqrt y$，密度 $\tfrac12$，概率 $=\tfrac12\times2\sqrt y=\sqrt y$。
\\[1pt]
\textbf{第 2 步}："密度是分布函数的导数"，而 $(\sqrt y)'=\dfrac{1}{2\sqrt y}$：因为 $\sqrt y=y^{1/2}$，
幂函数求导 $(y^{n})'=ny^{n-1}$，取 $n=\tfrac12$ 得 $\tfrac12y^{-1/2}$。于是
\[
f_Y(y)=\frac{d}{dy}\sqrt y=\frac{1}{2\sqrt y}=\frac12y^{-1/2}\quad(0<y\le1),
\qquad f_Y(y)=0\ (\text{其他}).
\]
\textbf{自检}：$\displaystyle\int_0^1\frac{dy}{2\sqrt y}=\frac12\Bigl[2\sqrt y\Bigr]_0^1=1$，总面积对了。
另外 $y\to0^+$ 时 $f_Y(y)\to+\infty$，这就是"在 $0$ 附近特别密"，也正是后面熵为负的根源。
\end{calc}

\begin{calc}{第 3 步：代定义算 $\Hc{X^2}$（算术一步不跳）}
\[
\Hc{X^2}=-\int_0^1f_Y(y)\lb f_Y(y)\,dy
=-\int_0^1\frac{1}{2\sqrt y}\,\lb\left(\frac{1}{2\sqrt y}\right)dy .
\]
\textbf{先化简对数}：$\dfrac{1}{2\sqrt y}=\dfrac12y^{-1/2}$，用 $\lb\tfrac12=-1$、
$\lb y^{\beta}=\beta\lb y$（这里 $\beta=-\tfrac12$）：
\[
\lb\left(\frac{1}{2\sqrt y}\right)=\lb\Bigl(\tfrac12y^{-1/2}\Bigr)
=\lb\tfrac12+\lb\bigl(y^{-1/2}\bigr)=-1-\tfrac12\lb y .
\]
代回去，外面的负号把它翻正：
\[
\Hc{X^2}=\int_0^1\frac{1}{2\sqrt y}\Bigl[1+\frac12\lb y\Bigr]dy
=\underbrace{\frac12\int_0^1y^{-1/2}dy}_{\text{第一部分}}
+\underbrace{\frac14\int_0^1y^{-1/2}\lb y\,dy}_{\text{第二部分}} .
\]
\textbf{第一部分}：用 $\int_0^1y^{-1/2}dy=2$，得 $\dfrac12\times2=1$。
\textbf{第二部分}：用主力积分 $\int_0^1y^{a-1}\ln y\,dy=-\dfrac{1}{a^2}$，取 $a=\tfrac12$ 得
$\int_0^1y^{-1/2}\ln y\,dy=-\frac{1}{(1/2)^2}=-4$。但这里对数底是 $\mathrm{e}$，我们要 $\lb$（底 $2$），
换算 $\lb y=\dfrac{\ln y}{\ln2}$：
\[
\int_0^1y^{-1/2}\lb y\,dy=\frac{1}{\ln2}\int_0^1y^{-1/2}\ln y\,dy
=\frac{-4}{\ln2}=-4\lb\mathrm{e}\approx-5.7708 ,
\]
于是第二部分 $=\dfrac14\times\Bigl(-\dfrac{4}{\ln2}\Bigr)=-\dfrac{1}{\ln2}=-\lb\mathrm{e}\approx-1.4427$
（这里用了 $\dfrac{1}{\ln2}=\lb\mathrm{e}$，即 $1/0.6931=1.4427$）。
\textbf{合并}：
$\boxed{\ \Hc{X^2}=1-\dfrac{1}{\ln2}=1-\lb\mathrm{e}\approx1-1.4427=-0.4427\bit\ }$
约 $-0.443\bit$，\textbf{是负数}——这就是本章那个著名的负熵结果。
\end{calc}

\textbf{负号到底从哪来}：$\Hc{X^2}=1-\lb\mathrm{e}$ 中 $+1$ 来自密度曲线"正常的那部分"，
$-\lb\mathrm{e}\approx-1.4427$ 来自 $y\to0^+$ 时 $f_Y(y)\to+\infty$ 的区域——那里 $f_Y>1$、
$\lb f_Y>0$、单项 $-\lb f_Y$ 为负，把总和拉到 $0$ 以下。所以负号\textbf{不是算错}，
它反映"$X^2$ 在 $0$ 附近极度集中"；而按第三节，真正的信息量是无穷大，$\Hc{X^2}$ 只是去掉
无穷大常数后的\textbf{相对度量}。\textbf{别}把它读成"负的信息"。

\subsubsection{第三问：$\Hc{X^3}$——用变换公式更省力}

$g(x)=x^3$ 在 $[-1,1]$ 上严格递增（一一映射），可用第三节的变换公式。

\begin{calc}{算 $\Hc{X^3}$（两个小块）}
\textbf{第 1 步（求导）}：$g'(x)=3x^2$，故 $|g'(x)|=3x^2$，用 $\lb(uv)=\lb u+\lb v$：
\[
\Hc{X^3}=\Hc{X}+\E\bigl[\lb\bigl(3X^2\bigr)\bigr]
=1+\E\bigl[\lb3+\lb X^2\bigr]=1+\lb3+\E\bigl[\lb X^2\bigr]
\]
（$\lb3$ 是常数可提出期望；$\Hc{X}=1$ 来自第一问，$\lb3\approx1.5850$）。
\textbf{第 2 步（算 $\E[\lb X^2]$）}：按期望定义 $\E\bigl[\lb X^2\bigr]
=\int_{-1}^{1}\lb\bigl(x^2\bigr)\cdot\frac12\,dx$。$\lb(x^2)$ 是\textbf{偶函数}
（$x$ 换成 $-x$ 不变），对称区间上积分 $=2$ 倍的 $[0,1]$ 上积分，再乘 $\tfrac12$ 恰好抵消：
\begin{align*}
\E\bigl[\lb X^2\bigr]&=\frac12\cdot2\int_0^1\lb\bigl(x^2\bigr)dx
=2\int_0^1\lb x\,dx=\frac{2}{\ln2}\int_0^1\ln x\,dx
=\frac{2}{\ln2}\times(-1)=-\frac{2}{\ln2},
\end{align*}
即 $\E[\lb X^2]=-2\lb\mathrm{e}\approx-2.8854$（$\int_0^1\ln x\,dx=-1$ 是主力积分取 $a=1$ 的特例：
$-1/a^2=-1$）。这个数本章反复出现，值得单独记。
\textbf{第 3 步（合并）}：$1+\lb3=\lb2+\lb3=\lb6\approx2.5850$，故
\begin{align*}
\Hc{X^3}&=1+\lb3-2\lb\mathrm{e}=\lb6-2\lb\mathrm{e}
\approx2.5850-2.8854=\boxed{\ -0.3004\bit\ }
\end{align*}
约 $-0.300\bit$，也是负的。\textbf{对照}：直接求 $X^3$ 的密度再积分结果一样
（$f_W(w)=\dfrac{1}{6|w|^{2/3}}$，$w\in[-1,1]$），但变换公式省掉一半计算量。
\end{calc}

\begin{qa}{这道题最容易错的地方（三问答案：$1\bit$、$-0.443\bit$、$-0.300\bit$）}
\begin{enumerate}[leftmargin=1.6em,itemsep=1pt,topsep=1pt]
\item \textbf{以为 $X^2$ 是均匀分布}（$X$ 均匀 $\ne$ $X^2$ 均匀）。必须用分布函数法求
$f_Y(y)=\dfrac{1}{2\sqrt y}$ 再代定义，写成 $f(x^2)=\tfrac12$ 直接零分。
\item \textbf{$\sqrt y$ 求导漏掉 $\tfrac12$}：$\sqrt y=y^{1/2}$，$(y^n)'=ny^{n-1}$ 给出
$\tfrac12y^{-1/2}$。
\item \textbf{$\ln$ 与 $\lb$ 混用、丢了 $\ln2$}：主力积分是 $\ln$ 版（$-1/a^2$），熵的定义是 $\lb$ 版，
必须除以 $\ln2$（即乘 $\lb\mathrm{e}\approx1.4427$），丢了答案就不是 $-0.443$。
\item \textbf{$\E[\lb X^2]$ 记错}：它是 $-2\lb\mathrm{e}\approx-2.8854$，不是 $-2$、也不是 $-1$。
\item \textbf{把负号"做没了"}：算出负值就以为错、硬加绝对值——负值是对的。
\end{enumerate}
\end{qa}

\subsection{第六节\quad 两个最大熵定理（连续版）}

离散情形是"等概率时熵最大"。连续情形不能这么说：连续熵可为负，你把区间缩到极小，
熵就趋近 $-\infty$，无从谈"最大"。所以必须先\textbf{加约束}，再问"这个约束下哪种分布熵最大"。

\begin{defbox}{两个最大熵定理（记住"约束 $\to$ 分布 $\to$ 公式"三件套）}
\textbf{定理 A（峰值 / 幅度受限）}：若 $X$ 只能取 $[a,b]$ 内的值，则
$\Hc{X}\le\lb(b-a)$，等号当且仅当 $X$ 在 $[a,b]$ 上\textbf{均匀分布}。
（教材用对称区间 $[-M,M]$ 写成 $\Hc{X}\le\lb(2M)$，是同一件事：$b-a=2M$。）
\textbf{记忆}：\textbf{幅度受限 $\Rightarrow$ 均匀分布熵最大 $=\lb(b-a)$}。
直觉是：幅度被框住了，把概率\textbf{摊得最平}（处处密度一样）最不确定，
任何"偏爱某些位置"的分布都会降低不确定度。\\[2pt]
\textbf{定理 B（平均功率 / 方差受限）}：若 $X$ 的\textbf{方差固定为 $\sigma^2$}，则
$\Hc{X}\le\frac12\lb\bigl(2\pi\mathrm{e}\,\sigma^2\bigr)$，等号当且仅当 $X$ 服从\textbf{正态分布}
$N(m,\sigma^2)$。\textbf{记忆}：\textbf{方差受限 $\Rightarrow$ 正态分布熵最大
$=\frac12\lb(2\pi\mathrm{e}\sigma^2)$}。直觉是：方差固定 $=$ 总能量固定，此时正态分布
把能量分得最"不浪费"——均值附近铺开，尾部又不过分拖长。
\end{defbox}

\textbf{为什么重要？} 它们是一切"最坏情形"分析的起点：\textbf{手里信息不够，就假定熵取到上限}；
所以它们回答"只知道幅度 / 只知道功率时，最难压缩、最难传输的信号长什么样"——
通信里"最坏噪声是高斯噪声"根就在这里。考试最常见问法是"\textbf{在 XX 约束下用什么分布
能达到最大熵？最大熵是多少？}"，答案必须给\textbf{分布名 + 公式}。证明（Lagrange 乘子法）
只作简答：把"总量为 $1$、方差固定"用乘子 $\lambda_1,\lambda_2$ 拼进目标函数，对密度求极值得
$f(x)\propto\exp\bigl(-\lambda_2(x-m)^2\bigr)$，用 $\int f=1$ 定系数、用方差定
$\lambda_2=\frac{1}{2\sigma^2}$，即为正态密度。

\subsection{第七节\quad 熵功率：把熵换算成"等效噪声功率"}

\begin{defbox}{熵功率（一公式一解释）}
对任意连续随机变量 $X$，设其微分熵为 $\Hc{X}$，称 $\bar N=\dfrac{1}{2\pi\mathrm{e}}\,2^{\,2\Hc{X}}$
为 $X$ 的\textbf{熵功率}。\textbf{人话}：它是一个"翻译器"，把\textbf{熵}（比特）翻译成
\textbf{功率 / 方差}（物理单位），含义是"\textbf{一个方差为 $\bar N$ 的正态分布，
其微分熵恰好等于 $\Hc{X}$}"——即"这个信号的熵，相当于多大功率的高斯噪声"。
\end{defbox}

公式怎么来的？反过来解一次就清楚。由定理 B，方差为 $\bar N$ 的正态分布的熵是
$\frac12\lb(2\pi\mathrm{e}\bar N)$，令它等于 $\Hc{X}$，逐步解：
\[
\Hc{X}=\frac12\lb\bigl(2\pi\mathrm{e}\bar N\bigr)
\ \Longrightarrow\
2\Hc{X}=\lb\bigl(2\pi\mathrm{e}\bar N\bigr)
\ \Longrightarrow\
2^{\,2\Hc{X}}=2\pi\mathrm{e}\bar N
\ \Longrightarrow\
\bar N=\frac{1}{2\pi\mathrm{e}}2^{\,2\Hc{X}} .
\]
（第二步两边乘 $2$；第三步把 $\lb$ 拆掉：$\lb u=v\iff u=2^v$；第四步两边除以 $2\pi\mathrm{e}$。）
两个能用的观察：\textbf{(1)} 正态分布的熵功率就是它的方差——代入
$\Hc{X}=\frac12\lb(2\pi\mathrm{e}\sigma^2)$ 得 $2^{2\Hc{X}}=2\pi\mathrm{e}\sigma^2$，故
$\bar N=\sigma^2$，与教材一致。\textbf{(2)} 熵功率相等 $\iff$ 微分熵相等，所以它是
"不确定度的物理单位版本"，在信源编码、率失真函数里反复出现。

\subsection{第八节\quad 考试怎么考}

\begin{exam}{三类题 + 答题模板}
\textbf{题型一：给密度求 $\Hc{X}$。}\ 三步：抄定义 $\Hc{X}=-\int f\lb f\,dx$；
代 $f$ 并\textbf{化简对数}（$\lb\frac1u=-\lb u$、$\lb uv=\lb u+\lb v$）；
用标准积分算出数值、写单位 bit。\emph{检查点}：常数密度提到积分号外后，$\int dx$ 就是区间长度。\\[3pt]
\textbf{题型二：给 $X$ 的分布，求 $g(X)$ 的熵}（作业 2.10 型）。两条路：
\textbf{路 A（分布函数法，万能）}——$F_Y(y)=P(g(X)\le y)$ 解出 $y$ 的有效范围 $\to$ 求导得 $f_Y$
$\to$ \textbf{自检 $\int f_Y=1$} $\to$ 代定义积分；\textbf{路 B（变换公式，$g$ 单调时最快）}——
$\Hc{g(X)}=\Hc{X}+\E[\lb|g'(X)|]$，前提是 $g$ \textbf{严格单调}且 $\Hc{X}$ 先会算。
两路结果必须一致，不一致就逐行回查。\\[3pt]
\textbf{题型三：简答"为什么连续熵可为负 / 为什么定义微分熵"。}\ 模板四句：
\begin{enumerate}[leftmargin=1.6em,itemsep=1pt,topsep=1pt]
\item 把连续取值切成宽 $\Delta x$ 的小格，按离散熵算得 $H_\Delta=\Hc{X}-\lb\Delta x$；
\item $\Delta x\to0$ 时 $-\lb\Delta x\to+\infty$，所以按离散熵定义连续随机变量的熵是
\textbf{无穷大}、失去意义；
\item Shannon 去掉这个无穷大常数、只留有限的 $\Hc{X}=-\int f\lb f\,dx$，称为\textbf{微分熵}；
\item 因此 $\Hc{X}$ 只是\textbf{相对的不确定度度量}（差一个与分布无关的常数），
故可以为负（$f>1$ 使单项 $-\lb f<0$），也只能比较相对大小。
\end{enumerate}
\textbf{题型四：问"什么分布达到最大熵"。}\ 见"幅度 / 峰值受限"$\to$ \textbf{均匀}，$\Hc{X}=\lb(b-a)$；
见"平均功率 / 方差受限"$\to$ \textbf{正态}，$\Hc{X}=\frac12\lb(2\pi\mathrm{e}\sigma^2)$。
\end{exam}

\subsection{第九节\quad 几个很可能冒出来的问题}

\begin{qa}{密度函数能不能大于 $1$？}
\textbf{能，而且经常。} $f(x)$ 的单位是"每单位长度的概率"，就像线密度可以是 $2$ 千克/米。
$X\sim U(0,\tfrac14)$ 时 $f(x)=4$，而 $P\bigl(0<X\le\tfrac18\bigr)=4\times\tfrac18=0.5$ 依然合格。
受 $1$ 限制的是\textbf{面积}（概率），不是\textbf{高度}（密度）；硬性要求只有 $f(x)\ge0$ 和
$\int_{-\infty}^{+\infty}f(x)dx=1$。正因为 $f$ 可以大于 $1$，微分熵才可能为负。
\end{qa}

\begin{qa}{为什么连续熵可以为负，而离散熵不能？}
\textbf{离散熵非负}的原因很清楚：$0\le p\le1$ 保证 $\lb p\le0$，故 $-p\lb p\ge0$，非负项求和当然非负。
\textbf{连续熵失守}要两条一起说：(1) \textbf{记账方式}——$\Hc{X}$ 是"真熵 $+\infty$ 减去无穷大常数
$-\lb\Delta x$"剩下的有限部分（见第三节），减掉无穷大后当然可能落到 $0$ 以下；
(2) \textbf{直接看公式}——密度可以大于 $1$，此时 $\lb f>0$、单项 $-\lb f$ 为负，积分就有负贡献。
比如 $X\sim U(0,\tfrac12)$ 时 $\Hc{X}=\lb\tfrac12=-1\bit$；作业里 $\Hc{X^2}\approx-0.443\bit$ 同理。
\textbf{一句话}：负号不表示"负的信息"，只表示"相对同单位基准分布而言更集中"，
所以微分熵只做\textbf{差值比较}，不当绝对信息量读。
\textbf{附带提醒}：算 $g(X)$ 的熵时\textbf{不能}把 $x$ 换成 $g(x)$ 塞进 $f$ 里——换了随机变量，
分布就变了：$X\sim U(-1,1)$ 时 $X^2$ 不均匀（在 $0$ 附近特别拥挤，密度 $\dfrac{1}{2\sqrt y}$），
直接写 $f(x^2)=\tfrac12$ 得熵 $=1$，与正确答案 $-0.443$ 差 $1.443$ bit。
\end{qa}

\begin{idea}{这一讲的骨架（复习按这个顺序回忆）}
(1) 连续 $\ne$ 离散：单点概率 $=0$，只有区间有概率，用\textbf{密度}描述，记住 $f$ 可以 $>1$；
(2) 微分熵 $\Hc{X}=-\int f\lb f\,dx=\E[-\lb f(X)]$，是离散熵的积分版；
(3) 切格子得 $H_\Delta=\Hc{X}-\lb\Delta x$，真熵无穷大，只留有限项 $\Rightarrow$ 可为负、只比相对大小；
(4) 两个必背：$U(a,b)\to\lb(b-a)$；$N(m,\sigma^2)\to\frac12\lb(2\pi\mathrm{e}\sigma^2)$；
(5) 求 $g(X)$ 的熵：分布函数法，或 $\Hc{X}+\E[\lb|g'(X)|]$（$g$ 单调）；
(6) 最大熵：幅度受限 $\to$ 均匀；方差受限 $\to$ 正态；(7) 熵功率
$\bar N=\frac{1}{2\pi\mathrm{e}}2^{2\Hc{X}}$。
\end{idea}

\clearpage
\section{信道与信道容量：一条线最多能传多少信息}

\begin{idea}{这一讲只回答一个问题}
前几讲算的是“一条消息本身值多少 bit”。这一讲把消息真的送出去：\textbf{经过一条会出错的线之后，最多还能有多少 bit 到对面？}
这个“最多能到多少”就叫\textbf{信道容量}，记作 $C$。一句话分工：前几讲管“消息值多少”，这一讲管“\textbf{线只肯给你多少}”。
\end{idea}

你在电话里给对面报一串 0/1。你说“0”，对面大多数时候记下 $0$，但信号一抖，他也可能记成 $1$。
他记成什么你事前不知道，只知道\textbf{各种结果的概率}——这就是信道：不是理想的铜线，而是一句“发了这个，收到那个的可能性多大”。

\subsection{信道就是一张“发什么、收到什么”的概率表}

符号先讲清（每个符号第一次出现我都用大白话说一遍）：
\begin{itemize}[leftmargin=1.6em,itemsep=1pt]
  \item $X$：我\textbf{发出去}的符号，是随机变量（这一格发什么，事前我只知道概率）。
  \item $Y$：对面\textbf{收回来}的符号，也是随机变量。
  \item $\p{y|x}$：读作“\textbf{在已经发了 $x$ 这个前提下，收到 $y$ 的概率}”。那个竖线是中文“在……的条件下”，不是除法、也不是“或者”。
\end{itemize}

\begin{calc}{先学会读这张表（本讲的基本功）}
某条线，输入 $\{0,1\}$，输出也是 $\{0,1\}$，转移概率表是：
\begin{center}
\begin{tabular}{ccc}
\toprule
我发的 $x$ & 收到 $y=0$ 的概率 & 收到 $y=1$ 的概率 \\
\midrule
$x=0$ & $0.9$ & $0.1$ \\
$x=1$ & $0.2$ & $0.8$ \\
\bottomrule
\end{tabular}
\end{center}

\textbf{三条读表规则（务必背下来）：}
\begin{enumerate}[leftmargin=1.8em,itemsep=2pt]
  \item \textbf{横着读一行}：第一行读作“我发 $0$，对面收到 $0$ 的概率 $0.9$、收到 $1$ 的概率 $0.1$”；
        第二行读作“我发 $1$，对面收到 $0$ 的概率 $0.2$、收到 $1$ 的概率 $0.8$”。
  \item \textbf{每行加起来必须等于 $1$}：$0.9+0.1=1$，$0.2+0.8=1$。道理很土——发出去一个符号，对面总得收到点东西，所有可能铺开就是 $100\%$。
  \item \textbf{每列加起来不必等于 $1$}：第一列 $0.9+0.2=1.1$，第二列 $0.1+0.8=0.9$。因为“收到 $0$ 的总概率”取决于我发 $0$、发 $1$ 各占多少，而那是\textbf{我}定的，不是线定的。
\end{enumerate}

写成矩阵：
\[
Q=\begin{pmatrix}
\p{y=0|x=0} & \p{y=1|x=0}\\
\p{y=0|x=1} & \p{y=1|x=1}
\end{pmatrix}
=\begin{pmatrix} 0.9 & 0.1\\ 0.2 & 0.8\end{pmatrix}.
\]
人话：\textbf{一行 = 一种“我发了什么”，一格 = 一个条件概率}；第 $i$ 行第 $j$ 列就是“发了 $i$ 收到 $j$”的概率。
\end{calc}

\begin{defbox}{信道（本课只考“离散无记忆信道”）}
给定输入字母表 $A$、输出字母表 $B$，把所有 $\p{y|x}$ 列成的那张表（记 $Q=(q(y|x))$）就完整描述了一个信道，它叫\textbf{转移概率矩阵}。
“无记忆”是说：这一格出不出错与上一格无关，所以整串符号的概率可以直接连乘，我们只需要研究“一格”的信道。
\end{defbox}

\subsubsection{从这张表走到 $\I{X;Y}$：只有两步}

你已经知道“互信息 = 知道 $Y$ 之后，$X$ 的不确定性减少了多少”。写成公式：
\[
\I{X;Y}=\Hx{Y}-\Hx{Y|X}.
\]
人话：\textbf{看着 $Y$ 我还糊涂多少}，\textbf{减去}“本来就知道会糊涂多少”，剩下的才是真传过来的。
其中 $\Hx{Y|X}$ 是“在知道发了什么的条件下，对收到什么的平均不确定性”，算法是\textbf{每行各自算熵，再按我发各符号的概率加权平均}：
\[
\Hx{Y|X}=\sum_{x}\p{x}\,\Hx{Y|X=x}
=\sum_{x}\p{x}\left[-\sum_{y}\p{y|x}\lb\p{y|x}\right].
\]
人话：$\Hx{Y|X}$ 就是“\textbf{这条线平均每格给你添多少乱}”。它是线自己的毛病，与你怎么发无关——这一点后面要用。

\begin{calc}{手算示范：上面那张表，输入等概，求 $\I{X;Y}$}
设 $\p{X=0}=\p{X=1}=0.5$。（约定 $\Hx{p}=-p\lb p-(1-p)\lb(1-p)$，读作“一个概率为 $p$ 的二元事件的熵”，单位 bit；故 $\Hx{0.5}=1$，$\Hx{0}=\Hx{1}=0$。）

\textbf{第 1 步，全概率求输出的分布}：$\p{Y=0}=0.5\times0.9+0.5\times0.2=0.45+0.10=0.55$，$\p{Y=1}=1-0.55=0.45$。

\textbf{第 2 步，两处熵}。用换底 $\lb n=\dfrac{\ln n}{\ln2}$，$\ln2=0.6931$，查得 $\lb0.55=-0.8625$、$\lb0.45=-1.1520$、$\lb0.9=-0.1520$、$\lb0.2=-2.3219$、$\lb0.8=-0.3219$，于是
\[
\Hx{Y}=0.55\times0.8625+0.45\times1.1520=0.4744+0.5184=0.9928\bit,
\]
\[
\Hx{Y|X}=0.5\times\Hx{0.1}+0.5\times\Hx{0.2}=0.5\times0.4690+0.5\times0.7219=0.5954\bit .
\]
（其中 $\Hx{0.1}=0.1\times3.3219+0.9\times0.1520=0.4690$，$\Hx{0.2}=0.2\times2.3219+0.8\times0.3219=0.7219$。）

\textbf{第 3 步，相减}：$\I{X;Y}=0.9928-0.5954=0.3974\bit$。
人话：每传一格平均只有 $0.397$ bit 真到了对面，而理论上限是 $\Hx{X}=\lb2=1$ bit，线吃掉了大头。
\end{calc}

\begin{idea}{记住这个“手算三步走”}
\textbf{第一步}全概率算 $\p{Y=y}$；\textbf{第二步}每行求熵再加权得 $\Hx{Y|X}$，顺带算 $\Hx{Y}$；\textbf{第三步}相减。最容易错的是第一步忘了乘 $\p{x}$。
\end{idea}

\subsection{三个具体信道：把 $\I{X;Y}$ 一个一个算出来}

\subsubsection{无噪信道：收到就等于知道发了什么}

发 $0$ 必收 $0$，发 $1$ 必收 $1$，表是

\begin{center}
\begin{tabular}{ccc}
\toprule
我发的 $x$ & 收到 $y=0$ & 收到 $y=1$ \\
\midrule
$x=0$ & $1$ & $0$ \\
$x=1$ & $0$ & $1$ \\
\bottomrule
\end{tabular}
\end{center}

每行都是“一个 $1$ 一个 $0$”，每行熵 $=\Hx{1}=-1\cdot\lb1-0\cdot\lb0=0$（约定 $0\lb0=0$），故 $\Hx{Y|X}=0$，
\[
\I{X;Y}=\Hx{Y}-0=\Hx{Y}=\Hx{X},
\qquad
C=\max_{\p{x}}\Hx{X}=\lb2=1\bit .
\]
人话：\textbf{线一点乱都不添，发多少 bit 就收到多少 bit}；等概时 $X$ 最“不确定”（$\Hx{X}=1$ bit），所以容量是 $1$ bit。
换成 $4$ 个符号的无噪信道（$00,01,10,11$ 都不错），等概时 $\Hx{X}=\lb4=2\bit$，故 $C=2$ bit。
规律：\textbf{无噪信道 $C=\lb|X|$}（$|X|$ 是输入符号个数）。

\subsubsection{二元对称信道 BSC：以概率 $p$ 翻错}

“二元”= 发 $0/1$ 收 $0/1$；“对称”= 两种错误机会一样，都以概率 $p$ 翻错。

\begin{center}
\begin{tabular}{ccc}
\toprule
我发的 $x$ & 收到 $y=0$ & 收到 $y=1$ \\
\midrule
$x=0$ & $1-p$ & $p$ \\
$x=1$ & $p$ & $1-p$ \\
\bottomrule
\end{tabular}
\end{center}

两行是\textbf{同一组数 $(1-p,\,p)$ 换了顺序}——这个性质到“三个必背结论”那一节会变成一条能直接套的公式。
设 $\p{X=0}=p_0$，则 $\p{Y=0}=p_0(1-p)+(1-p_0)p=p+p_0(1-2p)$（这行只是告诉你输出分布跟 $p_0$ 有关）。
而两行熵相同，加权平均还是它自己：
\[
\Hx{Y|X}=p_0\Hx{p}+(1-p_0)\Hx{p}=\Hx{p}
=-\big[p\lb p+(1-p)\lb(1-p)\big].
\]
人话：\textbf{这条线每格固定吃掉 $\Hx{p}$ bit}，不管你发得均不均匀。于是
\[
\I{X;Y}=\Hx{Y}-\Hx{p}\le\lb2-\Hx{p}=1-\Hx{p},
\]
因为 $Y$ 只有两个取值，熵最大是 $\lb2=1$。什么时候取等？当 $\Hx{Y}=1$，即 $Y$ 等概；取 $p_0=0.5$ 时 $\p{Y=0}=p+0.5(1-2p)=0.5$，等号能取到。所以
\[
\boxed{\ C=\max_{p_0}\big[\Hx{Y}-\Hx{p}\big]=1-\Hx{p}\ }
\]
人话：容量 = 天花板 $1$ bit 减去“翻错带来的糊涂”$\Hx{p}$，最佳输入是等概。

\begin{calc}{代入 $p=0.1$：$C\approx0.531$ bit（算术一步不跳）}
\[
\Hx{0.1}=-\big[\,0.1\lb0.1+0.9\lb0.9\,\big],
\]
先查对数：$\lb0.1=\dfrac{\ln0.1}{\ln2}=\dfrac{-2.3026}{0.6931}=-3.3219$，$\lb0.9=\dfrac{\ln0.9}{\ln2}=\dfrac{-0.1054}{0.6931}=-0.1520$。代入：
\[
\Hx{0.1}=0.1\times3.3219+0.9\times0.1520=0.3322+0.1368=0.4690\bit,
\]
\[
C=1-\Hx{0.1}=1-0.4690=0.5310\bit\approx0.531\bit .
\]
人话：一条“错 10\%”的线，每格平均只送到 $0.531$ bit，近一半被线吃掉。这个 $0.469$ 请记住，考试很可能要你写。

\textbf{两个极端}：$p=0$ 时 $\Hx{0}=-0\lb0-1\lb1=0$，$C=1-0=1\bit$（退化成无噪信道）；
$p=0.5$ 时 $\Hx{0.5}=1$，$C=1-1=0\bit$。中间再记两个点：$p=0.01$ 时 $\Hx{0.01}=0.0808$，$C=0.919\bit$；
$p=0.11$ 时 $\Hx{0.11}=0.4999\approx0.5$，$C\approx0.500\bit$——\textbf{错误率 11\% 就把容量砍掉一半}。
\end{calc}

\subsubsection{二元擦除信道 BEC：以概率 $p$ 变成“擦除符号 ?”}

这条线不偷偷改，而是\textbf{大声告诉你“这位我丢了”}。输出有三个符号：$0$、$1$、擦除记号 $?$。

\begin{center}
\begin{tabular}{cccc}
\toprule
我发的 $x$ & 收到 $y=0$ & 收到 $y=?$ & 收到 $y=1$ \\
\midrule
$x=0$ & $1-p$ & $p$ & $0$ \\
$x=1$ & $0$ & $p$ & $1-p$ \\
\bottomrule
\end{tabular}
\end{center}

每行是 $(1-p,p,0)$ 或 $(0,p,1-p)$，熵都是 $\Hx{p}$，所以 $\Hx{Y|X}=\Hx{p}$。取 $\p{X=0}=0.5$，则
\[
\p{Y=?}=0.5p+0.5p=p,\qquad \p{Y=0}=\p{Y=1}=\frac{1-p}{2},
\]
\[
\Hx{Y}=\underbrace{\Hx{p}}_{\text{是不是被擦}}+\underbrace{(1-p)\cdot\lb2}_{\text{没被擦时是 }0\text{ 还是 }1}
=\Hx{p}+(1-p).
\]
（第二步用了熵的链式展开：先问“擦没擦”，没擦的那部分占 $1-p$、里面 $0/1$ 各半、熵为 $\lb2=1$ bit，再乘 $1-p$。）相减：
\[
\I{X;Y}=\big[\Hx{p}+(1-p)\big]-\Hx{p}=(1-p)\bit,
\qquad
\boxed{\ C=1-p\ }
\]
人话：\textbf{擦掉 $p$ 的比例，就正好损失 $p$ 的容量}，一点不冤。$p=0.1$ 时 $C=0.9\bit$；$p=0.5$ 时 $C=0.5\bit$；$p=1$ 时 $C=0$。

\begin{idea}{BSC 和 BEC 都“错 50\%”，容量为什么差这么多？}
BSC 的 $C=0$，BEC 的 $C=0.5\bit$。差别在于\textbf{擦除是诚实的错误，翻错是撒谎}：擦掉的那位你明知没用可以扔掉，
剩下 $1-p$ 的比例仍敢用；而翻错 50\% 时你连“哪一位错了”都不知道，收到 $0$ 时发 $0$ 和发 $1$ 一样可能，一格信息都没有。
\end{idea}

\subsection{信道容量的定义：我能选的只有“发什么符号的概率”}

上面三个例子里 $\I{X;Y}$ 都随“我怎么发”变化。为了把线本身的好坏单独拎出来：
\begin{defbox}{信道容量}
\[
C=\max_{\p{x}}\I{X;Y}
\]
人话：\textbf{在“发各符号的概率分布”这个我能自由选择的旋钮上，挑一个让传过去的信息最大的，那个最大值就是这条线的容量。}
\end{defbox}

三个必须分清的点：
\begin{enumerate}[leftmargin=1.8em,itemsep=2pt]
  \item \textbf{我能改的只有 $\p{x}$。}转移概率 $\p{y|x}$ 是线本身的性质，我改不了（要改就得换线），所以“取最大”只能对 $\p{x}$ 取。
  \item \textbf{$C$ 是线的属性，与我的输入分布无关。}同一条线，我发得聪明就能传到 $C$、发得笨就传得少，但这条线的 $C$ 一直是那个数。
        像一条公路的“最高时速”是路和车的能力，不取决于你今天开多快。
  \item \textbf{求 $C$ 时又要对 $\p{x}$ 取最大。}两句不矛盾：$C$ \textbf{是}一个不依赖 $\p{x}$ 的数，但它的\textbf{算法}是对 $\p{x}$ 求最大值。考试最爱在这儿挖坑。
\end{enumerate}

\begin{defbox}{最佳输入分布（考试会直接问这个词）}
使 $\I{X;Y}$ 取到最大值 $C$ 的那个输入分布 $\p{x}$，叫\textbf{最佳输入分布}。
$\I{X;Y}$ 关于 $\p{x}$ 是上凸的（中间高两头低，不会有两个一样高的山头），所以最佳分布只有一个。
\end{defbox}

\begin{defbox}{判别最佳分布的条件（好用，且能当验算）}
记 $I(X=k;Y)$ 为“\textbf{固定发第 $k$ 个符号时}这一格能传过去的互信息”（把“$x=k$ 板上钉钉”代入 $\I{X;Y}$ 就能算）。则最佳分布 $\p{x}^*$ 满足：
凡 $\p{x=k}>0$（用得上）的，$I(X=k;Y)=C$；凡 $\p{x=k}=0$（不用）的，$I(X=k;Y)\le C$。
\end{defbox}

人话：\textbf{最佳方案让“每个被用到的符号都同样划算”}，不划算的符号干脆不给概率。
它有两个用处：一是反手验算（算出 $p^*$ 后代回去看是否都等于 $C$），二是找 $p^*$（让各 $I(X=k;Y)$ 相等，解方程）——后面“非对称信道”那一节就这么干。

\subsection{三个必背结论}

\subsubsection{结论一：$C\le\lb|X|$（容量有天花板）}

推理一行：$\I{X;Y}\le\Hx{X}$（知道 $Y$ 能让 $X$ 的不确定性减少，但不会减少成负的、也不会比原来更多），而 $\Hx{X}\le\lb|X|$（$|X|$ 个符号的熵最大是等概的 $\lb|X|$）。所以
\[
\boxed{\ C\le\lb|X|\ }
\qquad
\text{人话：一格最多传 $\lb|X|$ bit，因为这一格我最多只能“有 $\lb|X|$ bit 的疑问”。}
\]
\textbf{什么时候取等}：$\I{X;Y}=\Hx{X}$，即 $\Hx{X|Y}=0$（收到 $Y$ 就完全确定 $X$），也就是\textbf{无噪信道}且输入等概。
\textbf{数值}：输入 $3$ 个符号则 $C\le\lb3=1.585\bit$；$4$ 个符号则 $C\le\lb4=2\bit$。
\textbf{常常够不着的例子}（课本例 5.1.2“复制信道”）：$26$ 个字母，每个以 $0.5$ 保持原样、以 $0.5$ 变成下一个字母，则 $C=\lb13=3.700\bit$，
而天花板是 $\lb26=4.700\bit$，少了整整 $1$ bit——正好是“到底有没有错一步”这个二元问题的 $1$ bit。

\subsubsection{结论二：强对称（对称）信道 $C=\lb|X|-\Hx{\text{一行}}$}

\textbf{什么叫对称}：转移概率矩阵的\textbf{每一行都是第一行换了顺序}（行可排列），并且\textbf{每一列元素之和相同}。二元对称信道、$K$ 元对称信道都属于这类。

为什么公式长这样：行可排列 $\Rightarrow$ 各行熵相同 $\Rightarrow$ $\Hx{Y|X}=\Hx{\text{一行}}$，与输入分布无关；
列和相等 $\Rightarrow$ 输入等概时输出也等概 $\Rightarrow$ $\Hx{Y}$ 顶到天花板 $\lb|X|$（这类信道输入输出符号数相同）。于是
\[
\boxed{\ C=\lb|X|-\Hx{\text{一行}},\qquad \text{最佳输入分布 = 等概分布}\ }
\]
人话：\textbf{容量 = 天花板减掉“一行那点平均糊涂”}，$\Hx{\text{一行}}$ 就是这条线平均每格自己添的乱。

\textbf{用它重算 BSC，检验公式靠不靠谱}：BSC 每行都是 $(1-p,p)$，这一行的熵就是 $\Hx{p}$，故 $C=\lb2-\Hx{p}=1-\Hx{p}$，
与上面 BSC 硬算的结果\textbf{完全一样}；$p=0.1$ 时 $C=1-0.4690=0.531\bit$ 也一样。公式和手算对上过一次，以后就能放心套。

\textbf{再看一个三符号的例子。}某 $3$ 元对称信道：发对的概率是 $0.7$，错成另外两个符号各 $0.15$，
因此它的每一行都是 $(0.7,\;0.15,\;0.15)$。由 $\lb0.7=-0.5146$、$\lb0.15=-2.7370$ 得
\[
\Hx{\text{一行}}=0.7\times0.5146+2\times0.15\times2.7370=0.3602+0.8211=1.1813\bit,
\]
\[
C=\lb3-1.1813=1.5850-1.1813=0.4037\bit ,
\]
最佳输入是三个符号各 $1/3$。注意它远小于 $1$ bit——三个符号看着热闹，其实穿不过 $0.41$ bit。

\subsubsection{结论三：准对称信道，等概输入也最佳}

\textbf{什么叫准对称}：行可排列（每行都是第一行换顺序），但\textbf{列和不相等}。此时输出符号可分成几组 $B_1,\dots,B_S$，使每组对应的列组成的子矩阵仍对称。容量公式：
\[
C=\sum_{s=1}^{S}N_s\lb\frac{|B_s|}{N_s}-\Hx{\text{一行}},
\qquad N_s=\text{一行里落进第 }s\text{ 组的概率（对任何输入都相同）}.
\]
人话：\textbf{把输出分成几个“小组”，每组的容量照对称公式算，再按概率加起来}。最佳输入仍是等概（题目只要求你会用这一条）。

\textbf{例子}（课本例 5.1.7，既有擦除又有误传）：每行是 $(1-p-q,\ q,\ p)$，取 $p=q=0.1$。
分两组：$B_1=\{0,1\}$ 是“收到真符号”（$|B_1|=2$，$N_1=(1-p-q)+p=1-q=0.9$），$B_2=\{?\}$ 是“擦除”（$|B_2|=1$，$N_2=q=0.1$），
$N_1+N_2=1$。代入，其中 $\lb\dfrac{2}{0.9}=\lb2.2222=1.1520$、$\lb\dfrac{1}{0.1}=\lb10=3.3219$、
$\Hx{0.8,0.1,0.1}=0.8\times0.3219+2\times0.1\times3.3219=0.2575+0.6644=0.9219$：
\[
C=0.9\times1.1520+0.1\times3.3219-0.9219=1.0368+0.3322-0.9219=0.4471\bit .
\]
\textbf{验算}（输入等概，直接算 $\Hx{Y}-\Hx{Y|X}$）：$\p{Y=0}=\p{Y=1}=\dfrac{1-p}{2}=0.45$、$\p{Y=?}=p=0.1$，故
$\Hx{Y}=1.3690$、$\Hx{Y|X}=0.9219$、$\I{X;Y}=0.4471\bit$，与公式结果一样，说明公式没记错。
\textbf{考试拿不准就用这招验算。}

\subsection{非对称信道：套不了公式的时候怎么办}

把最常见的坑摆出来。考虑 $2\times2$ 信道
\[
Q=\begin{pmatrix}1-\alpha & \alpha\\ \beta & 1-\beta\end{pmatrix}
\quad\text{取 }\alpha=0.1,\ \beta=0.3,\qquad
Q=\begin{pmatrix}0.9 & 0.1\\ 0.3 & 0.7\end{pmatrix},
\]
第一行是“发 $0$：$90\%$ 收对、$10\%$ 收成 $1$”，第二行是“发 $1$：$30\%$ 收成 $0$、$70\%$ 收对”。

\textbf{第一件事：判对称性。}两行 $(0.9,0.1)$ 与 $(0.3,0.7)$ \textbf{不是同一组数的重新排列}（数都不一样），列和 $0.9+0.3=1.2\ne0.1+0.7=0.8$，
所以它\textbf{既不对称也不准对称，前面的公式一个都不能用}。这就是“先判对称性”的意思。

\textbf{第二件事：把 $\I{X;Y}$ 写成一个只含 $p$ 的函数}，其中 $p=\p{X=0}$，$1-p=\p{X=1}$。每行熵是常数：
由 $\lb0.3=-1.7370$、$\lb0.7=-0.5146$ 得 $\Hx{0.3}=0.3\times1.7370+0.7\times0.5146=0.5211+0.3602=0.8813$，故
\[
\Hx{Y|X}=0.4690p+0.8813(1-p).
\]
输出分布\textbf{这次随 $p$ 变}（小心）：
\[
q:=\p{Y=0}=0.9p+0.3(1-p)=0.3+0.6p,
\qquad
\p{Y=1}=1-q=0.7-0.6p .
\]
于是
\[
\I{X;Y}=\Hx{q}-0.4690p-0.8813(1-p)
=\Hx{0.3+0.6p}-0.4690p-0.8813(1-p).
\]
人话：一个自变量、一条曲线，最大值只可能在两端或导数为 $0$ 处。

\textbf{第三件事：先扫描看大概。}（考试时间紧时，扫描表本身就是分。）

\begin{center}
\begin{tabular}{cccccc}
\toprule
$p$ & $0$ & $0.2$ & $0.4$ & $0.5$ & $0.8$ \\
\midrule
$q=0.3+0.6p$ & $0.30$ & $0.42$ & $0.54$ & $0.60$ & $0.78$ \\
$\Hx{q}$ & $0.8813$ & $0.9815$ & $0.9954$ & $0.9710$ & $0.7602$ \\
$\Hx{Y|X}$ & $0.8813$ & $0.7988$ & $0.7164$ & $0.6751$ & $0.5515$ \\
$\I{X;Y}$ & $0$ & $0.1826$ & $0.2790$ & $0.2958$ & $0.2087$ \\
\bottomrule
\end{tabular}
\end{center}

（$\Hx{Y|X}$ 那一行举例：$p=0.2$ 时 $0.4690\times0.2+0.8813\times0.8=0.0938+0.7050=0.7988$。表中数字四舍五入到四位，故末行与前几行相减可能差 $0.0001$。）
看得出：$p=0$ 与 $p=1$ 时 $\I{X;Y}=0$（只发一种符号，等于什么也没说），最大值在 $0.5$ 偏右一点。

\textbf{第四件事：求导定精确值。}只需一个工具：
\[
\frac{\mathrm{d}}{\mathrm{d}u}\Hx{u}=\lb\frac{1-u}{u},
\qquad\text{人话：熵函数在 } u \text{ 处的斜率是 } \lb\frac{1-u}{u}.
\]
（对 $\Hx{u}=-u\lb u-(1-u)\lb(1-u)$ 直接求导，两项合起来正好是上式。）链式法则给出 $\dfrac{\mathrm{d}}{\mathrm{d}p}\Hx{q}=\dfrac{\mathrm{d}q}{\mathrm{d}p}\lb\dfrac{1-q}{q}=0.6\lb\dfrac{1-q}{q}$，于是
\[
\frac{\mathrm{d}\I{X;Y}}{\mathrm{d}p}=0.6\ \lb\frac{1-q}{q}-0.4690+0.8813=0.6\ \lb\frac{1-q}{q}+0.4123 .
\]
令它等于 $0$：$\lb\dfrac{1-q}{q}=-\dfrac{0.4123}{0.6}=-0.6872$，即 $\dfrac{1-q}{q}=2^{-0.6872}=0.6211$。
解出 $1-q=0.6211q\Rightarrow q=\dfrac{1}{1.6211}=0.6169$，回代 $q=0.3+0.6p$：
\[
p^*=\frac{0.6169-0.3}{0.6}=\frac{0.3169}{0.6}=0.5281 .
\]
（二阶导数为负，所以这是最大值点而不是最小值点。）

\textbf{第五件事：算容量}，把 $p^*=0.5281$、$q^*=0.6169$ 代回。先算 $\Hx{0.6169}$：$\lb0.6169=-0.6969$、$\lb0.3831=-1.3842$，故
\[
\Hx{0.6169}=0.6169\times0.6969+0.3831\times1.3842=0.4299+0.5303=0.9602\bit,
\]
再把线吃掉的部分加权：$0.46900\times0.52812=0.24769$（发 $0$ 那行）、$0.88129\times0.47188=0.41586$（发 $1$ 那行），两项和 $0.66355$。相减：
\[
\boxed{\ C=0.96020-0.66355=0.29667\approx0.2967\bit,\qquad p^*\approx0.5281\ }
\]
对照：若硬算等概输入（$p=0.5$），只能得 $\I{X;Y}=0.2958\bit$，比最佳值少约 $0.001\bit$——
本例最优点离等概很近，但\textbf{严格说不是等概}。

\textbf{第六件事：用判别条件验算}（强烈建议写进卷面）。把 $p^*$ 代回，分别算两个符号的 $I(X=k;Y)$：
\begin{align*}
I(X=0;Y)&=0.9\lb\frac{0.9}{0.6169}+0.1\lb\frac{0.1}{0.3831}\\
&=0.9\times0.5449+0.1\times(-1.9378)=0.4904-0.1938=0.2966,
\end{align*}
\begin{align*}
I(X=1;Y)&=0.3\lb\frac{0.3}{0.6169}+0.7\lb\frac{0.7}{0.3831}\\
&=0.3\times(-1.0400)+0.7\times0.8695=-0.3120+0.6086=0.2966,
\end{align*}
两者都等于 $C\approx0.2967$（四位小数下都写成 $0.2966$，差别纯属四舍五入）；两个符号都在使用（$p^*$ 与 $1-p^*$ 都大于 $0$），正好符合判别条件。

\begin{idea}{这个例子要记住的三句话}
（1）\textbf{不对称就不能套公式}，先看两行是不是同一组数的排列；
（2）非对称时\textbf{把 $\I{X;Y}$ 写成只含 $p$ 的一元函数，再对 $p$ 求最大}（求导或扫描）；
（3）答案通常是“$p^*$ 略偏等概、$C$ 明显偏小”——不对称本身就是一种额外损失。
\end{idea}

\subsection{级联信道：串联两道关卡只会更差}

“级联”就是串起来：$X\to$ 信道 1 $\to Y\to$ 信道 2 $\to Z$。真实通信几乎都这样（发端、线路、收端各有瑕疵）。
\textbf{关键事实}：给定 $Y$ 之后，$Z$ 不再额外知道 $X$ 的任何事——$X$ 的信息只能“经 $Y$ 转手”。于是
\[
\I{X;Z}\le\I{X;Y},
\qquad
\text{人话：从终点 } Z \text{ 看到的 } X,\text{ 不会比中间站 } Y \text{ 看到的更多。}
\]
这就是有名的\textbf{数据处理不等式}：任何“后处理”（复印、排序、画图、压缩）都不能凭空造出信息。两条不等式一夹：
\[
\boxed{\ C_{\text{级联}}\le\min\{C_1,\ C_2\}\ }
\qquad
\text{人话：串起来不超过任何一环，瓶颈在最短的那块板。}
\]
证明只需两行：$\I{X;Z}\le\I{X;Y}\le C_1$，同时 $\I{X;Z}\le\I{Y;Z}\le C_2$（$X\to Y\to Z$ 也构成马尔可夫链），
故 $\I{X;Z}\le\min\{C_1,C_2\}$；对输入分布取最大，右边是常数，不等号保住。

\begin{calc}{两个 BSC 级联，手算一次}
单条线都是 BSC，翻错概率 $\varepsilon=0.1$，$C_1=1-\Hx{0.1}=0.531\bit$。串两条后，只有“先对后错”或“先错后对”才净翻错：
\[
\varepsilon_2=0.9\times0.1+0.1\times0.9=0.18,
\qquad\text{一般式 } \varepsilon_2=\frac{1-(1-2\varepsilon)^2}{2}=\frac{1-0.64}{2}=0.18 .
\]
（$(1-2\varepsilon)^2=0.8^2=0.64$ 叫“可靠性”，串级时相乘，所以错误率被拉向 $0.5$。）于是
\[
C_{\text{级联}}=1-\Hx{0.18}=1-0.6801=0.3199\bit<0.531\bit,
\]
其中 $\Hx{0.18}=0.18\lb\dfrac{1}{0.18}+0.82\lb\dfrac{1}{0.82}=0.18\times2.4739+0.82\times0.2863=0.4453+0.2348=0.6801$。
确实更差，符合 $C_{\text{级联}}\le\min\{0.531,0.531\}=0.531$。串 $N$ 条时 $\varepsilon_N=\frac{1-(1-2\varepsilon)^N}{2}\to0.5$，
故 $C_N\to1-\Hx{0.5}=0$——\textbf{无限串联等于完全失联}。
\end{calc}

\begin{idea}{为什么这个结论值得单独记}
它解释了一件反直觉的事：\textbf{多一道环节永远不是好事}。所以工程上能省一环就省一环，
也说明“反复转发、先解密再加密”救不了已经丢掉的信息。
\end{idea}

\subsection{高斯信道：连续世界的容量公式（只记结论和直觉）}

如果信号是连续电压、噪声是热噪声（高斯分布），就是\textbf{高斯信道}。结论漂亮得不像话：
\[
\boxed{\ C=\frac12\lb\!\left(1+\frac{P}{N}\right)\ \text{bit/符号}\ }
\]
$P$ 是信号平均功率（我能花多大力气说话），$N$ 是噪声平均功率（环境多吵），$\dfrac{P}{N}$ 叫\textbf{信噪比}，
读作“\textbf{信号功率比噪声功率}”：$P/N=1000$ 就是信号比噪声强 $1000$ 倍，$P/N=1$ 就是声音和噪音一样大。

\begin{calc}{代入三个数}
\textbf{（i）$P/N=1000$}：$C=\frac12\lb1001$，而 $\lb1001=\dfrac{\ln1001}{\ln2}=\dfrac{6.9088}{0.6931}=9.9672$，故 $C=\frac12\times9.9672=4.9836\approx4.98$ bit/符号。\\
\textbf{（ii）$P/N=1$}：$C=\frac12\lb2=\frac12\times1=0.5$ bit/符号。\\
\textbf{（iii）两个整数好算的}：$P/N=3$ 时 $C=\frac12\lb4=\frac12\times2=1$ bit；$P/N=15$ 时 $C=\frac12\lb16=2$ bit。
\end{calc}

\begin{idea}{为什么是 $\frac12$ 和 $\lb$（一句话直觉）}
\textbf{“为什么是 $\lb$”}：容量只能随信噪比\textbf{对数}增长。看上面的数：$P/N$ 从 $1$ 涨到 $3$（涨 3 倍）多 $0.5$ bit；
从 $3$ 涨到 $15$（涨 5 倍）只多 $1$ bit。\textbf{功率翻一番，只多 0.5 bit}——这就是“\textbf{对数瓶颈}”：靠加功率换速率，代价是指数级的。
\textbf{“为什么有 $\frac12$”}：连续信号取值无穷多，直接算熵没法比；$\frac12$ 是“每自由度”的折算（功率受限时高斯分布的熵最大，代进去系数正好是 $\frac12$）。
本讲只要求记公式、会代入、会解释 $P/N$ 是什么。
\end{idea}

\subsection{应试与答疑}

\begin{exam}{考试怎么考：见到信道题的固定动作}
\textbf{标准框架（按顺序写，一行都别省）：}
\begin{enumerate}[leftmargin=1.8em,itemsep=2pt]
  \item \textbf{抄转移矩阵}，检查每行和是否为 $1$（不是就先怀疑抄错题）。
  \item \textbf{判对称性}（这步决定后面用公式还是列方程）：每行都是第一行换顺序 $\Rightarrow$ 输入对称；
        行可排列\textbf{且}列和相等 $\Rightarrow$ 强对称（可直接套公式）；行可排列但列和不等 $\Rightarrow$ 准对称（等概最佳，按分组公式或直接算 $\Hx{Y}-\Hx{Y|X}$）；
        连行可排列都不是（如 $\left(\begin{smallmatrix}0.9&0.1\\0.3&0.7\end{smallmatrix}\right)$）$\Rightarrow$ \textbf{不许套公式}。
  \item \textbf{对称/准对称}：写 $C=\lb|X|-\Hx{\text{一行}}$（或分组公式）$\to$ 代数 $\to$ 补一句“最佳输入分布为等概分布，即 $p=1/|X|$”。
  \item \textbf{非对称}：设 $p=\p{X=0}$ $\to$ 全概率写 $q=\p{Y=0}$（检查 $\p{Y=1}=1-q$）$\to$ 写 $\I{X;Y}=\Hx{q}-\sum_x\p{x}\Hx{\text{该行}}$
        $\to$ 求导令 $0$（或给扫描表）$\to$ 解 $p^*$ $\to$ 代回得 $C$ $\to$ 用“各 $I(X=k;Y)=C$”验算。
  \item \textbf{级联题}：先写 $C_{\text{级联}}\le\min\{C_1,C_2\}$，再按需把矩阵相乘求等效信道；问为什么就一句“级联是马尔可夫链 + 数据处理不等式”。
  \item \textbf{高斯题}：$C=\frac12\lb(1+P/N)$，单位 bit/符号；先把 $P/N$ 算出来。
\end{enumerate}
\textbf{三个最常见的扣分点}：（1）不判对称性就套公式（本例会虚高到 $0.53$，正确是 $0.297$）；
（2）算 $\Hx{Y|X}$ 忘了乘 $\p{x}$，直接把两行熵相加；（3）说“$C$ 跟输入分布有关”——$C$ 与输入分布无关，只是求它时要对输入分布取最大。
\end{exam}

\begin{qa}{为什么 $C$ 是对“输入分布”取最大，而不是对“转移概率”取最大？}
因为\textbf{转移概率是线的物理性质，我改不了}——它由设备、线路、噪声决定，给定之后这条线的身份就定了；
而 $\p{x}$（我打算按什么比例发 $0$ 和 $1$）是我完全能控制的旋钮。$C$ 要回答“\textbf{这条线最好能给我多少}”，
所以要在“我能动的那一维”上取最好。反过来若对 $\p{y|x}$ 取最大，等于假设我换一条完全不同的线，跟这条线的容量无关；
而且 $\I{X;Y}$ 对 $\p{y|x}$ 是凸函数（不是凹的），最大值只会出现在“极端”信道上（比如无噪），毫无意义。
\end{qa}

\begin{qa}{$p=0.5$ 的 BSC 为什么容量是 $0$？我明明收到了东西啊。}
要分清“收到了符号”和“收到了\textbf{信息}”。$p=0.5$ 时 $\p{y=0|x=0}=\p{y=0|x=1}=0.5$、$\p{y=1|x=0}=\p{y=1|x=1}=0.5$，
即\textbf{知道收到什么，对“发的是什么”毫无提示}（看到 $Y$ 后 $X$ 的不确定性不减少）。公式也给 $C=1-\Hx{0.5}=1-1=0$。
物理上这叫“输出与输入统计独立”：线在替你掷硬币，结果跟你说了什么无关。对比同是“错一半”的 BEC 容量是 $0.5\bit$，
因为它诚实地说“这位丢了”，没丢的那一半你仍敢用。\textbf{“错”不可怕，“错得让你不知道错了”才可怕。}
\end{qa}

\begin{qa}{等概输入是最佳分布，那输出也一定等概吗？}
不一定。\textbf{对称信道}里两者同时成立，所以 $C=\lb|X|-\Hx{\text{一行}}$ 里的“天花板”项能拿到；
\textbf{非对称信道}就未必：本讲的例子最佳输入 $p^*=0.5281$（不是 $0.5$），对应输出 $q^*=0.3+0.6\times0.5281=0.6169$（也不是 $0.5$）。
一句口诀：\textbf{对称看等概，非对称要列方程}。
\end{qa}

\clearpage
\section{信源编码：用最短的平均码长把消息写下来}

\begin{idea}{这一讲在解决什么问题（先看地图）}
前几讲是\textbf{算账}：用熵 $\Hx{X}$ 量出一个信源到底有多少不确定性。这一讲开始干活：
硬盘和网线只认 $0$ 和 $1$，所以得把信源符号\textbf{换成一串 0/1}，这叫\textbf{编码}。
换法有无穷多种，凭什么说一种比另一种好？两条标准：
\textbf{(1) 能还原}——收到 0/1 串后必须一模一样地恢复原文（本章叫"无错编码"）；
\textbf{(2) 在能还原的前提下尽量省}——平均每个符号花的比特数尽量少。
于是全讲主线只有一句话：\textbf{在"能唯一还原"的码里把平均码长压到最小，而这条平均码长的下限就是熵。}
码的分类（什么叫能还原）$\to$ Kraft 不等式（什么码长可行）$\to$ 两个界限（熵是下限）
$\to$ Huffman（怎么把最小值做出来）$\to$ 熵率（有前后文时下限变成什么）。
\end{idea}

%==============================================================
\subsection{编码到底在干什么：把符号换成 0/1 串}

存一篇文章进硬盘，硬盘里没有"字"，只有能存高/低电平的格子，即只认 $0$ 和 $1$。
所以必须事先约定一张\textbf{对照表}：这个字对应这串 0/1，那个字对应那串 0/1；
存的时候查表写，读的时候反查表还原。最熟的例子是 ASCII 码：字母 \texttt{A} 记作 $1000001$（7 位）。

生活比方：\textbf{寄快递}。信源符号形状各异，不能直接扔上传送带，得先装进标准纸箱（码字）。
箱子有大有小（码字有长有短），\textbf{装得越紧凑，同样的车皮运得越多}——这就是"平均码长越小越好"。

\begin{defbox}{五个名词（定义 4.3.2--4.3.4）}
设信源符号集 $S=\{s_1,\dots,s_q\}$，$s_i$ 出现的概率为 $p_i=\p{s_i}$，$\sum_i p_i=1$；
码符号表（字母表）为 $A=\{a_1,\dots,a_r\}$。
\begin{itemize}[leftmargin=1.3em,itemsep=1pt,topsep=2pt]
  \item $r$ 叫码的\textbf{基数}：只用 $0,1$ 时 $r=2$（\textbf{二元码}），用 $0,1,2$ 时 $r=3$（\textbf{三元码}）。
  \item \textbf{码字} $w_i$：分给 $s_i$ 的那串码符号，如 $w_2=0110$。
  \item \textbf{码长} $l_i=\mathrm{len}(w_i)$：码字里码符号的个数，上例 $l_2=4$。
  \item \textbf{码} $C=\{w_1,\dots,w_q\}$：整张对照表。
  \item \textbf{编码规则} $(C,f)$，$f:S\to C$：把每个符号送到它的码字，
        且必须是\textbf{单射}（不同符号不能共用一个码字），否则一开始就没法还原。
\end{itemize}
\end{defbox}

\begin{defbox}{平均码长（本讲第一主角，必背）}
\begin{align*}
\bar L=\mathrm{Avelen}(C,f)=\sum_{i=1}^{q}p_i\,l_i=p_1l_1+p_2l_2+\cdots+p_ql_q .
\end{align*}
人话：\textbf{把各码字的长度按它出现的频繁程度加权平均}。出现得多的符号，
它的码长在平均值里说话更有分量。所以想把 $\bar L$ 压低，最有效的办法就是
\textbf{给常出现的符号配短码字}。注意 $\bar L$ 只与"概率 + 码长"有关，
跟码字具体长什么样（是 $0110$ 还是 $1011$）无关。
\end{defbox}

\begin{calc}{手算一个三符号的平均码长（不跳步）}
$p_1=0.5,\ p_2=0.3,\ p_3=0.2$（和 $=1$），二元码如下表。

\begin{center}
\begin{tabular}{@{}cccc@{}}
\toprule
符号 & 概率 $p_i$ & 码字 $w_i$ & 码长 $l_i$\\
\midrule
$s_1$ & $0.5$ & $0$  & $1$\\
$s_2$ & $0.3$ & $10$ & $2$\\
$s_3$ & $0.2$ & $11$ & $2$\\
\bottomrule
\end{tabular}
\end{center}

\textbf{第一步抄公式}：$\bar L=\sum_i p_il_i$。\textbf{第二步逐项相乘}：
$\bar L=0.5\times1+0.3\times2+0.2\times2=0.5+0.6+0.4=1.5$ bit/符号。
\textbf{第三步跟熵比}（用 $\Hx{X}=\sum_i p_i\lb\frac{1}{p_i}$，先算每个 $\lb\frac1{p_i}$）：
\begin{align*}
\lb\tfrac{1}{0.5}=\lb 2=1,\quad
\lb\tfrac{1}{0.3}=\lb 3.333\approx1.7370,\quad
\lb\tfrac{1}{0.2}=\lb 5\approx2.3219,
\end{align*}
$\Hx{X}=0.5\times1+0.3\times1.7370+0.2\times2.3219=0.5+0.5211+0.4644=1.4855$ bit/符号。
\textbf{第四步效率}：$\eta=\dfrac{\Hx{X}}{\bar L}=\dfrac{1.4855}{1.5}\approx99.03\%$，
即平均每符号只比理论下限多花 $1.5-1.4855=0.0145$ 比特。
\end{calc}

%==============================================================
\subsection{码的三层分类：非奇异 $\subset$ 唯一可译 $\subset$ 前缀码}

\subsubsection{先讲你天天遇到的问题：电话号码}
假设公司分机号是 $1$、$12$、$123$。你拨 "$12$"，交换机听见 "1" 就以为你要转 $1$ 号分机，
"2" 被当成下一个动作——\textbf{永远拨不到 $12$ 号}。这就是"码字不能互为前缀"的意思：
\textbf{短号码不能是长号码的开头}，否则一听到短号码就以为该停了，后面几位没法解释。

\begin{defbox}{三层定义（4.3.5--4.3.6，必背）}
\textbf{(1) 非奇异码}：$f$ 是单射，不同符号给不同码字。
人话：\textbf{一张表里没有两个相同的码字}。\\
\textbf{(2) 唯一可译码}：任一串码字 $c_{i_1}\cdots c_{i_k}$ 的分割方式唯一。
人话：\textbf{收到长 0/1 串后能一刀一刀切回原文，且切法只有一种}。\\
\textbf{(3) 前缀码（异字头码、即时码）}：$C$ 中没有任何码字是另一个码字的前缀。
人话：\textbf{一个码字一收完就知道它结束了}，可以立刻译出，不用等后面（"即时"、"无译码延时"）。
\end{defbox}

三者\textbf{层层包含}（选择题最爱考）：
\begin{align*}
\underbrace{\text{前缀码（即时码）}}_{\text{最强，无延时}}\;\subsetneq\;
\underbrace{\text{唯一可译码}}_{\text{能还原，可能要等}}\;\subsetneq\;
\underbrace{\text{非奇异码}}_{\text{最弱，只是不重样}} .
\end{align*}
一句话：\textbf{前缀码一定唯一可译且无延时；唯一可译码不一定是前缀码。}

\subsubsection{自己动手判四张码表（A/B/C/D 四型）}
同一信源、概率是"折半"型 $p=(0.5,\,0.25,\,0.125,\,0.125)$（和 $=1$），四张二元码表：

\begin{center}
\begin{tabular}{@{}ccccc|c@{}}
\toprule
符号 & 概率 $p_i$ & 码 A & 码 B & 码 C & 码 D\\
\midrule
$s_1$ & $0.5$   & $0$  & $0$  & $0$    & $0$\\
$s_2$ & $0.25$  & $1$  & $1$  & $01$   & $10$\\
$s_3$ & $0.125$ & $1$  & $00$ & $011$  & $110$\\
$s_4$ & $0.125$ & $1$  & $11$ & $0111$ & $111$\\
\midrule
平均码长 $\bar L$ & & $1.0$ & $1.25$ & $1.875$ & $1.75$\\
\bottomrule
\end{tabular}
\end{center}

\textbf{码 A}（$0,1,1,1$）：$s_2,s_3,s_4$ 共用码字 "$1$"，收到 "$1$" 分不清是哪个，
连\textbf{非奇异码}都不是。$\bar L_A=0.5+0.25+0.125+0.125=1$。\\
\textbf{码 B}（$0,1,00,11$）：四个码字互不相同，\textbf{非奇异} $\checkmark$。
但 "$11$" 既可切成 $s_4$、也可切成 $s_2s_2$，两种切法都合法，
故\textbf{不是唯一可译码}。$\bar L_B=0.5\times1+0.25\times1+0.125\times2+0.125\times2=1.25$。\\
\textbf{码 C}（$0,01,011,0111$）："$0$" 是 "$01$" 的前缀，故\textbf{不是前缀码}；
但它\textbf{是唯一可译的}——四个码字都只有一个开头的 $0$，后面 $1$ 的连续个数决定是哪个符号
（$0$ 个 $1$ 是 $s_1$，$1$ 个是 $s_2$，$2$ 个是 $s_3$，$3$ 个是 $s_4$）。
麻烦在于收到 "$011$" 时不能确定它是 $s_3$ 到头了、还是 $s_4=0111$ 的前三位，
\textbf{必须等下一个符号才敢拍板}——这就是\textbf{有译码延时}。
$\bar L_C=0.5\times1+0.25\times2+0.125\times3+0.125\times4=1.875$。\\
\textbf{码 D}（$0,10,110,111$）：$0$ 开头一个、$1$ 后接 $0/10/11$ 各不相同，无前缀关系，
是\textbf{前缀码}，收完即译。$\bar L_D=0.5\times1+0.25\times2+0.125\times3+0.125\times3=1.75$。

\begin{idea}{平均码长 $=1$ 的那个码（码 A）是不是最好的？不是，而且不能用}
表里 $\bar L$ 最小的恰是码 A（$=1$），只看数字 A 好像最好。但这个"最好"是\textbf{假的最好}：
\begin{itemize}[leftmargin=1.3em,itemsep=1pt,topsep=2pt]
  \item \textbf{它根本不能用}：三个符号共用 "$1$"，收到无法判断原文，连非奇异码都不满足。
        还原不回来的码，$\bar L$ 再小也没人要。
  \item \textbf{平均码长只在"唯一可译码"这一族里比才有意义}：把不能还原的码也算进来，
        $\bar L$ 可以任意小（所有符号都编成空串就 $\bar L=0$），但那等于什么都没写，信息全丢了。
\end{itemize}
\textbf{正确的问法是}：在唯一可译码（工程上用前缀码）这一族里，$\bar L$ 最小的是谁？
本表唯一可译的只有 C 和 D，D 更小（$1.75<1.875$）且无延时，
所以\textbf{四张表里最好的是 D}，而且它恰好达到熵（见下），一分钱没浪费。
\end{idea}

\textbf{顺手验证 D 为什么能到 $100\%$}：这个信源概率都是 $2$ 的负整数次幂，
\begin{align*}
\Hx{X}=0.5\lb 2+0.25\lb 4+0.125\lb 8+0.125\lb 8
=0.5+0.5+0.375+0.375=1.75=\bar L_D .
\end{align*}
效率 $\eta=100\%$。规律：当 $p_i=2^{-k_i}$ 时最优码长就是 $l_i=\lb\frac{1}{p_i}$（整数），
此时 $\bar L=\sum p_i\lb\frac1{p_i}=\Hx{X}$，一分冗余都没有。

%==============================================================
\subsection{Kraft 不等式：码字空间到底够不够用}

\subsubsection{人话版：一分四、四分十六}
拿一张纸代表"所有可能的 0/1 串"，只能靠\textbf{二分}来划分：中间剪一刀，左边串以 $0$ 开头、
右边以 $1$ 开头。于是\textbf{长度 1 的码字（如 "$0$"）占掉半张纸}（份额 $1/2$），
\textbf{长度 2 的码字（如 "$01$"）只占 $1/4$，长度 $l$ 的占 $\frac{1}{2^l}$}——
因为长度 $l$ 的二进制串共有 $2^l$ 个，挑走一个就是切走 $\frac{1}{2^l}$；$r$ 元码同理占 $\frac{1}{r^l}$。

\textbf{码字之间不能互相盖住}（某个码字一旦成了别人的前缀，就把别人那块纸整个吞了），
而纸总共只有 $1$ 张。所以\textbf{所有码字占的份额加起来不能超过 $1$}——
这就是 Kraft 不等式，一点也不神秘。

\begin{defbox}{Kraft 不等式（定理 4.3.1，必背）}
设码字长度为 $l_1,\dots,l_q$，则
\begin{align*}
\sum_{i=1}^{q}r^{-l_i}=r^{-l_1}+r^{-l_2}+\cdots+r^{-l_q}\ \le\ 1 .
\end{align*}
\textbf{必要性}：若 $C$ 是 $r$ 元即时码（前缀码），上式必成立——\textbf{真编得出来就一定没超支}。\\
\textbf{充分性}：若一组正整数 $l_1,\dots,l_q$ 满足上式，则\textbf{一定存在}具有这些码长的
$r$ 元前缀码——\textbf{没超支就一定编得出来}。\\
并且 McMillan 定理保证：对\textbf{唯一可译码}该不等式同样成立，所以这套判据对唯一可译码也管用。
\end{defbox}

\textbf{一句话翻译}：$\sum_i r^{-l_i}\le1$ 说的是"你打算用的这些码长，占的位置\textbf{不超过}整个码字空间"。
长码字占得少（$\frac{1}{r^l}$ 小），短码字占得多（$\frac{1}{r^l}$ 大），
所以"用太多短码字"必然超支——这正是"不可能让所有符号都用很短的码字"的原因。

\begin{calc}{判断二元/三元前缀码的存在性（考试常见小问）}
\textbf{(1) 码长 $\{1,2,3,3\}$ 能否构成二元前缀码？}（$r=2$，逐项算 $2^{-l}$）
\begin{align*}
\sum_i2^{-l_i}=\frac12+\frac14+\frac18+\frac18
=\frac{4}{8}+\frac{2}{8}+\frac{1}{8}+\frac{1}{8}=\frac{8}{8}=1\le1 .
\end{align*}
\textbf{能。}和恰为 $1$，说明纸被切得一点不剩，叫\textbf{满树码（完全码）}。
具体例子就是上面的码 D：$0,10,110,111$。

\textbf{(2) 码长 $\{1,2,2,3\}$ 能否？}
\begin{align*}
\sum_i2^{-l_i}=\frac12+\frac14+\frac14+\frac18
=\frac{4}{8}+\frac{2}{8}+\frac{2}{8}+\frac{1}{8}=\frac{9}{8}=1.125>1 .
\end{align*}
\textbf{不能。}超支 $1/8$ 张纸。直观：两个长度 $2$ 的码字吃掉 $1/2$，再加一个长度 $1$ 的又是 $1/2$，
这三个就占满 $1$ 了，第四个（哪怕长度 $3$）无处安放。

\textbf{(3) 三元码、长度 $\{1,2,2,2,2,2,3,3\}$}（$r=3$）：
\begin{align*}
\sum_i3^{-l_i}=\frac13+5\times\frac19+2\times\frac{1}{27}
=\frac{9}{27}+\frac{15}{27}+\frac{2}{27}=\frac{26}{27}\approx0.963\le1 .
\end{align*}
成立，故存在这样的三元前缀码（下面三元 Huffman 的例子正好给出它）。注意 $0.963<1$，
还剩一点空间没用——只有概率都是 $3$ 的负整数次幂时才可能刚好填满。
\end{calc}

\textbf{怎么用它判断"某码长分配可行不可行"}（三步走）：
(1) 数清是几元码，定下 $r$；
(2) 把每个 $\frac{1}{r^{l_i}}$ 算出来相加（\textbf{通分}成同一分母最不易错，上例都化成 $/8$、$/27$）；
(3) 和 $\le1$ $\Rightarrow$ \textbf{可行}（存在这样的前缀码）；和 $>1$ $\Rightarrow$ \textbf{不可行}；
若和 $=1$，补一句"满树码"，老师会加分。

\begin{qa}{Kraft 不等式成立，是不是就说明"这张码表"是前缀码？}
\textbf{不是，这是最常见的误解。} Kraft 定理管的是\textbf{"码长这组数字"}：它保证存在某个前缀码
其码长恰是这些数，但\textbf{不保证你手上这张具体码表}是前缀码。
反例（教材例 4.3.3）：$C=\{0,\,11,\,100,\,1100\}$，长度 $1,2,3,4$，
$2^{-1}+2^{-2}+2^{-3}+2^{-4}=\frac12+\frac14+\frac18+\frac1{16}=\frac{15}{16}\le1$
完全成立，但 "$11$" 是 "$1100$" 的前缀，所以这个码\textbf{不是即时码}。
\textbf{分工要记清}：判断一组\textbf{长度数字}用 Kraft；判断一张\textbf{具体码表}就逐个检查"谁是谁的前缀"。
\end{qa}

%==============================================================
\subsection{变长编码的两个界限：熵是理论下限}

\begin{defbox}{无失真变长编码定理（Shannon 第一定理，必背）}
对离散无记忆信源 $X$ 作 $r$ 元变长编码、平均码长 $\bar L$，则\textbf{任何}唯一可译码都满足
$\bar L\ge\dfrac{\Hx{X}}{\lb r}$，而且\textbf{一定存在}唯一可译码使其满足
$\bar L<\dfrac{\Hx{X}}{\lb r}+1$。二元码（$r=2$，$\lb2=1$）就是最常用的版本：
\begin{align*}
\boxed{\ \Hx{X}\ \le\ \bar L\ <\ \Hx{X}+1\ }
\end{align*}
\end{defbox}

\textbf{逐句翻译}：
\textbf{左边 $\Hx{X}\le\bar L$} 是\textbf{硬下限，谁都突破不了}。直观理由：$N$ 个符号的消息共 $q^N$ 种组合，
要区分它们至少需要约 $\lb q^N$ 比特，而熵正是"平均到底需要多少比特"的精确答案。
\textbf{右边 $\bar L<\Hx{X}+1$} 说\textbf{好码一定存在，且最多只比熵多不到 $1$ 比特}：
多出的那点可理解为"取整的代价"，因为码长必须是整数，而 $-\lb p_i$ 一般是小数，向上取整最多浪费 $1$ 位。
合起来：$\bar L$ 被夹在宽度为 $1$ 的窗口 $[\Hx{X},\Hx{X}+1)$ 里；
想再逼近就得\textbf{分块编码}（把 $N$ 个符号当一个符号编，界限变成
$\Hx{X}\le\frac{\bar L_N}{N}<\Hx{X}+\frac1N$，$N$ 越大越贴近熵）。

\textbf{"为什么能这么近"的构造思路}（一句话版）：直接把码长取成
$l_i=\left\lceil\lb\frac{1}{p_i}\right\rceil$（$\lceil\cdot\rceil$ 是\textbf{向上取整}，
如 $\lceil2.32\rceil=3$）。于是 $\lb\frac1{p_i}\le l_i<\lb\frac1{p_i}+1$，推出
$2^{-l_i}\le p_i$，求和得 $\sum_i2^{-l_i}\le\sum_ip_i=1$，由 Kraft 知这样的前缀码存在；
而 $\bar L=\sum_ip_il_i<\sum_ip_i\left(\lb\frac1{p_i}+1\right)=\Hx{X}+1$。
\textbf{右边的不等号就是这么来的：取整最多浪费 1 位。}

\begin{calc}{$H=2.6087$、$\bar L=2.72$ 的码到底有多好}
拿本节后面 Huffman 那个 $7$ 符号例子的结果核对：
$\Hx{X}=2.6087$ bit/符号，$\bar L=2.72$ bit/符号。
\textbf{先验界限}：$2.6087\le2.72<2.6087+1=3.6087$，两边都满足 $\checkmark$。
\textbf{再算效率}（实际离下限有多近）：
$\eta=\dfrac{\Hx{X}}{\bar L}=\dfrac{2.6087}{2.72}\approx0.9591=95.91\%$。
\textbf{再算冗余度}（浪费了多少）：$1-\eta=1-0.9591=0.0409=4.09\%$，
即每符号多花 $2.72-2.6087=0.1113$ 比特。
\end{calc}

读法：$\bar L$ 只比熵多 $0.111$ 比特，远小于定理允许的 $1$ 比特，说明 Huffman 码已非常接近极限。
考试凡让你"算平均码长并评价"，就要把 $\Hx{X}$、$\bar L$、$\eta$、冗余度\textbf{四个数都写出来}，
再补一句"$\Hx{X}\le\bar L<\Hx{X}+1$ 成立"。

%==============================================================
\subsection{Huffman 编码：把最小平均码长真的做出来}

\subsubsection{一句话原理：每次把最轻的两个箱子捆一起}
还是寄快递：$7$ 个包裹重量（概率）不同，要装成一个大包裹，路径用 $0/1$ 标记。
显然\textbf{越重的包裹该越早被分拣出去（路径越短）}；
而\textbf{最轻的两个最该被"捆在一起"共享同一段路径}——它们最少出现，让它们多走几步代价最小。
做法：\textbf{反复把概率最小的两个合并成新节点（概率相加），直到只剩一个节点（概率 $=1$）；
再从根往回走，左枝记 $0$、右枝记 $1$，路径上的 $0/1$ 序列就是码字。}

\begin{idea}{四条铁律（记住这四条就不会算错）}
(1) 每次取\textbf{当前序列中最小的两个}相加，把和\textbf{插回}序列并\textbf{重新排序}；
(2) 二元码合并次数固定为 $n-1$ 次；
(3) 码符号分配方向可任选（左 $0$ 右 $1$ 或反之），\textbf{所以 Huffman 码不唯一}；
(4) \textbf{必须画码树或写完整合并过程}，只给最终码表老师看不出你是怎么来的。
\end{idea}

\subsubsection{手把手全过程：7 个符号的经典例子}

\begin{calc}{完整 Huffman 编码（考试最可能的计算题，照着抄一遍）}
概率（已降序）$P=(0.20,\,0.19,\,0.18,\,0.17,\,0.15,\,0.10,\,0.01)$，和 $=1$，
对应符号 $s_1,\dots,s_7$。

\textbf{第一步：合并过程}（老师要看的"过程"，一行都不能省）：

\begin{center}
\begin{tabular}{@{}clc@{}}
\toprule
轮次 & 当前概率序列（从大到小） & 本次合并最小的两个\\
\midrule
初始 & $0.20,\ 0.19,\ 0.18,\ 0.17,\ 0.15,\ 0.10,\ 0.01$ & ---\\
第 1 轮 & $0.20,\ 0.19,\ 0.18,\ 0.17,\ 0.15,\ \mathbf{0.11}$ & $0.10+0.01=0.11$\\
第 2 轮 & $0.26,\ 0.20,\ 0.19,\ 0.18,\ \mathbf{0.17}$ & $0.15+0.11=0.26$\\
第 3 轮 & $0.35,\ 0.26,\ 0.20,\ \mathbf{0.19}$ & $0.17+0.18=0.35$\\
第 4 轮 & $0.39,\ 0.35,\ \mathbf{0.26}$ & $0.19+0.20=0.39$\\
第 5 轮 & $0.61,\ \mathbf{0.39}$ & $0.26+0.35=0.61$\\
第 6 轮 & $\mathbf{1.00}$ & $0.39+0.61=1.00$\\
\bottomrule
\end{tabular}
\end{center}

注意每轮都只取"当前最小的两个"，合并后\textbf{重新按大小插回}；
合并次数 $6=n-1$ 次，最后只剩根节点 $1.00$ $\checkmark$。

\textbf{第二步：倒着展开，分配码字}（从根往下，左枝 $0$、右枝 $1$）：
\begin{align*}
1.00&=\underbrace{0.61}_{0}+\underbrace{0.39}_{1}, &
0.61&=\underbrace{0.35}_{00}+\underbrace{0.26}_{01},\\
0.39&=\underbrace{0.19}_{10}+\underbrace{0.20}_{11}, &
0.35&=\underbrace{0.17}_{000}+\underbrace{0.18}_{001},\\
0.26&=\underbrace{0.15}_{010}+\underbrace{0.11}_{011}, &
0.11&=\underbrace{0.10}_{0110}+\underbrace{0.01}_{0111}.
\end{align*}

\textbf{第三步：读出码表}（顺着下标就能读出码字）：

\begin{center}
\begin{tabular}{@{}ccccc@{}}
\toprule
符号 & 概率 $p_i$ & 码字 $w_i$ & 码长 $l_i$ & $p_il_i$\\
\midrule
$s_1$ & $0.20$ & $11$   & $2$ & $0.40$\\
$s_2$ & $0.19$ & $10$   & $2$ & $0.38$\\
$s_3$ & $0.18$ & $001$  & $3$ & $0.54$\\
$s_4$ & $0.17$ & $000$  & $3$ & $0.51$\\
$s_5$ & $0.15$ & $010$  & $3$ & $0.45$\\
$s_6$ & $0.10$ & $0110$ & $4$ & $0.40$\\
$s_7$ & $0.01$ & $0111$ & $4$ & $0.04$\\
\midrule
合计 & $1.00$ & & $\bar L=$ & $2.72$\\
\bottomrule
\end{tabular}
\end{center}

码长向量 $(2,2,3,3,3,4,4)$——\textbf{概率越大码越短}，符合"大配短、小配长"。
\textbf{顺手用 Kraft 验算}（能立刻发现算错）：
\begin{align*}
\sum_i2^{-l_i}=2\times2^{-2}+3\times2^{-3}+2\times2^{-4}
=\frac{2}{4}+\frac{3}{8}+\frac{2}{16}
=\frac{8}{16}+\frac{6}{16}+\frac{2}{16}=1\ \checkmark
\end{align*}
和恰为 $1$，是一棵\textbf{满树}、没有浪费位置（这是 Huffman 码的特征：
只要所有概率都参与合并，树一定是满的）。
\end{calc}

\begin{calc}{四个数一个不能少（考试得分点）}
\textbf{(1) 平均码长}（上表最后一列相加）：
\begin{align*}
\bar L&=0.20\times2+0.19\times2+0.18\times3+0.17\times3+0.15\times3+0.10\times4+0.01\times4\\
&=0.40+0.38+0.54+0.51+0.45+0.40+0.04=2.72\ \text{bit/符号}.
\end{align*}
\textbf{(2) 熵}（用 $\Hx{X}=\sum_ip_i\lb\frac1{p_i}$，先把每个 $\lb\frac1{p_i}$ 单独算出来）：
\begin{align*}
\lb\tfrac{1}{0.20}&=\lb5\approx2.3219, &
\lb\tfrac{1}{0.19}&\approx2.3959, &
\lb\tfrac{1}{0.18}&\approx2.4739,\\
\lb\tfrac{1}{0.17}&\approx2.5564, &
\lb\tfrac{1}{0.15}&\approx2.7370, &
\lb\tfrac{1}{0.10}&=\lb10\approx3.3219,\\
\lb\tfrac{1}{0.01}&=\lb100\approx6.6439 . &&&
\end{align*}
（校验：$2^{2.3219}\approx5$，$2^{6.6439}\approx100$，对得上。）逐项相乘再求和：
\begin{align*}
\Hx{X}&=0.20\times2.3219+0.19\times2.3959+0.18\times2.4739+0.17\times2.5564\\
&\quad+0.15\times2.7370+0.10\times3.3219+0.01\times6.6439\\
&=0.4644+0.4552+0.4453+0.4346+0.4105+0.3322+0.0664\\
&=2.6087\ \text{bit/符号}.
\end{align*}
（中间相加：$0.4644+0.4552=0.9196$，$+0.4453=1.3649$，$+0.4346=1.7995$，
$+0.4105=2.2100$，$+0.3322=2.5422$，$+0.0664=2.6086$，末位舍入差 $0.0001$。）

\textbf{(3) 效率}：$\eta=\dfrac{\Hx{X}}{\bar L}=\dfrac{2.6087}{2.72}\approx0.9591=95.91\%$。\\
\textbf{(4) 冗余度}：$1-\eta=1-0.9591=0.0409=4.09\%$。\\
\textbf{(5) 对照定理}：$2.6087=\Hx{X}\le\bar L=2.72<3.6087=\Hx{X}+1$ $\checkmark$，
且只比下限多 $0.1113$ 比特，比定理允许的 $1$ 比特好得多。
\end{calc}

\begin{idea}{为什么 Huffman 码最优（一句话版）}
教材定理 4.3.3：Huffman 码在所有 $r$ 元即时码中\textbf{平均码长最小}；
再由"所有唯一可译码的最小平均码长与所有即时码的最小平均码长\textbf{相同}"（推论 4.3.2），
知它在所有唯一可译码里也最优。考试一般不要求证明，但要求\textbf{会用}。
\end{idea}

\subsubsection{三元 Huffman 编码：首步要合并几个？}
二元码每次合 $2$ 个，一路下去总会剩 $1$ 个；但 $r$ 元码每次合 $r$ 个，
\textbf{可能最后剩下的不是 $1$ 个}，树就长不完整。解决办法：\textbf{第一次先少合几个}。

\begin{defbox}{首步合并个数（必背公式）}
对 $n$ 个符号作 $r$ 元 Huffman 编码，\textbf{第一轮合并}
\begin{align*}
s=1+\big[(n-1)\bmod(r-1)\big]
\end{align*}
个（若算出 $s=1$，说明不用调整，每轮都合 $r$ 个即可），之后\textbf{每轮都合 $r$ 个}。
人话：\textbf{要让每轮把"待合并个数"减少 $r-1$、最后正好减到 $1$，
就得先用一次小合并把个数调成 $1\bmod(r-1)$ 的样子。}
\end{defbox}

\begin{calc}{三元 Huffman 的小例子（$n=8$，$r=3$）}
概率（降序）$0.22,0.20,0.18,0.15,0.10,0.08,0.05,0.02$（和 $=1$）。
\textbf{先算首步个数}：$n=8$、$r=3$、$r-1=2$，
$s=1+[(8-1)\bmod2]=1+[7\bmod2]=1+1=2$，
所以\textbf{第一轮只合 $2$ 个}，之后每轮合 $3$ 个。走一遍：

\begin{center}
\small
\begin{tabular}{@{}clc@{}}
\toprule
轮次 & 当前概率序列（从大到小） & 本次合并\\
\midrule
初始 & $0.22,\ 0.20,\ 0.18,\ 0.15,\ 0.10,\ 0.08,\ 0.05,\ 0.02$ & ---\\
第 1 轮（合 $2$ 个） & $0.22,\ 0.20,\ 0.18,\ 0.15,\ 0.10,\ 0.08,\ \mathbf{0.07}$ & $0.05+0.02=0.07$\\
第 2 轮（合 $3$ 个） & $0.25,\ 0.22,\ 0.20,\ 0.18,\ \mathbf{0.15}$ & $0.10+0.08+0.07=0.25$\\
第 3 轮（合 $3$ 个） & $0.53,\ 0.25,\ \mathbf{0.22}$ & $0.15+0.18+0.20=0.53$\\
第 4 轮（合 $3$ 个） & $\mathbf{1.00}$ & $0.22+0.25+0.53=1.00$\\
\bottomrule
\end{tabular}
\end{center}

\textbf{分配码字}（三元码用 $0,1,2$）：
\begin{align*}
1.00&=\underbrace{0.53}_{0}+\underbrace{0.25}_{1}+\underbrace{0.22}_{2}, &
0.53&=\underbrace{0.20}_{00}+\underbrace{0.18}_{01}+\underbrace{0.15}_{02},\\
0.25&=\underbrace{0.10}_{10}+\underbrace{0.08}_{11}+\underbrace{0.07}_{12}, &
0.07&=\underbrace{0.05}_{120}+\underbrace{0.02}_{121}.
\end{align*}
码长：$0.22\to1$；$0.20,0.18,0.15,0.10,0.08\to2$；$0.05,0.02\to3$，即 $(1,2,2,2,2,2,3,3)$。
\textbf{平均码长}（单位是三进制位 trit）：
\begin{align*}
\bar L_3&=0.22\times1+2\times(0.20+0.18+0.15+0.10+0.08)+3\times(0.05+0.02)\\
&=0.22+2\times0.71+3\times0.07=0.22+1.42+0.21=1.85 .
\end{align*}
\textbf{熵与界限}：$\Hx{X}=2.7535$ bit，$\lb3\approx1.5850$，
$\dfrac{\Hx{X}}{\lb3}=\dfrac{2.7535}{1.5850}\approx1.7373\le\bar L_3=1.85<2.7373$ $\checkmark$
效率 $\eta=1.7373/1.85\approx93.91\%$。\textbf{注意}：三元码每符号位数看着少（$1.85<2.72$），
但一个 trit 的信息量比一个 bit 多，\textbf{不能直接比大小}，要像上面除以 $\lb r$ 换成比特单位再比。
\end{calc}

\begin{qa}{为什么 Huffman 码不唯一？考试怎么答？}
三个原因，任一条都足以让答案不唯一：
(1) \textbf{左右可互换}——把 $0$ 和 $1$ 对调，码表全变但码长一点没变，$\bar L$ 不变；
(2) \textbf{概率相等时可随便排}——出现两个相等概率时谁左谁右无所谓，码字不同、码长相同；
(3) \textbf{中间序列的排序写法不唯一}——只要守住"取最小两个"的规则，插回位置可以有不同的写法。
标准答法：\textbf{最优码长分配 $(l_1,\dots,l_q)$ 唯一（去掉相等概率的对称性后），
但具体码字不唯一；所有 Huffman 码的平均码长都相同，都等于最优值。}
\end{qa}

%==============================================================
\subsection{其他几种常用变长编码（一句话认识）}

\begin{center}
\small
\begin{tabular}{@{}p{2.7cm}p{6.6cm}p{5.6cm}@{}}
\toprule
方法 & 做法（一句话） & 平均码长性质（一句话）\\
\midrule
Shannon 码 &
概率降序，取 $l_k=\left\lceil\lb\frac{1}{p_k}\right\rceil$，再用累积概率的二进制展开截 $l_k$ 位当码字。 &
即时码，但\textbf{不一定最优}；$\bar L<\Hx{X}+1$，$p$ 极端时浪费大
（如 $p=(0.9999,0.0001)$ 码长为 $1$ 和 $14$）。\\
\addlinespace
Fano 码 &
把符号\textbf{分成两组}使两组概率和尽量接近，一组记 $0$、一组记 $1$，再对每组递归地这样分。 &
不一定最优，只能保证 $\bar L\le\Hx{X}+2$ 量级，比 Huffman 略差。\\
\addlinespace
Shannon--Fano--Elias 码 &
用\textbf{累积概率中点} $\bar F(k)=F(k-1)+\frac{p_k}{2}$ 的二进制展开，
取 $\left\lceil\lb\frac{1}{p_k}\right\rceil+1$ 位当码字。 &
即时码，$\bar L<\Hx{X}+2$；价值在于\textbf{可递推实现}，是算术编码的前身。\\
\addlinespace
算术编码 &
把整条消息映射成 $[0,1)$ 里的一段\textbf{小区间}，用区间内某数的二进制展开当码字；
消息越长区间越窄，需要的位数自动跟着变。 &
\textbf{能逼近极限}：$N\to\infty$ 时冗余度趋于 $0$，极限下实现理想压缩，
这是它强于 Huffman 之处。\\
\addlinespace
游程编码（RLE） &
把二元序列中连续相同的符号换成"长度记数"，如 $00101110010001\cdots\to3,1,1,3,2,1,3,\cdots$
（"0 游程 3 个、1 游程 1 个……"）。 &
一一对应、可逆；本身是变长码，游程长度还需二次编码（Huffman 等），
对\textbf{长串重复}压缩效果极好（传真 MH 码就这么干）。\\
\bottomrule
\end{tabular}
\end{center}

\textbf{怎么记}：Shannon、Fano、SFE 都是"历史方案，都比 Huffman 差一点"；
SFE 的贡献是"递推 + 累积概率"的思想，被算术编码继承；游程编码是另一个赛道（先换表示，再压）。

%==============================================================
\subsection{平稳信源与马尔可夫信源（浅讲，够应付考试）}

\subsubsection{为什么还要熵率}
前面结论都建立在"符号之间独立"（无记忆）上。但真实的中文、英文显然不独立：
看到"信"，下一个字很可能是"息"、是"号"，不可能是"狗"——\textbf{前后文本身就带着信息}。
有记忆时每来一个符号带来的\textbf{新}信息比"独立时"少，因为一部分已被前文猜到了。
所以要一个新指标：\textbf{平均每个符号真正带来了多少新信息}。

\begin{defbox}{熵率与冗余度（4.4.1 节，必背）}
长度 $N$ 的符号块 $X_1X_2\cdots X_N$ 的联合熵为 $\Hx{X_1X_2\cdots X_N}$，
\textbf{平均每符号的熵}为 $H_N(X)=\frac1N\Hx{X_1X_2\cdots X_N}$，\textbf{熵率}定义为
\begin{align*}
H_\infty(X)=\lim_{N\to\infty}H_N(X)=\lim_{N\to\infty}\frac1N\Hx{X_1X_2\cdots X_N},
\end{align*}
而对平稳信源它还等于 $H_\infty(X)=\lim_{N\to\infty}\Hx{X_N\mid X_1\cdots X_{N-1}}$。
\textbf{冗余度}与\textbf{相对冗余度}（$K$ 为源字母表大小）：
\begin{align*}
\text{冗余度}=\lb K-H_\infty(X),\qquad
\text{相对冗余度}=1-\frac{H_\infty(X)}{\lb K}.
\end{align*}
\end{defbox}

\textbf{人话}：$H_\infty$ 就是"\textbf{把前后文都考虑进去以后，每个符号还剩多少真·新信息}"。
它一定不超过 $\Hx{X_1}$（独立时的熵）：\textbf{知道上下文只会减少不确定性}。
定理 4.4.1 把这个直觉说成三件事：$H_N(X)$ 随 $N$ 增大\textbf{单调不增}；
$H_N(X)$ 介于 $0$ 与 $\Hx{X_1}$ 之间；极限存在且有上面两个等价算法。
\textbf{冗余度}的人话：$\lb K$ 是"每符号最多能装多少比特"，$H_\infty$ 是"它实际装了多少"，
差出来的就是\textbf{浪费掉的部分}——也就是"这句话里有多少是语法和用词习惯替你写好的"。

\begin{calc}{英文的冗余度（教材例 4.4.1，背结论）}
英文 $26$ 个字母加空格共 $27$ 个字符，等概率时
$H_0=\lb27=\lb3^3=3\lb3=3\times1.5850=4.7549$ bit/字母。
只考虑字母频率得 $H_1=4.036$；考虑 $2$ 阶、$3$ 阶依赖得 $H_2=3.326$、$H_3=3.31$；
一般认为 $H_\infty\approx1.4$ bit/字母。于是
$\dfrac{H_\infty}{H_0}=\dfrac{1.4}{4.7549}\approx0.294$，$R=1-0.294\approx0.706=70.6\%$。
\textbf{读法}：写英文时约 $71\%$ 的内容是被语言结构"预先规定"的，只有约 $29\%$ 是作者自由选的。
\end{calc}

\subsubsection{马尔可夫信源：状态 + 转移矩阵 + 平稳分布}
\textbf{什么是马尔可夫信源}：信源任一时刻处于某个\textbf{状态} $s_j$，在 $s_j$ 下按
$\p{x_k\mid s_j}$ 发出符号 $x_k$，发完按 $p(s_j\mid s_i)$ \textbf{转移}到新状态。
关键特征是\textbf{下个状态和输出只取决于"现在处于哪个状态"，跟更早的历史无关}
（马尔可夫性/无后效性）。$m$ 阶马尔可夫信源（输出只与前 $m$ 个符号有关）的状态数是 $K^m$。

\textbf{状态转移矩阵} $P=(p_{ij})_{J\times J}$，$p_{ij}=p(s_j\mid s_i)$ 表示"从 $i$ 跳到 $j$"。
按概率的规矩，\textbf{每一行加起来必须等于 $1$}（从某状态出发总要跳到某处去）。
作业 4.12(1) 被点名"转移矩阵写错"，常见两种错：\textbf{行列搞反}
（$p_{ij}$ 是"$i\to j$"，\textbf{行=出发、列=到达}；搞反后"每行和为 $1$"变成"每列和为 $1$"，
后面平稳分布方程整体错位）；\textbf{把"该状态下输出符号的概率"当成转移概率}
（这是两套数）。\textbf{自查办法：填完矩阵立刻检查每行和是否为 $1$。}

\begin{defbox}{平稳分布怎么求（考试标准套路）}
稳态状态概率 $p_1,\dots,p_J$ 必须满足
\begin{align*}
\sum_{i=1}^{J}p_i\,p_{ij}=p_j\ (j=1,\dots,J),\qquad \sum_{j=1}^{J}p_j=1,
\end{align*}
矩阵写法即 $P^{\mathsf T}\boldsymbol p=\boldsymbol p$。
人话：稳态就是"转移一轮之后各状态的分布不再变"——\textbf{进这个状态的概率流正好等于出它的概率流}。
注意这 $J$ 个方程\textbf{线性相关}（只有 $J-1$ 个独立），
所以\textbf{必须配上 $\sum_jp_j=1$} 才解得出唯一答案，这一步千万别漏。
\end{defbox}

\begin{calc}{两状态小例子：手解平稳分布 + 算熵率}
两状态 $s_1,s_2$，转移矩阵（行=出发，列=到达）
$P=\begin{pmatrix}\frac34&\frac14\\[1pt]\frac12&\frac12\end{pmatrix}$。
先自查：第 1 行 $\frac34+\frac14=1$ $\checkmark$，第 2 行 $\frac12+\frac12=1$ $\checkmark$。

\textbf{(1) 解平稳分布}。设稳态概率 $p_1,p_2$，按定义列方程：
\begin{align*}
p_1&=\tfrac34p_1+\tfrac12p_2,\qquad
p_2=\tfrac14p_1+\tfrac12p_2,\qquad
p_1+p_2=1 .
\end{align*}
用第二式（最简单）：$p_2-\frac12p_2=\frac14p_1\Rightarrow\frac12p_2=\frac14p_1\Rightarrow p_2=\frac12p_1$。
代进 $p_1+p_2=1$：$p_1+\frac12p_1=1\Rightarrow\frac32p_1=1\Rightarrow p_1=\frac23,\ p_2=\frac13$。
\textbf{验算}（一定做）：$\frac34p_1+\frac12p_2=\frac34\cdot\frac23+\frac12\cdot\frac13=\frac12+\frac16=\frac46=\frac23$ $\checkmark$。

\textbf{(2) 各状态的条件熵}。设 $s_1$ 下发出 $0$ 的概率 $\frac34$、发出 $1$ 的概率 $\frac14$；
$s_2$ 下发出 $0,1$ 各 $\frac12$。则
\begin{align*}
\Hx{X\mid s_1}&=-\tfrac34\lb\tfrac34-\tfrac14\lb\tfrac14
=0.75\times0.4150+0.25\times2=0.3113+0.5=0.8113,\\
\Hx{X\mid s_2}&=-\tfrac12\lb\tfrac12-\tfrac12\lb\tfrac12=1 .
\end{align*}
（用到 $\lb\frac{4}{3}=\lb1.3333\approx0.4150$、$\lb4=2$。）

\textbf{(3) 熵率}（稳态下按状态概率加权）：
\begin{align*}
H_\infty(X)=\sum_{j=1}^{2}p(s_j)\Hx{X\mid s_j}
=\tfrac23\times0.8113+\tfrac13\times1=0.5409+0.3333\approx0.8742\ \text{bit/符号}.
\end{align*}

\textbf{(4) 冗余度}。源字母表 $\{0,1\}$，$K=2$、$\lb K=1$，
故冗余度 $=1-0.8742=0.1258$ bit，相对冗余度 $=1-0.8742/1=12.58\%$。

\textbf{(5) 与逐符号编码对比}。稳态下 $0,1$ 的边缘概率为
$\p{0}=\tfrac23\cdot\tfrac34+\tfrac13\cdot\tfrac12=\tfrac12+\tfrac16=\tfrac23$，$\p{1}=\tfrac13$。
对 $(\frac23,\frac13)$ 直接做二元 Huffman 只能得到两个长度 $1$ 的码字，$\bar L=1$ bit/符号，
效率 $0.8742/1=87.42\%$。这正好说明：\textbf{熵率 $0.8742$ 靠逐符号编码达不到，
必须把若干符号合起来编（扩展编码）才能逼近}——这就是下面平稳信源编码定理的意思。
\end{calc}

\begin{defbox}{平稳/马尔可夫信源的编码定理（结论记住即可）}
对离散平稳信源作 $r$ 元（教材记 $D$ 元）变长编码，把 $L$ 个源符号当一块编、
每符号平均码长记 $\bar l$，则 $\forall\varepsilon>0$，只要 $L$ 足够大，就存在唯一可译码使
\begin{align*}
\frac{H_\infty(X)}{\lb D}\ \le\ \bar l\ <\ \frac{H_\infty(X)}{\lb D}+\varepsilon .
\end{align*}
\textbf{翻译}：块越长，每符号平均码长就越能\textbf{任意逼近}熵率除以 $\lb D$。
无记忆时 $H_\infty(X)=\Hx{X_1}$，就退回第 5 节的 $\Hx{X}\le\bar L<\Hx{X}+1$。
\end{defbox}

%==============================================================
\begin{exam}{考试怎么考（第4章必考清单）}
\textbf{最高频的五类题}：
(1) \textbf{Huffman 全过程 + 平均码长 + 效率}（大题，$\sim10$ 分）。答题七步：
概率降序列表 $\to$ \textbf{写出每一轮合并过程}（$n-1$ 行）$\to$ \textbf{画码树}并在叶上标符号写码字
$\to$ 列码表（含 $l_i$、$p_il_i$）$\to$ $\bar L=\sum p_il_i$ $\to$ $\Hx{X}=\sum p_i\lb\frac1{p_i}$
$\to$ $\eta=\Hx{X}/\bar L$ 并补一句 $\Hx{X}\le\bar L<\Hx{X}+1$。
\textbf{不画树、只给码表要扣分}（作业 4.4 就栽在这里）。
(2) \textbf{Kraft 不等式判断}：给一组码长问能否构成 $r$ 元前缀码——写出通分后的和、与 $1$ 比较、
下结论；等于 $1$ 补"满树码"。
(3) \textbf{给码表判前缀码/唯一可译码}：$4$ 个小码表逐条判断并各算 $\bar L$；
按"非奇异 $\subset$ 唯一可译 $\subset$ 前缀"的层次说清结论。
(4) \textbf{求熵率与冗余度}：先写平稳分布方程（$P^{\mathsf T}\boldsymbol p=\boldsymbol p$ 加 $\sum p_j=1$）
解出 $p_j$，再算 $H_\infty=\sum_jp(s_j)\Hx{X\mid s_j}$，最后 $R=\lb K-H_\infty$；
\textbf{别忘了检查转移矩阵每行和为 $1$}。
(5) \textbf{概念小问}：$\bar L$ 的定义、三种码的包含关系、熵为什么是下限、
Huffman 码为何不唯一、三元码首步合并个数 $s=1+[(n-1)\bmod(r-1)]$。

\textbf{易丢分细节}：三元（$r$ 元）编码别忘\textbf{先算 $s$}，直接每轮合 $3$ 个会剩两个、树画不出来；
平均码长写清单位（bit/符号 或 trit/符号），效率写成百分数；
熵的 $\lb$ 是以 $2$ 为底（不是 $\ln$、不是 $\lg$）；
算 $\eta$ 时分子是熵、分母是平均码长，\textbf{别写反}（$\eta\le1$ 是天然检查）。
\end{exam}

\begin{qa}{为什么熵是平均码长的下限？}
分两层答。\textbf{(1) 数数}：$N$ 个符号的消息共 $q^N$ 种，要全部区分开至少需 $q^N$ 个不同码字，
即至少 $\lb q^N=N\lb q$ 位、每符号至少 $\lb q$ 位——正是"等概率时熵 $=\lb q$"。
\textbf{(2) 精确版}：设唯一可译码码长为 $l_i$，令 $q_i=\dfrac{r^{-l_i}}{\sum_jr^{-l_j}}$
（把 Kraft 的份额归一化成概率）。由 Kraft 不等式 $\sum_jr^{-l_j}\le1$ 得 $r^{-l_i}\le q_i$，
取对数得 $\lb\frac{1}{q_i}\le l_i\lb r$。于是
\begin{align*}
\sum_ip_i\lb\frac{1}{p_i}-\bar L\lb r
&\le\sum_ip_i\lb\frac{1}{p_i}-\sum_ip_i\lb\frac{1}{q_i}
=\sum_ip_i\lb\frac{q_i}{p_i}=-D(p\Vert q)\le0,
\end{align*}
最后一步用\textbf{相对熵（KL 散度）$D(p\Vert q)\ge0$}。故 $\bar L\ge\Hx{X}/\lb r$；
等号成立当且仅当 $q_i=p_i$，即 $l_i=\lb\frac{1}{p_i}$ 全是整数（概率都是 $2$ 的负整数次幂）。
\textbf{一句话}：熵是"消息本来需要的比特数"，码长是"你实际花掉的比特数"，
Kraft 不等式卡住了你能花的总额度，所以怎么编都不可能比熵更省。
\end{qa}

\begin{qa}{Huffman 为什么一定要"每次合并最小的两个"？合最大的会怎样？}
\textbf{合最小的两个}，等于让最少出现的两个符号共享同一段长路径，
它们占用的加权代价 $p_i\times(\text{多走的位数})$ 最小，对 $\bar L$ 伤害最小。
\textbf{反过来合最大的两个}，就等于给最常见的符号最长码字。拿第 6 节的分布真试一遍：
先合 $0.20$ 与 $0.19$ 得 $0.39$，再合 $0.39$ 与 $0.18$ 得 $0.57$……一直"合最大的"下去，
得到码长 $(6,6,5,4,3,2,1)$——\textbf{恰好倒过来了}，最常见的 $s_1$ 反而要 $6$ 位。
此时
$\bar L'=1.20+1.14+0.90+0.68+0.45+0.20+0.01=4.58$，
效率只剩 $2.6087/4.58\approx57\%$。所以"取最小两个"不是随便规定的，
它是能证明最优（定理 4.3.3）的关键一步。
（顺带一句：\textbf{"大概率配短码"只是必要条件，不等于 Huffman}。比如把码长配成 $(2,2,3,3,4,4,4)$，
虽然大的短小的长、Kraft 也过关（$\frac{2}{4}+\frac{2}{8}+\frac{3}{16}=\frac{15}{16}\le1$），
但 $\bar L=2.87$、$\eta\approx90.9\%$，明显不如 Huffman 的 $2.72$、$95.9\%$。）
\end{qa}

\clearpage
\section{信道编码与汉明码：让收到的 0/1 能自己纠错}

这一讲只干一件事：\textbf{让收到的 0/1 自己把错改回来}。学完你应该能拿一支笔，把 8 位信息
\texttt{10011100} 手工变成 12 位码字 \texttt{111100101100}；并且在路上被翻错一个 bit 之后，
只看收到的 12 位，就反推出错在第几位、把它翻回来。全程只需要你会三样东西：
数\textbf{数}、判\textbf{奇偶}、外加一个叫\textbf{异或}的运算，不涉及概率论。

\subsection{先讲个故事：为什么要"故意多写几位"}

室友在电话里让你帮他记一个 11 位手机号，线不太好，你听到的数字里可能有 1 位是错的。
他念完就挂了。问题来了：你手里这 11 位，\textbf{你完全不知道其中有没有错}。
因为手机号本来就"可以是"任何 11 位数字，你没有任何东西可以拿来核对。也就是说，
\textbf{线路确实可能翻错 1 个 bit，但收方根本发现不了}——这就是"直接发原码"的处境：
错误发生了，却没有留下任何痕迹；更别提确定是哪一位错了。

现实世界的答案是"多带几位没用的数字"，而且你天天在用：身份证号 18 位，前 17 位是信息，
\textbf{最后 1 位是校验码}，它不是新信息，只用来查前面 17 位有没有被抄错；书的 ISBN、
快递单号、银行卡号，最后一位通常也是这个角色。这些"多出来的位"有个正经名字：
\textbf{校验位}。它的取值不自由，而是\textbf{完全由信息位算出来}的，于是收发双方多了一条约定：
\emph{"你算出来的和我算出来的对不上，就说明路上出错了"}。这就把"发现错误"从不可能变成了可能。

本讲处理的就是教材 5.6 节那套办法：要把 8 位信息送过去，但路上每一位都可能以概率 $p$
被翻错（这个"会翻错的线"叫\textbf{二元对称信道}，简称 BSC；设 $p<1/2$，意思是线虽然会错，
但错得不算多——若 $p>1/2$，把收到的每一位都反过来反而更好，这种信道没有讨论价值）。
做法：\textbf{不直接发这 8 位，而是再补 4 位校验，一共发 12 位}。这 4 位不携带新信息，
纯粹是"为了纠错而花掉的开销"。

\begin{idea}{一句话记住：冗余 = 用更长的码换可靠}
\textbf{冗余}（redundancy）就是那些"重复的、能被猜出来"的部分。原码没有冗余，错一位无从判断；
多带 4 位校验，等于给了收方一把尺子，让它能量出"这串码字不对劲"。代价很直白：本来 8 位就够，
现在要占 12 个位置，线要多花 $12/8=1.5$ 倍时间。\textbf{可靠性是用长度换来的。}
\end{idea}

特别要注意：\textbf{信道编码和信源编码的目标正好相反}。

\begin{center}\small
\setlength{\tabcolsep}{5pt}
\begin{tabular}{@{}lll@{}}
\toprule
 & 信源编码 & 信道编码\\
\midrule
要解决的问题 & 数据太大，传得慢 & 数据会传错，不可靠\\
目标 & 让码\textbf{变短} & 让码\textbf{变长}\\
对待冗余 & 尽可能\textbf{挖掉} & 尽可能\textbf{加入}（而且要有规律）\\
衡量指标 & 平均码长、编码效率 & 码率 $R$、纠错能力 $t$\\
典型例子 & Huffman 编码 & 重复码、汉明码\\
\bottomrule
\end{tabular}
\end{center}

一句话：\textbf{信源编码是"把水挤干"，信道编码是"往水里掺沙子"}；掺的沙子必须有规律，
否则收方也不知道哪些是沙子。

\begin{defbox}{四个术语（后面全靠它们说话）}
\textbf{信息位} $D_1,D_2,\dots,D_k$：真正想传的内容，本讲 $k=8$；$D$ 是 data，
$D_1$ 即"第 1 个数据位"。\quad
\textbf{校验位} $C_1,C_2,\dots,C_r$：由信息位算出来的位，本讲 $r=4$；$C$ 是 check，
$C_1$ 即"第 1 个校验位"。\quad
\textbf{码长} $n=k+r$：一个码字共多少位，本讲 $n=8+4=12$。\quad
\textbf{码字}：这 $n$ 位合在一起的整体，比如 \texttt{111100101100} 就是一个 12 位码字。
\end{defbox}

\begin{qa}{先把符号扫盲一遍：$\oplus$ 是什么？}
叫\textbf{异或}。规则只有两条：两个数\textbf{相同}得 0、\textbf{不同}得 1，即
$0\oplus0=0$，$0\oplus1=1$，$1\oplus0=1$，$1\oplus1=0$。要用的就两个性质：
(1) \textbf{连加看奇偶}——一串 bit 异或起来得 1 表示"里面有奇数个 1"，得 0 表示"偶数个 1"；
(2) \textbf{自己消掉自己}——$x\oplus x=0$，$x\oplus 0=x$，所以 $x\oplus y=0$ 等价于 $x=y$。
这两条就是整套汉明码的地基。
\end{qa}

\subsection{三个必须先钉死的概念}

\subsubsection{码率 $R=k/n$}

\begin{defbox}{码率}
一次发 $n$ 位，其中只有 $k$ 位是真信息，则 $R=k/n$。
\end{defbox}

人话：\textbf{每占用 1 个信道位置，平均能送进去多少位真信息}。$R$ 越接近 1 越"划算"，
越小说明为了纠错花的代价越大。三个码比一比：

\begin{center}\small
\setlength{\tabcolsep}{6pt}
\begin{tabular}{@{}lccccl@{}}
\toprule
码 & $k$ & $r$ & $n$ & $R=k/n$ & 说明\\
\midrule
不编码（裸传） & 8 & 0 & 8 & $8/8=1$ & 一点不浪费，错一位就完蛋\\
3 重复码 & 1 & 2 & 3 & $1/3\approx0.333$ & 每位说三遍，最稳也最亏\\
本讲的汉明码 & 8 & 4 & 12 & $8/12=2/3\approx0.667$ & 用一半浪费换"纠 1 位错"\\
教材 (7,4) 汉明码 & 4 & 3 & 7 & $4/7\approx0.571$ & 教材 5.6.2 节的例子\\
\bottomrule
\end{tabular}
\end{center}

注意后两行：同样是"纠 1 位错"，汉明码的 $2/3$ 是重复码 $1/3$ 的两倍。
\textbf{汉明码值得学的理由不是它更可靠，而是它更划算。}

\subsubsection{汉明距离：两个码字"差几位"}

\begin{defbox}{汉明距离}
两个等长的 0/1 串 $x$、$y$ 逐位对齐比一比，\textbf{不同的位置有几个}，
就叫它们的\textbf{汉明距离}，记作 $d(x,y)$。就是"数不同"。
\end{defbox}

例：$x=1011$ 与 $y=1101$。

\begin{center}\small
\setlength{\tabcolsep}{6pt}
\begin{tabular}{@{}lcccc@{}}
\toprule
位置 & 第 1 位 & 第 2 位 & 第 3 位 & 第 4 位\\
\midrule
$x$ & 1 & 0 & 1 & 1\\
$y$ & 1 & 1 & 0 & 1\\
一样吗 & 一样 & \textbf{不一样} & \textbf{不一样} & 一样\\
\bottomrule
\end{tabular}
\end{center}

2 个位置不同，故 $d(1011,1101)=2$。三条必记的事实：$d(x,y)=0$ 当且仅当 $x=y$；
$d(x,y)=d(y,x)$；若发 $x$ 路上被翻了 $t$ 位收到 $y$，则 $d(x,y)=t$——
\textbf{"错了几位"和"汉明距离"是同一件事}。

对一个码（一堆码字组成的集合）还有个关键数字：\textbf{最小距离} $d_{\min}$，即把所有两两之间的
汉明距离都算一遍、取最小的那个。整个码的抗错能力由它决定。

\subsubsection{纠错能力：$d_{\min}\ge 2t+1$ 才能纠 $t$ 位错}

\begin{defbox}{纠错能力判据}
$d_{\min}\ge 2t+1$ 是"能保证纠正 $t$ 位错误"的充要条件。
反过来读最常用：$d_{\min}=3$ 能纠 1 位错，$d_{\min}=5$ 能纠 2 位，$d_{\min}=7$ 能纠 3 位。
\end{defbox}

为什么是 $2t+1$ 而不是 $t+1$？分三步看。

\textbf{第一步（场景）。} 发送码字 $x$，路上被翻 $t$ 位，收到 $y$，则 $d(x,y)=t$——
收方手里这个 $y$ 离真正的答案 $x$ 有 $t$ 步远。

\textbf{第二步（收方干什么）。} 收方不知道 $x$，只能拿 $y$ 去和码本里所有码字比，
挑\textbf{最近的}当答案（叫"最小距离译码"）。要成功，就得保证 $y$ 离 $x$ 比离任何别的码字都近。

\textbf{第三步（算）。} 设另一个码字 $x'$ 满足 $d(x,x')\ge d_{\min}$。三角不等式给出
$d(x,x')\le d(x,y)+d(y,x')$，于是
\[
d(y,x')\ \ge\ d(x,x')-d(x,y)\ \ge\ d_{\min}-t .
\]
要 $d(y,x')>d(x,y)=t$，即 $d_{\min}-t>t$，也就是 $d_{\min}>2t$；整数条件下写作
$d_{\min}\ge 2t+1$。

\textbf{人话版。} 这个 $2t+1$ 是把距离劈成三段用：

\begin{center}\small
\setlength{\tabcolsep}{6pt}
\begin{tabular}{@{}lll@{}}
\toprule
段 & 用掉多少 & 干什么\\
\midrule
给错误留的余地 & $t$ & 收到的东西可能已被踹出去 $t$ 步，这段是"宽容区"\\
给对手留的余地 & $t$ & 别的码字也得允许别人误认到 $t$ 步以内，不能挤着\\
必须多出的 & $1$ & 保证"我最先赢"是\textbf{严格}赢，不是打平\\
\bottomrule
\end{tabular}
\end{center}

若只留 $2t$ 步，就会出现"$x$ 错了 $t$ 位、$x'$ 也错了 $t$ 位，两边正好撞在同一个 $y$ 上"，
收方彻底没主意、只能瞎猜——\textbf{所以必须再拉开 1 步}，这就是"留 2 倍余量再多 1"的来历。
$d_{\min}=3$ 时 $2t+1=3$ 给出 $t=1$：错 1 位能纠，错 2 位不敢保证（可能离另一个码字更近，
纠错反而纠歪）。汉明码"只纠 1 位错"不是它不努力，是它只有 3 的距离。

\subsection{重复码：最笨但最管用的办法}

\textbf{编码规则。} 每个 bit 说三遍：想发 \texttt{0} 就发 \texttt{000}，想发 \texttt{1} 就发
\texttt{111}。码本只有两个码字 $\{$\texttt{000},\texttt{111}$\}$，$n=3$、$k=1$、$R=1/3$。
两者汉明距离是 3（三位全不同），由 $3\ge 2\times1+1$ 知\textbf{能纠 1 位错}。

\textbf{译码规则：多数表决}——收到 3 位，\textbf{哪个数字多就判哪个}：

\begin{center}\small
\setlength{\tabcolsep}{6pt}
\begin{tabular}{@{}cl@{}}
\toprule
收到的 3 位 & 判成\\
\midrule
\texttt{000} & \texttt{0}（没错）\\
\texttt{001}、\texttt{010}、\texttt{100} & \texttt{0}（纠了 1 位错）\\
\texttt{011}、\texttt{101}、\texttt{110} & \texttt{1}（纠了 1 位错）\\
\texttt{111} & \texttt{1}（没错）\\
\bottomrule
\end{tabular}
\end{center}

看第二行：本来发 \texttt{000}，第 3 位被翻成 \texttt{001}，收方数出"0 有两个、1 只有一个"，
判 \texttt{0}，\textbf{纠错成功}。这就是多数表决的全部内容。

\begin{calc}{算重复码的误码率（代入 $p=0.01$，算术一步不跳）}
设每位独立地以概率 $p$ 被翻错。

\textbf{（1）什么时候表决会出错？} 只有一种情况：3 位里\textbf{至少错 2 位}
（错的位数不超过 1，正确的数字一定占多数）。

\textbf{（2）分类计数。}"恰好错 2 位"：从 3 个位置里挑 2 个出错，有 $\binom{3}{2}=3$ 种挑法，
每种概率是"这 2 位错（$p\cdot p$）、剩下 1 位对（$1-p$）"，即 $p^{2}(1-p)$，
合计 $3p^{2}(1-p)$。"3 位全错"只有 1 种情况，概率 $p^{3}$。两类互斥，相加：
\[
P_e=3p^{2}(1-p)+p^{3}.
\]
人话：\textbf{误码率 = "错俩"的概率 + "错仨"的概率}。

\textbf{（3）化简}（考试写出这步会加分）：$P_e=3p^{2}-3p^{3}+p^{3}=3p^{2}-2p^{3}$。

\textbf{（4）代入 $p=0.01$。} 先算中间量：
\[
p^{2}=0.01\times0.01=10^{-4},\qquad p^{3}=p^{2}\cdot p=10^{-4}\times10^{-2}=10^{-6}.
\]
第一项 $3p^{2}(1-p)=3\times10^{-4}\times0.99=2.97\times10^{-4}$；
第二项 $p^{3}=10^{-6}=0.01\times10^{-4}$。相加：
\[
P_e=2.97\times10^{-4}+0.01\times10^{-4}=2.98\times10^{-4}.
\]

\textbf{（5）对比原来。} $P_e/p=2.98\times10^{-4}\big/10^{-2}=2.98\times10^{-2}\approx3\%$，
即错误概率降到原来的约\textbf{三十分之一}。

\textbf{（6）代价。} 码率从 1 掉到 $1/3$：传同样多的信息要多花 3 倍时间。推广到"每 bit 发
$2m+1$ 遍"，则 $R=1/(2m+1)$，误码率趋于 0 但码率也趋于 0——\textbf{这条路走不通}。
\end{calc}

\begin{idea}{重复码的教训}
重复码靠"把信息抄 3 遍"，思路对但极浪费。汉明码的高明之处在于：\textbf{它不抄信息，
而是让 $r$ 个校验位各自"监视一片位置"，再把 $r$ 个监视结果按二进制拼起来，
拼出的数字直接就是"错在第几位"}。于是 $r$ 位校验能照看约 $2^{r}$ 个位置，而不是 1 个。
\end{idea}

\subsection{汉明码的构造（本讲核心，四步模板）}

下面这套流程照着抄，任何"$k$ 位信息、偶校验、构造汉明码"的题都能做。
例题即作业那道：\textbf{信息位 $k=8$，采用偶校验}。先把偶校验讲清楚：

\begin{defbox}{偶校验 / 奇校验}
\textbf{偶校验}：要求"一组 bit 里 1 的个数是\textbf{偶数}"，用异或写就是"这组 bit 异或起来
等于 0"。\textbf{奇校验}：要求"1 的个数是\textbf{奇数}"，即异或起来等于 1。
两者只是目标奇偶不同，其余全部照搬；本讲若不特别说明一律用偶校验。
\end{defbox}

\subsubsection{第一步：定校验位个数 $r$}

$r$ 不能随便挑，要满足
\[
2^{r}\ \ge\ n+1,\qquad\text{其中 } n=k+r .
\]
$n$ 里含 $r$，是自己包含自己的式子，标准做法是\textbf{逐个试}。

\textbf{人话翻译。} $r$ 个校验位最后给出 $r$ 个 0/1 答案，拼起来是 $r$ 位二进制，最多表示
$2^{r}$ 种情况；而收方要区分的状况有 $n+1$ 种：\textbf{"没错"}（1 种）加
\textbf{"第 1 位错"}、$\dots$、\textbf{"第 $n$ 位错"}（$n$ 种）。要能一一报到，就得
$2^{r}\ge n+1$。把 $n=8+r$ 代进去逐个试：

\begin{center}\small
\setlength{\tabcolsep}{6pt}
\begin{tabular}{@{}ccccccl@{}}
\toprule
试的 $r$ & $n=8+r$ & 需区分 $n+1$ 种 & 可用状态 $2^{r}$ & 够不够\\
\midrule
$r=1$ & $n=9$ & $10$ & $2^{1}=2$ & $2<10$，不够\\
$r=2$ & $n=10$ & $11$ & $2^{2}=4$ & $4<11$，不够\\
$r=3$ & $n=11$ & $12$ & $2^{3}=8$ & $8<12$，不够\\
$r=4$ & $n=12$ & $13$ & $2^{4}=16$ & $16\ge13$，\textbf{可行}\\
\bottomrule
\end{tabular}
\end{center}

故 \fbox{$r=4$}，\fbox{$n=k+r=8+4=12$}，这是一个 \textbf{(12,8) 码}，码率 $R=8/12=2/3$。

交叉验证一下教材：$k=4$ 时 $r=2$ 得 $n=6$、$2^{2}=4<7$ 不够；$r=3$ 得 $n=7$、$2^{3}=8\ge8$
恰好够，故 (7,4) 汉明码，与教材 "$n=2^{r}-1$、$k=2^{r}-1-r$" 一致。我们这题 $k=8$ 不是这种
"刚好填满"的值，所以 $2^{4}=16$ 有富余（只用了 13 个状态），13、14、15 号位置不存在，
永远不会被报到，不影响使用。

\subsubsection{第二步：定位置（校验位放在 1、2、4、8 位）}

\begin{defbox}{位置规则（必须背下来）}
把码字的 $n$ 个位置从 1 开始编号。\textbf{编号是 2 的整数次幂的位置放校验位}，即第
$1,2,4,8,16,\dots$ 位（也就是 $2^{0},2^{1},2^{2},2^{3},2^{4},\dots$）放 $C_1,C_2,C_4,C_8,\dots$；
\textbf{其余位置按从小到大的顺序}放信息位 $D_1,D_2,\dots,D_k$。
\end{defbox}

注意校验位的\textbf{下标就是它所在的位置号}：$C_1$ 在第 1 位，$C_4$ 在第 4 位，$C_8$ 在第 8 位。
好处是"名字即地址"，纠错时不容易搞错。把 $n=12$ 排出来：

\begin{center}\small
\setlength{\tabcolsep}{4pt}
\begin{tabular}{@{}l*{12}{c}@{}}
\toprule
位置（十进制） & 1&2&3&4&5&6&7&8&9&10&11&12\\
\midrule
位置（4 位二进制地址） & 0001&0010&0011&0100&0101&0110&0111&1000&1001&1010&1011&1100\\
\midrule
这一位是谁 & $C_1$&$C_2$&$D_1$&$C_4$&$D_2$&$D_3$&$D_4$&$C_8$&$D_5$&$D_6$&$D_7$&$D_8$\\
\bottomrule
\end{tabular}
\end{center}

\textbf{读法。} 第 1、2、4、8 位是校验位 $C_1,C_2,C_4,C_8$（下标即位置号）；剩下的第
3、5、6、7、9、10、11、12 位是信息位，依次是 $D_1$ 到 $D_8$。

\textbf{这里的坑，老师说第 7 位和第 10 位最常写错}：因为第 8 位被 $C_8$ 占掉了，
所以 \textbf{$D_5$ 不是第 8 位，而是第 9 位}，往后全部顺延一位。一旦从这里错位，
$D_5\sim D_8$ 的位置会全错，整个码字就废了。口诀：\textbf{"逢 2 的幂次让给校验位，
数据位见空就填"}。

\subsubsection{第三步：写校验方程（按二进制地址分组）}

\textbf{分组规矩。} 每个位置号写成 4 位二进制（因为 $r=4$）。地址的\textbf{第 1 位}（从右数，
权值 1）是 1 的进 $G_1$、由 $C_1$ 管；\textbf{第 2 位}（权值 2）是 1 的进 $G_2$、由 $C_2$ 管；
\textbf{第 3 位}（权值 4）是 1 的进 $G_4$、由 $C_4$ 管；\textbf{第 4 位}（权值 8）是 1 的进
$G_8$、由 $C_8$ 管。照位置表逐位检查，得四组：

\begin{center}\small
\setlength{\tabcolsep}{4pt}
\begin{tabular}{@{}lll@{}}
\toprule
组 & 包含的位置 & 成员\\
\midrule
$G_1$（地址第 1 位为 1） & $\{1,3,5,7,9,11\}$ & $C_1,\ D_1,\ D_2,\ D_4,\ D_5,\ D_7$\\
$G_2$（地址第 2 位为 1） & $\{2,3,6,7,10,11\}$ & $C_2,\ D_1,\ D_3,\ D_4,\ D_6,\ D_7$\\
$G_4$（地址第 3 位为 1） & $\{4,5,6,7,12\}$ & $C_4,\ D_2,\ D_3,\ D_4,\ D_8$\\
$G_8$（地址第 4 位为 1） & $\{8,9,10,11,12\}$ & $C_8,\ D_5,\ D_6,\ D_7,\ D_8$\\
\bottomrule
\end{tabular}
\end{center}

\textbf{动手核对，别信脑袋。} 拿 $G_2$ 说：$2=0010$、$3=0011$、$6=0110$、$7=0111$、
$10=1010$、$11=1011$，第 2 位都是 1，全部进组；而 $1=0001$、$4=0100$、$5=0101$、
$8=1000$、$9=1001$、$12=1100$ 的第 2 位都是 0，全部不进。一个不多一个不少。

\textbf{偶校验方程}：每组里 1 的个数必须是偶数，即每组异或起来等于 0：
\begin{align*}
G_1:\quad & C_1\oplus D_1\oplus D_2\oplus D_4\oplus D_5\oplus D_7=0,\\
G_2:\quad & C_2\oplus D_1\oplus D_3\oplus D_4\oplus D_6\oplus D_7=0,\\
G_4:\quad & C_4\oplus D_2\oplus D_3\oplus D_4\oplus D_8=0,\\
G_8:\quad & C_8\oplus D_5\oplus D_6\oplus D_7\oplus D_8=0 .
\end{align*}
人话：\textbf{每一行是一个"检查组"，组里的位凑在一起，1 的个数必须凑成偶数}；
校验位的任务就是"补足那个差额"——组里信息位已有奇数个 1 则 $C$ 取 1，
已有偶数个 1 则 $C$ 取 0。由 $x\oplus y=0\Leftrightarrow x=y$，把 $C$ 解出来：
\begin{align*}
C_1&=D_1\oplus D_2\oplus D_4\oplus D_5\oplus D_7,\\
C_2&=D_1\oplus D_3\oplus D_4\oplus D_6\oplus D_7,\\
C_4&=D_2\oplus D_3\oplus D_4\oplus D_8,\\
C_8&=D_5\oplus D_6\oplus D_7\oplus D_8 .
\end{align*}
（在异或里"移项"就是两边同时异或同一个东西，所以把 $C$ 挪到等号右边即可。）

\subsubsection{第四步：代入信息码 \texttt{10011100}，求四个校验位}

题目给的 8 位信息码 \texttt{10011100}（最高位是 $D_1$），逐位拆开：

\begin{center}\small
\setlength{\tabcolsep}{6pt}
\begin{tabular}{@{}lcccccccc@{}}
\toprule
符号 & $D_1$ & $D_2$ & $D_3$ & $D_4$ & $D_5$ & $D_6$ & $D_7$ & $D_8$\\
\midrule
取值 & 1 & 0 & 0 & 1 & 1 & 1 & 0 & 0\\
\bottomrule
\end{tabular}
\end{center}

\begin{calc}{手把手求 $C_1,C_2,C_4,C_8$（异或中间结果全写出）}
\textbf{求 $C_1$}（$G_1$ 组，$1,3,5,7,9,11$ 位）：
\[
C_1=D_1\oplus D_2\oplus D_4\oplus D_5\oplus D_7=1\oplus0\oplus1\oplus1\oplus0 .
\]
从左往右一个个吃掉：$1\oplus0=1$；$1\oplus1=0$；$0\oplus1=1$；$1\oplus0=1$。
故 $C_1=1$。

\textbf{求 $C_2$}（$G_2$ 组，$2,3,6,7,10,11$ 位）：
\[
C_2=D_1\oplus D_3\oplus D_4\oplus D_6\oplus D_7=1\oplus0\oplus1\oplus1\oplus0 .
\]
和上面碰巧是同一串数，同样五步得 $C_2=1$。（组里 $D_1,D_4,D_6$ 三个 1，奇数个，
所以 $C_2$ 取 1 凑偶数。）

\textbf{求 $C_4$}（$G_4$ 组，$4,5,6,7,12$ 位）：
\[
C_4=D_2\oplus D_3\oplus D_4\oplus D_8=0\oplus0\oplus1\oplus0 .
\]
$0\oplus0=0$；$0\oplus1=1$；$1\oplus0=1$。故 $C_4=1$。（组里只有 $D_4$ 一个 1。）

\textbf{求 $C_8$}（$G_8$ 组，$8,9,10,11,12$ 位）：
\[
C_8=D_5\oplus D_6\oplus D_7\oplus D_8=1\oplus1\oplus0\oplus0 .
\]
$1\oplus1=0$；$0\oplus0=0$；$0\oplus0=0$。故 $C_8=0$。
（组里已有 $D_5,D_6$ 两个 1，偶数个，不用再添。）

\textbf{汇总}：\fbox{$C_1=1$}\quad\fbox{$C_2=1$}\quad\fbox{$C_4=1$}\quad\fbox{$C_8=0$}
\end{calc}

\textbf{拼成 12 位码字}：按位置表一位一位填格。

\begin{center}\small
\setlength{\tabcolsep}{4pt}
\begin{tabular}{@{}l*{12}{c}@{}}
\toprule
位置 & 1&2&3&4&5&6&7&8&9&10&11&12\\
\midrule
身份 & $C_1$&$C_2$&$D_1$&$C_4$&$D_2$&$D_3$&$D_4$&$C_8$&$D_5$&$D_6$&$D_7$&$D_8$\\
取值 & 1&1&1&1&0&0&1&0&1&1&0&0\\
\bottomrule
\end{tabular}
\end{center}

从左到右读出来：$\boxed{\ \texttt{1111\,0010\,1100}\ }$，即 \fbox{\texttt{111100101100}}。

\begin{calc}{回代验算：四组里 1 的个数是否都为偶数？（不写这步不给分）}
把校验位\textbf{也算进去}，数每组 1 的个数。

$G_1=\{1,3,5,7,9,11\}$，取值 $C_1,D_1,D_2,D_4,D_5,D_7=1,1,0,1,1,0$，
1 的个数 $=1+1+0+1+1+0=4$，\textbf{偶数} $\checkmark$

$G_2=\{2,3,6,7,10,11\}$，取值 $C_2,D_1,D_3,D_4,D_6,D_7=1,1,0,1,1,0$，
1 的个数 $=4$，\textbf{偶数} $\checkmark$

$G_4=\{4,5,6,7,12\}$，取值 $C_4,D_2,D_3,D_4,D_8=1,0,0,1,0$，
1 的个数 $=1+0+0+1+0=2$，\textbf{偶数} $\checkmark$

$G_8=\{8,9,10,11,12\}$，取值 $C_8,D_5,D_6,D_7,D_8=0,1,1,0,0$，
1 的个数 $=0+1+1+0+0=2$，\textbf{偶数} $\checkmark$

四组分别是 $4/4/2/2$，全为偶数，验算通过。
\end{calc}

\subsubsection{反过来用：收到的东西错在哪一位？}

设发送方发出 \texttt{111100101100}，路上第 10 位被咬了一口，$1$ 翻成 $0$，收方拿到
\[
y=\texttt{1111\,0000\,1100}\quad(\text{第 10 位由 }1\text{ 变成 }0).
\]

\begin{calc}{手算反推出错位置（综合征 / 伴随式）}
收方\textbf{并不知道}原码字，只有 $y$ 和那四条校验方程。做法：拿 $y$ 重新算四个组的奇偶性，
\textbf{算出来是奇记 1、是偶记 0}，四个 0/1 合起来叫\textbf{综合征}（伴随式），记
$(S_1,S_2,S_4,S_8)$。$y$ 从左到右十二位是 $1,1,1,1,0,0,1,0,1,0,0,0$。

\textbf{$S_1$（看 $G_1=\{1,3,5,7,9,11\}$）}：
$1\oplus1\oplus0\oplus1\oplus1\oplus0$，即 $1\oplus1=0$，$0\oplus0=0$，
$0\oplus1=1$，$1\oplus1=0$，$0\oplus0=0$，得 $S_1=0$（组内 4 个 1，偶数，没被破坏）。

\textbf{$S_2$（看 $G_2=\{2,3,6,7,10,11\}$）}：
$1\oplus1\oplus0\oplus1\oplus0\oplus0$，即 $1\oplus1=0$，$0\oplus0=0$，
$0\oplus1=1$，$1\oplus0=1$，$1\oplus0=1$，得 $S_2=1$（组内 3 个 1，奇数，\textbf{被破坏}）。

\textbf{$S_4$（看 $G_4=\{4,5,6,7,12\}$）}：
$1\oplus0\oplus0\oplus1\oplus0$，即 $1\oplus0=1$，$1\oplus0=1$，$1\oplus1=0$，
$0\oplus0=0$，得 $S_4=0$（组内 2 个 1，偶数，没被破坏）。

\textbf{$S_8$（看 $G_8=\{8,9,10,11,12\}$）}：
$0\oplus1\oplus0\oplus0\oplus0$，即 $0\oplus1=1$，$1\oplus0=1$，$1\oplus0=1$，
$1\oplus0=1$，得 $S_8=1$（组内 1 个 1，奇数，\textbf{被破坏}）。

\textbf{拼地址}：按权值 $1,2,4,8$ 把四个结果装成十进制位置号，
\[
S_1\cdot1+S_2\cdot2+S_4\cdot4+S_8\cdot8=0\times1+1\times2+0\times4+1\times8=10,
\]
写成二进制即 $(S_8S_4S_2S_1)=(1010)_{2}=10$。

\textbf{结论：第 10 位错了。} 把 $y_{10}$ 由 0 翻回 1，得 \texttt{111100101100}，
与原码字一模一样。\textbf{纠错完成，收方不需要回问发送方。}
\end{calc}

\begin{defbox}{综合征为什么能"报门牌号"}
因为\textbf{校验矩阵的每一列就是该位置的二进制地址}（见后面 $H$ 那一节）。某一位出错，
只会让"地址里为 1 的那几个组"的奇偶性翻掉，于是综合征正好等于那一列的地址。
所以\textbf{综合征 = 出错位置的二进制地址，直接读出来就行}，不用查表、不用试。
\end{defbox}

两种边界情形要会区分：综合征 $(0,0,0,0)$ 说明\textbf{没有错}（或错了至少 2 位恰好互相抵消，
汉明码管不了）；综合征指向第 1、2、4、8 位（地址 $0001,0010,0100,1000$）说明错的是
\textbf{校验位自己}，信息位没坏，把那个校验位翻回来即可，或干脆忽略。

\subsubsection{对照练习：老师课件里的那道例题}

课件与作业讲评那道题是：8 位信息码 \texttt{11101100}（最高位为 $D_1$），偶校验，构造汉明码。
做法\emph{一个字都不用改}，只换数据。数据位 $D_1$ 到 $D_8$ 依次取 $1,1,1,0,1,1,0,0$，
代入第三步解出的四式：
\begin{align*}
C_1&=D_1\oplus D_2\oplus D_4\oplus D_5\oplus D_7=1\oplus1\oplus0\oplus1\oplus0=1,\\
C_2&=D_1\oplus D_3\oplus D_4\oplus D_6\oplus D_7=1\oplus1\oplus0\oplus1\oplus0=1,\\
C_4&=D_2\oplus D_3\oplus D_4\oplus D_8=1\oplus1\oplus0\oplus0=0,\\
C_8&=D_5\oplus D_6\oplus D_7\oplus D_8=1\oplus1\oplus0\oplus0=0 .
\end{align*}
（前两行：$1\oplus1=0$，$0\oplus0=0$，$0\oplus1=1$，$1\oplus0=1$，得 1；后两行：
$1\oplus1=0$，$0\oplus0=0$，$0\oplus0=0$，得 0。）故 $C_1=1,C_2=1,C_4=0,C_8=0$，填格：

\begin{center}\small
\setlength{\tabcolsep}{4pt}
\begin{tabular}{@{}l*{12}{c}@{}}
\toprule
位置 & 1&2&3&4&5&6&7&8&9&10&11&12\\
\midrule
身份 & $C_1$&$C_2$&$D_1$&$C_4$&$D_2$&$D_3$&$D_4$&$C_8$&$D_5$&$D_6$&$D_7$&$D_8$\\
取值 & 1&1&1&0&1&1&0&0&1&1&0&0\\
\bottomrule
\end{tabular}
\end{center}

得 \fbox{\texttt{111011001100}}，与课件答案一致。回代验算（逐组数 1 的个数）：
\begin{itemize}[leftmargin=1.6em,itemsep=0pt]
  \item $G_1$：$C_1,D_1,D_2,D_4,D_5,D_7$ 取 $1,1,1,0,1,0$，共 4 个 1，偶数；
  \item $G_2$：$C_2,D_1,D_3,D_4,D_6,D_7$ 取 $1,1,1,0,1,0$，共 4 个 1，偶数；
  \item $G_4$：$C_4,D_2,D_3,D_4,D_8$ 取 $0,1,1,0,0$，共 2 个 1，偶数；
  \item $G_8$：$C_8,D_5,D_6,D_7,D_8$ 取 $0,1,1,0,0$，共 2 个 1，偶数。
\end{itemize}
四组全偶，验算通过。

\textbf{讲评提醒}：老师说这道题"汉明码第 7 位和第 10 位编码错误"是常见失分点。对照上表，
第 7 位是 $D_4=0$、第 10 位是 $D_6=1$。根源就是第二步那个坑：第 8 位被 $C_8$ 占了，
$D_5\sim D_8$ 整体后移一位。\textbf{填格之前先把位置表抄一遍、填完再对一遍，就能避开。}

\subsection{线性分组码的一点皮毛：$G$ 和 $H$}

上面的构造看似手工，背后是线性代数。记住两个矩阵的名字和用法就够用了。

\subsubsection{校验矩阵 $H$：把"每组异或为 0"写成一个矩阵}

把第三步的四个组直接写成矩阵：行是组、列是位置，属于该组的位置填 1、不属于填 0，
就得到\textbf{校验矩阵 $H$}（$4\times12$；这样的宽矩阵用表来排更清楚，数的是同一个东西）：

\begin{center}\small
\setlength{\tabcolsep}{4pt}
\begin{tabular}{@{}c*{12}{c}@{}}
\toprule
 & 1&2&3&4&5&6&7&8&9&10&11&12\\
\midrule
$G_1$ 行 & 1&0&1&0&1&0&1&0&1&0&1&0\\
$G_2$ 行 & 0&1&1&0&0&1&1&0&0&1&1&0\\
$G_4$ 行 & 0&0&0&1&1&1&1&0&0&0&0&1\\
$G_8$ 行 & 0&0&0&0&0&0&0&1&1&1&1&1\\
\bottomrule
\end{tabular}
\end{center}

读法：第 1 行对应 $G_1$，它哪几列是 1 就说明 $G_1$ 管哪几位（第 $1,3,5,7,9,11$ 列），
和前面完全一致——一行顶替了一句话。

\textbf{最重要的观察}：竖着看，\textbf{第 $i$ 列就是 $i$ 的 4 位二进制地址}（自下而上是高位
到低位）。如第 10 列是 $(0,1,0,1)^{\mathsf T}$，正好是 $10=8+2$；第 12 列是
$(0,0,1,1)^{\mathsf T}$，正好是 $12=8+4$。\textbf{$H$ 的第 $i$ 列 $=i$ 的二进制表示}——
这就是汉明码的全部秘密。有了 $H$，"12 位向量 $c$ 是合法码字"有了判据：
\[
cH^{\mathsf T}=\mathbf{0}\quad(\text{即四个组异或全为 }0).
\]
人话：\textbf{码字乘校验矩阵的转置得全零，就说明四条方程全满足}。
（所有加法乘法都在二元域上做：$1+1=0$、$1\times1=1$，也就是"异或"和"与"。）

\subsubsection{生成矩阵 $G$：把"信息位 $\to$ 码字"写成一个乘法}

$G$ 是 $8\times12$ 矩阵，\textbf{每一行是一个"基底码字"}：第 $i$ 行是"只有 $D_i=1$、
其余信息位为 0"时算出来的完整 12 位码字。有了这 8 行，任何信息向量都能拼出来：
\[
c=d\,G,\qquad d=(D_1,D_2,\dots,D_8)\in\{0,1\}^{8}.
\]
人话：\textbf{码字 = 信息向量乘生成矩阵}；$d$ 里哪几位是 1，就把 $G$ 对应的几行加起来（异或）。

\begin{center}\small
\setlength{\tabcolsep}{4pt}
\begin{tabular}{@{}c*{12}{c}@{}}
\toprule
 & 1&2&3&4&5&6&7&8&9&10&11&12\\
\midrule
$D_1$ 行 & 1&1&1&0&0&0&0&0&0&0&0&0\\
$D_2$ 行 & 1&0&0&1&1&0&0&0&0&0&0&0\\
$D_3$ 行 & 0&1&0&1&0&1&0&0&0&0&0&0\\
$D_4$ 行 & 1&1&0&1&0&0&1&0&0&0&0&0\\
$D_5$ 行 & 1&0&0&0&0&0&0&1&1&0&0&0\\
$D_6$ 行 & 0&1&0&0&0&0&0&1&0&1&0&0\\
$D_7$ 行 & 1&1&0&0&0&0&0&1&0&0&1&0\\
$D_8$ 行 & 0&0&0&1&0&0&0&1&0&0&0&1\\
\bottomrule
\end{tabular}
\end{center}

\textbf{用法示范（顺便复核前面的答案）。} $d=\texttt{10011100}$ 即 $D_1,D_4,D_5,D_6$ 为 1，
把这四行竖着异或：第 1 位 $1\oplus1\oplus1\oplus0=1$、第 2 位 $1\oplus1\oplus0\oplus1=1$、
第 3 位 $1\oplus0\oplus0\oplus0=1$、第 4 位 $0\oplus1\oplus0\oplus0=1$、
第 7 位 $0\oplus1\oplus0\oplus0=1$、第 9 位 $0\oplus0\oplus1\oplus0=1$、
第 10 位 $0\oplus0\oplus0\oplus1=1$；其余位置这四行都是 0，故异或为 0。
结果 $C_1=1,C_2=1,C_4=1,C_8=0$，码字仍为 \texttt{111100101100}，
\textbf{与第四步手算完全一致}。

\textbf{两个矩阵的关系。} $GH^{\mathsf T}=\mathbf{0}$。人话：\textbf{$G$ 的每一行都是合法码字，
而合法码字过一遍 $H$ 一定是全零}，所以乘出来是零矩阵。这条式子就是"生成矩阵"与
"校验矩阵"的定义关系；考试让你验证时逐行验（每行与 $H$ 各行的内积都为 0）即可。

\subsubsection{用 $H$ 算综合征：$S=yH^{\mathsf T}$}

收到 $y$ 后算 $S=yH^{\mathsf T}=(S_1,S_2,S_4,S_8)$。人话：\textbf{就是把 $y$ 代进四个校验
方程，把四个奇偶结果排成一个 4 位向量}。设错误图样 $e$ 只在第 $i$ 位为 1（其余为 0），
则 $y=c\oplus e$，于是
\[
S=yH^{\mathsf T}=(c\oplus e)H^{\mathsf T}
=\underbrace{cH^{\mathsf T}}_{=\,\mathbf{0}}\oplus\,eH^{\mathsf T}
=H\text{ 的第 }i\text{ 列}.
\]
人话：\textbf{综合征只跟"错在哪"有关，跟"发了什么"完全无关}——所以同一位置的错，
不管发的是什么码字，综合征都一样，这就是它能当"门牌号"的原因。于是 $S=\mathbf{0}$ 表示
无错；$S\ne\mathbf{0}$ 时 $S$ 就是 $H$ 的某一列，列号即出错位置——前面算出的
$S=(0,1,0,1)^{\mathsf T}$ 正是 $H$ 的第 10 列，所以错在第 10 位。

\begin{idea}{教材 5.6.2 与我们这题的关系}
教材用 $r=3$ 的 $(7,4)$ 汉明码，$H$ 是 $3\times7$，每列是 3 位地址，思路完全一样；
而且它是"填满"的（$2^{3}-1=7$，所有非零 3 位地址都用上），故称\textbf{完备码}。
我们这题 $k=8$ 得 $r=4$、$(12,8)$，$2^{4}=16>13$，有 13、14、15 三个地址空着，
所以不是完备码。\textbf{构造方法一字不差，只是矩阵大小不同}。
\end{idea}

\subsection{手把手：已知转移矩阵，分别用两种准则译码（作业 5.11 类型）}

这题与汉明码无关，但同属第 5 章作业且讲评点名过，一并处理。\textbf{题目}：
离散无记忆信道，转移矩阵（行是输入 $x$，列是输出 $y$）为
\[
\bigl(q(y|x)\bigr)=
\begin{pmatrix}
\dfrac{1}{2}&\dfrac{1}{3}&\dfrac{1}{6}\\[2mm]
\dfrac{1}{6}&\dfrac{1}{2}&\dfrac{1}{3}\\[2mm]
\dfrac{1}{3}&\dfrac{1}{6}&\dfrac{1}{2}
\end{pmatrix},
\qquad
p(x)=\bigl(0.5,\ 0.25,\ 0.25\bigr).
\]
求最佳（最大后验）译码规则、极大似然译码规则，以及各自的平均差错率。

\textbf{先把矩阵读顺}：第 1 行第 2 列那个 $\frac13$ 意为"发 $x_1$ 而收到 $y_2$ 的概率是
$\frac13$"，记作 $q(y_2|x_1)=\frac13$。每行加起来都是 1（$\frac12+\frac13+\frac16=1$），
符合"发出一个符号必定收到某个符号"。注意它\textbf{不对称}：$q(y_2|x_1)=\frac13$ 而
$q(y_1|x_2)=\frac16$，所以不能偷懒套对称信道的结论。

\subsubsection{极大似然译码：只比 $q(y|x)$，每列取最大}

\textbf{规则}：收到 $y$ 后把这一列三个数比一比，\textbf{挑最大的}，它对应的 $x$ 就是判决。
注意：\textbf{只看 $q(y|x)$，不许乘 $p(x)$}（讲评原话）。

\begin{center}\small
\setlength{\tabcolsep}{6pt}
\begin{tabular}{@{}lcccc@{}}
\toprule
收到 $y$ & $q(y|x_1)$ & $q(y|x_2)$ & $q(y|x_3)$ & 挑最大 $\Rightarrow$ 判成\\
\midrule
$y_1$ & $\mathbf{1/2}$ & $1/6$ & $1/3$ & $x_1$\\
$y_2$ & $1/3$ & $\mathbf{1/2}$ & $1/6$ & $x_2$\\
$y_3$ & $1/6$ & $1/3$ & $\mathbf{1/2}$ & $x_3$\\
\bottomrule
\end{tabular}
\end{center}

每列最大值恰好都在对角线上（都是 $\frac12$），故 $g_{\mathrm{ML}}(y_1)=x_1$，
$g_{\mathrm{ML}}(y_2)=x_2$，$g_{\mathrm{ML}}(y_3)=x_3$，即"收到什么就判成什么"。

\textbf{平均差错率}用"1 减译对的总概率"最好算：
\[
P_{e,\mathrm{ML}}=1-\sum_{y}p\bigl(g(y)\bigr)\,q\bigl(y\,|\,g(y)\bigr).
\]
人话：\textbf{每个收到值，用真正那个 $x$ 的概率 $p(x)$ 乘"从它发出却收到 $y$"的概率，
加起来是译对的概率}，1 减掉就是错的概率。逐项：
\begin{align*}
y_1:\ & p(x_1)q(y_1|x_1)=\tfrac12\times\tfrac12=\tfrac14=0.25,\\
y_2:\ & p(x_2)q(y_2|x_2)=\tfrac14\times\tfrac12=\tfrac18=0.125,\\
y_3:\ & p(x_3)q(y_3|x_3)=\tfrac14\times\tfrac12=\tfrac18=0.125 .
\end{align*}
译对总概率 $=0.25+0.125+0.125=0.5$，故 $P_{e,\mathrm{ML}}=1-0.5=\boxed{0.5}$。

\textbf{换个算法验一遍}（按"发的是谁"分）：$x_1$ 只在收到 $y_2,y_3$ 时判错，
$\frac12\times(\frac13+\frac16)=\frac14$；$x_2$ 在收到 $y_1,y_3$ 时判错，
$\frac14\times(\frac16+\frac13)=\frac18$；$x_3$ 在收到 $y_1,y_2$ 时判错，
$\frac14\times(\frac13+\frac16)=\frac18$。相加 $\frac14+\frac18+\frac18=\frac12$，一致。

\subsubsection{最大后验（最佳）译码：比 $p(x)\,q(y|x)$}

\textbf{规则}：收到 $y$ 后比的是\textbf{联合量} $p(x)q(y|x)$（等价地比 $p(x|y)$，
因为 $p(x|y)$ 只是 $p(x)q(y|x)$ 除以 $p(y)$，而 $p(y)$ 对同一次收到的 $y$ 是常数，
不影响谁最大），挑最大的那个 $x$。

\begin{center}\small
\setlength{\tabcolsep}{5pt}
\begin{tabular}{@{}lcccc@{}}
\toprule
收到 $y$ & $p(x_1)q(y|x_1)$ & $p(x_2)q(y|x_2)$ & $p(x_3)q(y|x_3)$ & 判成\\
\midrule
$y_1$ & $\tfrac12\cdot\tfrac12=\tfrac14$ & $\tfrac14\cdot\tfrac16=\tfrac1{24}$ & $\tfrac14\cdot\tfrac13=\tfrac1{12}$ & $\mathbf{x_1}$\\
$y_2$ & $\tfrac12\cdot\tfrac13=\mathbf{\tfrac16}$ & $\tfrac14\cdot\tfrac12=\tfrac18$ & $\tfrac14\cdot\tfrac16=\tfrac1{24}$ & $\mathbf{x_1}$\\
$y_3$ & $\tfrac12\cdot\tfrac16=\tfrac1{12}$ & $\tfrac14\cdot\tfrac13=\tfrac1{12}$ & $\tfrac14\cdot\tfrac12=\mathbf{\tfrac18}$ & $\mathbf{x_3}$\\
\bottomrule
\end{tabular}
\end{center}

故 $g_{\mathrm{MAP}}(y_1)=x_1$，$g_{\mathrm{MAP}}(y_2)=x_1$，$g_{\mathrm{MAP}}(y_3)=x_3$。

\textbf{与极大似然只差一处：$y_2$ 的判决变了。} 原来 $q(y_2|x_2)=\frac12$ 是列最大，
但乘上 $p(x_2)=0.25$ 后只剩 $\frac18$；而 $x_1$ 虽然 $q(y_2|x_1)=\frac13$ 不是列最大，
可它的 $p(x_1)=0.5$ 最大，乘出来 $\frac16>\frac18$ 于是反超。
\textbf{这就是"先验概率把结论掰过来"的活例子}，也是最大后验优于极大似然的原因：
它知道 $x_1$ 本来就更容易出现。

\textbf{平均差错率}同样用 1 减译对概率：
\begin{align*}
y_1:\ & p(x_1)q(y_1|x_1)=\tfrac14,\\
y_2:\ & p(x_1)q(y_2|x_1)=\tfrac12\times\tfrac13=\tfrac16,\\
y_3:\ & p(x_3)q(y_3|x_3)=\tfrac14\times\tfrac12=\tfrac18 .
\end{align*}
译对总概率 $=\frac14+\frac16+\frac18=\frac{6}{24}+\frac{4}{24}+\frac{3}{24}=\frac{13}{24}$，故
\[
P_{e,\mathrm{MAP}}=1-\frac{13}{24}=\boxed{\frac{11}{24}\approx0.4583}.
\]
验算（按发送方分类）：$x_1$ 只有收到 $y_3$ 时判错，$\frac12\times\frac16=\frac{2}{24}$；
$x_2$ 三个 $y$ 全判成别人，$\frac14\times1=\frac{6}{24}$；$x_3$ 在收到 $y_1,y_2$ 时判错，
$\frac14\times(\frac13+\frac16)=\frac18=\frac{3}{24}$。相加 $\frac{2+6+3}{24}=\frac{11}{24}$，一致。

\begin{exam}{讲评点名的那句话：极大似然别乘 $p(x)$}
这道题最常见的错法是算极大似然时悄悄乘上了 $p(x)$，于是把 $P_{e,\mathrm{ML}}$ 算成 $11/24$
——\textbf{那是最大后验的答案}。区别只有一句话：\textbf{极大似然只看 $q(y|x)$、每列挑最大，
默认各 $x$ 等可能出现}；\textbf{最大后验看 $p(x)q(y|x)$，把"谁更容易出现"也算进去}。
所以\textbf{当且仅当 $p(x)$ 等概时两者给出同一套规则}。本题 $p(x)=(0.5,0.25,0.25)$ 不等概，
答案自然不同：$0.5$ 与 $11/24$，后者更小，符合"最佳译码不差于极大似然"。
\end{exam}

\subsection{考试怎么考}

\begin{exam}{构造汉明码的完整答题步骤（照抄这四个小标题，分步给分）}
\textbf{1. 定 $r$ 与 $n$。} 写不等式 $2^{r}\ge n+1$ 且 $n=k+r$，然后\textbf{把 $r=1,2,3,\dots$
逐个代进去试}，每步写清"此时 $n=?$，$2^{r}=?$ 与 $n+1=?$ 谁大"，最后写"故取 $r=\square$，
$n=k+r=\square$"。\emph{只写答案不写试算过程会扣分。}

\textbf{2. 列位置表。} 列表列出 $1\sim n$ 位，标清哪几位是 $C_1,C_2,C_4,C_8$（2 的幂次位）、
哪几位是 $D_1\sim D_k$（其余按顺序）。\emph{这张表是后面所有计算的依据，必须写全。}

\textbf{3. 写校验方程。} 把每个位置号写成 $r$ 位二进制地址，按"地址第 $j$ 位为 1"分成 $r$ 组，
写 $r$ 个方程（偶校验 $=0$、奇校验 $=1$），再解成 $C_{2^{j}}=\bigoplus(\text{组内信息位})$。

\textbf{4. 代入求校验位、拼码字、回代验算。} 逐个算 $C$（写出异或中间步骤），按位置表拼成
$n$ 位，\textbf{最后一定回代 $r$ 个方程、确认每组 1 的个数满足奇偶要求}，并写一句"验算通过"。

\medskip
\textbf{若还要纠错}：多写一步——把收到的 $y$ 代进同样的 $r$ 个方程得综合征
$(S_{1},S_{2},S_{4},S_{8})$，按权值 $1,2,4,8$ 组装成十进制位置号，指出第几位错并翻转。
\end{exam}

\begin{exam}{四个常见扣分点}
\textbf{1. 校验位位置放错。} 必须放在第 $1,2,4,8,\dots$ 位。若把校验位全堆在末尾
（前 8 位数据、后 4 位校验），那就不是教材这套汉明码了，位置表、分组、地址全部对不上。
\emph{对策：动手算之前先把位置表写出来给阅卷老师看。}

\textbf{2. 偶校验 / 奇校验混淆。} 偶校验要求"每组 1 的个数为偶数"、方程写 $=0$；
奇校验要求"奇数"、方程写 $=1$。\emph{混了就有一半校验位全反，白算。}

\textbf{3. 回代不验算。} 只写最终码字、不代回方程检查，一个位置填错就整题错。
\emph{对策：算完逐组数 1 的个数，把"4/4/2/2，全偶，通过"这行写上去}——既是检查也是给分点。

\textbf{4. 把 $2^{r}\ge n+1$ 写错。} 常见三种：写成 $2^{r}\ge n$（少 1，会漏掉"没错"这个状态）；
忘了 $n$ 里含 $r$（应 $n=k+r$）；不等号方向写反。\emph{对策：写不等式时旁边注明
"$n+1$ 是因为要区分'没错'+$n$ 个'第几位错'"。}
\end{exam}

\subsection{你可能想问}

\begin{qa}{为什么校验位一定要放在 1、2、4、8 位？放在别处不行吗？}
不是"不行"，而是"放这里最省事、最漂亮"。\textbf{第一，只有放在 2 的幂次位，位置的二进制
地址才刚好各缺一位。} 第 1 位地址 $0001$、第 2 位 $0010$、第 4 位 $0100$、第 8 位 $1000$，
它们各自\textbf{只属于一个组}，于是 $C_1$ 只负责 $G_1$ 一组的奇偶性，责任划分干净。
若把校验位放到第 3 位（地址 $0011$），它同时属于 $G_1$ 和 $G_2$，一错就同时弄脏两组结果，
综合征不再"干净地等于出错位置"，纠错就乱了。

\textbf{第二，唯一性。} 要求是"$n$ 个位置各有一个不同的非零地址，好当门牌号"；
$1,2,4,8,\dots$ 正好是能自由取值、互不冲突的位置，而数据位塞在剩下的位置就自动获得
"被多个组同时监控"的待遇，一错有多组报警，地址才排得出来。所以这不是规定，
而是\textbf{让"综合征 = 出错位置"这个性质成立的必要条件}。
\end{qa}

\begin{qa}{偶校验和奇校验有什么区别？选哪个？}
\textbf{区别只一处}：偶校验要求每组 1 的个数是\textbf{偶数}（方程右端 $=0$），奇校验要求是
\textbf{奇数}（方程右端 $=1$）；于是求校验位时"补差"方向反过来——组内信息位已有奇数个 1 时，
偶校验要 $C=1$，奇校验要 $C=0$。

\textbf{纠错能力完全一样}，都能纠 1 位错、$d_{\min}$ 都是 3。区别只在"全 0 码字"：
偶校验下 \texttt{00}$\cdots$\texttt{0} 是合法码字（0 个 1 是偶数），奇校验下不是，
所以奇校验能顺带发现"全 0"这种明显异常（比如线路断了收不到信号）。工程上两种都常见，
\textbf{考试时题目说哪个就用哪个}，把方程右端写成对应的 0 或 1 即可，其余步骤一字不改。
\end{qa}

\begin{qa}{汉明码既然厉害，为什么只纠 1 位错？错 2 位怎么办？}
因为它 $d_{\min}=3$，由 $d_{\min}\ge2t+1$ 得 $t=1$，\textbf{这是距离决定的硬上限，
跟算法聪不聪明无关}。错 2 位时两个错误会弄脏两组校验结果，综合征指向"第 3 个位置"——
既不是真正的出错位置也不表示无错，收方会\textbf{自信地纠错，结果越纠越错}。
这是汉明码的固有风险：\textbf{它能纠 1 位错，但错 2 位时可能纠歪。}

要纠更多错就得拉开距离：把汉明码再加一位总校验（扩展汉明码）可把 $d_{\min}$ 从 3 提到 4，
从而"纠 1 位 + 检出 2 位"；要真纠 2 位需 $d_{\min}\ge5$，得用 BCH 码、RS 码、卷积码这类
更强的码，代价是码率进一步下降——又回到那句老话：\textbf{可靠性是用长度换的}。
\end{qa}

\begin{idea}{这一讲的三句话总结}
\textbf{1. 为什么要编码}：线会翻错 bit，裸传原码收方无从察觉；加冗余位（校验位）让错误留下
痕迹。冗余 = 用更长的码换可靠，方向与信源编码正好相反。

\textbf{2. 汉明码怎么造}：$2^{r}\ge n+1$ 定 $r$ $\to$ 校验位放 $1,2,4,8$ 位 $\to$ 按二进制
地址分 $r$ 组写偶校验方程 $\to$ 代入求校验位、回代验算
（$k=8$ 得 $r=4,n=12$；\texttt{10011100} $\to$ \texttt{111100101100}）。

\textbf{3. 怎么纠错}：把收到的 $y$ 代进同样 $r$ 个方程得综合征，\textbf{综合征就是出错位置的
二进制地址}，翻回来即可（$(0,1,0,1)\to$ 第 10 位）。
\end{idea}

\clearpage
\section{最大熵原理与最小鉴别信息原理}

\begin{idea}{开场：这一讲解决什么问题}
前几讲的 $p(a_k)$ 都是\textbf{题目给好}的。现实恰恰相反——\textbf{分布不知道}，手上只有几条零碎情报（"平均长度 $5.02$ mm""方差是 $3$"）。于是问题来了：\emph{和这几条情报相符的分布有无穷多个，我该挑哪一个？}

第 8 章给两条互相呼应的挑法：\textbf{最大熵原理}（手上没有"老经验"时用，Jaynes 1957）和\textbf{最小鉴别信息原理}（手上有一个先验分布时用）。本讲讲清它们\emph{是什么、为什么合理、怎么算}，以及配套的\textbf{鉴别信息}怎么手算——这片是简答题高发区。
\end{idea}

\subsection{先把符号翻译成人话}

\begin{center}\small
\begin{tabular}{lll}
\toprule
符号 & 念法 & 人话 \\
\midrule
$X$，$a_k$ & "随机变量 X""第 $k$ 个取值" & 还没定下来的结果，和它的第 $k$ 种可能（掷骰子的点数） \\
$p(a_k)$ & "p of a-k" & $X$ 取到 $a_k$ 的概率，$\sum_k p(a_k)=1$ \\
$\Hx{X}$ & "H of X" & 平均不确定性 = 平均信息量，单位 bit \\
$\E f(X)$ & "$f(X)$ 的期望" & 把 $f(a_k)$ 按概率加权平均：$\sum_k f(a_k)\p{a_k}$ \\
约束 & constraint & "已知的那几条情报"，都写成某个期望等于常数 \\
\bottomrule
\end{tabular}
\end{center}

熵就是"每个结果的不确定度 $-\lb\p{a_k}$ 按概率加权平均"：$\Hx{X}=-\sum_k\p{a_k}\lb\p{a_k}$。概率越小越意外、信息量越大，所以熵 = 平均意外程度。这条你早就知道，本讲不重复。

\subsection{最大熵原理：情报不全时选最"散"的那个分布}

\textbf{硬币的例子。} 掷一枚硬币，不知道是否均匀，只测到"平均正面率是 $0.6$"。这条情报把唯一自由度用光了：$p(\text{正})=0.6$、$p(\text{反})=0.4$，熵
\begin{align*}
\Hx{X}=-0.6\lb 0.6-0.4\lb 0.4=0.6\times0.7370+0.4\times1.3219=0.4422+0.5288=0.9710\bit ,
\end{align*}
这就是最大熵解。真正需要"挑"的是\emph{约束比自由度少}的情形。

\textbf{骰子的例子（教材例 8.1.2）。} 掷骰子，只知道点数平均值是 $4.5$，即 $\sum_{k=1}^{6}k\,f_k=4.5$。可行解有无穷多个（均匀分布均值只有 $3.5$，不合格；把概率往大点数挪一点就合格，挪法无数）。按最大熵原理挑出来的解是

\begin{center}\small
\begin{tabular}{lcccccc}
\toprule
点数 $a_k$ & 1 & 2 & 3 & 4 & 5 & 6 \\
\midrule
$p(a_k)$ & 0.0543 & 0.0788 & 0.1142 & 0.1654 & 0.2398 & 0.3476 \\
\bottomrule
\end{tabular}
\end{center}

当场验两件事：\textbf{(1) 满足情报}——按点数加权求和
$1\times0.0543+2\times0.0788+3\times0.1142+4\times0.1654+5\times0.2398+6\times0.3476
=0.0543+0.1576+0.3426+0.6616+1.1990+2.0856=4.5007\approx4.5$，合格。\textbf{(2) 熵是多少}——$\Hx{X}=-\sum_k\p{a_k}\lb\p{a_k}\approx2.327\bit$（教材此处印 $1.61358$，是按\emph{自然对数}算的：$2.327\times\ln 2=1.6133$，同一个数、不同单位；答题按题目要求的底走）。

\begin{defbox}{最大熵原理（Jaynes, 1957）}
在只掌握部分信息（只有若干约束）的情况下要对分布作推断时，应取\textbf{满足全部约束、但熵值最大}的那个分布。这是唯一"不偏不倚"的选择，别的选择都等于偷偷添加了额外的约束或假设。数学上是一个约束优化问题：
\begin{align*}
\max_{p}\ &\Hx{X}=-\sum_{k=1}^{K}\p{a_k}\lb\p{a_k}\\
\text{s.t.}\ &\sum_{k=1}^{K}\p{a_k}f_m(a_k)=C_m,\quad m=1,2,\dots,M;
\qquad \sum_{k=1}^{K}\p{a_k}=1,\ \ \p{a_k}\ge 0 .
\end{align*}
\end{defbox}

\textbf{人话}：把情报翻译成函数 $f_m$ 和常数 $C_m$——给了均值就取 $f_1(a_k)=a_k$，再给二阶矩就取 $f_2(a_k)=a_k^2$；然后在所有满足这几条式子的分布里，挑\emph{平均不确定性最大}的那个。

\textbf{为什么合理}：凡是没有被情报约束到的方向，我们没有理由偏向任何一边；熵最大的分布恰是对这些方向"最不表态"的分布。教材 8.1.2 用\textbf{熵集中定理}把它做实：熵稍小于最大值的可行解，个数少到可忽略——教材对上面的掷骰子例算出，$95\%$ 的可行解满足 $1.609\le H\le H_{\max}$、$99.99\%$ 满足 $1.602\le H\le H_{\max}$（与 $1.61358$ 同单位）。所以最大熵解是\emph{最保险}的解。

不过要记住：只有"仅归一约束"时最大熵解才均匀。掷骰子例里均匀分布均值只有 $3.5$，\emph{根本不满足} $4.5$ 这条情报，压根不在可行解集合里——最大熵是"在可行解里挑熵最大的"，不是"在所有分布里挑熵最大的"。

\subsection{怎么算出来：拉格朗日乘子法}

直接令 $\partial\Hx{X}/\partial\p{a_k}=0$ 是解不出来的：$p$ 之间被"和为 1""加权和等于 $C_m$"绑着，动一个别的都得跟着动。标准套路是\textbf{把约束"罚"进目标函数}：引入 $M+1$ 个待定乘子 $\beta,\lambda_1,\dots,\lambda_M$，构造

\begin{align*}
F=\Hx{X}-\beta\Bigl(\sum_{k}\p{a_k}-1\Bigr)-\sum_{m=1}^{M}\lambda_m\Bigl(\sum_{k}\p{a_k}f_m(a_k)-C_m\Bigr),
\end{align*}

对每个 $\p{a_k}$ 求偏导并令其为零（底取 $e$）：

\begin{align*}
\frac{\partial F}{\partial \p{a_k}}=-\ln\p{a_k}-1-\beta-\sum_{m=1}^{M}\lambda_m f_m(a_k)=0
\;\Longrightarrow\;
\log \p{a_k}=-1-\beta-\sum_{m=1}^{M}\lambda_m f_m(a_k),
\end{align*}

于是离散情形的最大熵解（教材式 8.1.7）是

\begin{align*}
\boxed{\;\p{a_k}=e^{\,-1-\beta-\sum_{m=1}^{M}\lambda_m f_m(a_k)}\;}
\end{align*}

\textbf{人话}：每个概率都长成"$e$ 的某个线性组合次方"。$f_m$ 决定\emph{形状}——给均值则指数上是 $x$ 的一次式，概率按指数型升或降；给方差则带平方，概率成钟形。$\beta,\lambda_1,\dots,\lambda_M$ 这 $M+1$ 个常数不神秘：\textbf{由那 $M+1$ 个约束方程（$M$ 条情报 $+$ 归一）唯一确定}。不必会解这个方程组（一般要数值解），但要会\emph{写出形式}和\emph{数未知数个数}。

连续情形（求和换积分）形式一模一样：$p(x)=\exp\bigl(-\lambda_0-\sum_m\lambda_m f_m(x)\bigr)$，其中 $-\lambda_0$ 顶替了 $-1-\beta$。教材提醒：此时严格说叫"微分熵最大原理"，但习惯上仍称最大熵原理。底数取 $e$ 或 $2$ 只差一个固定因子，会被归一化吸进 $\beta$，分布形状不变。

\subsection{三个必须背的结论（每个配一句为什么）}

\begin{center}\small
\begin{tabular}{>{\raggedright\arraybackslash}p{0.22\textwidth}>{\raggedright\arraybackslash}p{0.21\textwidth}>{\raggedright\arraybackslash}p{0.20\textwidth}>{\raggedright\arraybackslash}p{0.25\textwidth}}
\toprule
已知情报 & 最大熵分布 & 熵值 & 为什么 \\
\midrule
只有"概率和为 1" & 均匀 $p(a_k)=1/K$ & $\Hx{X}=\lb K$ &
熵在概率全相等时最大，边界（某个 $p=1$）处为 $0$ \\
只给均值 $\E X$ & 连续：指数；离散非负整数：几何 & 见下 &
$f_1(x)=x$，指数上是线性的 \\
给均值与方差 & 正态 $N(\mu,\sigma^2)$ & $\Hc{X}=\tfrac12\lb(2\pi e\sigma^2)$ &
$f_1=x,f_2=x^2$，指数上是二次式，配方即高斯 \\
\bottomrule
\end{tabular}
\end{center}

\textbf{(a) 只有"总和为 1"这一约束 $\Rightarrow$ 均匀分布，$\Hx{X}=\lb K$。}
为什么：$-\sum p\lb p$ 是凹函数，在"概率全相等"这个中心取最大；一旦某个 $p(a_k)=1$（其余全 $0$），熵就跌到 $0$。所以没有别的约束时，"最散"就是"最平"。连续版：$X$ 落在 $[a,b]$ 上时 $p(x)=1/(b-a)$，$\Hc{X}=\lb(b-a)$。

\textbf{(b) 只给定均值 $\Rightarrow$ 指数分布（离散时是几何分布）。}
为什么：$M=1$、$f_1(x)=x$，代进解式就是 $p(x)\propto e^{-\lambda_1x}$——正是"指数型衰减"。具体地：取非负整数、$\E X=m$ 时是几何分布 $p_k=(1-\theta)\theta^k$，其中 $\theta=\frac{m}{1+m}$（因为几何分布均值是 $\frac{\theta}{1-\theta}$，解出来就是它）；$x\ge0$、$\E X=\mu$ 时 $p(x)=e^{-x/\mu}/\mu$，$\Hc{X}=\log_2(e\mu)$。

\textbf{(c) 给定均值与方差 $\Rightarrow$ 正态分布，$\Hc{X}=\frac12\lb(2\pi e\sigma^2)$。}
为什么：$M=2$、$f_1=x$、$f_2=x^2$，指数上是 $-\lambda_1x-\lambda_2x^2$，二次式配方必然得到 $e^{-(x-\mu)^2/(2\sigma^2)}$ 的形状。注意熵里只有 $\sigma^2$、没有 $\mu$：\emph{搬动钟形的位置不改变不确定性，把钟形拉胖才改变}——方差翻倍，熵增加 $\tfrac12$ bit。

\begin{idea}{和第 6 讲（连续熵最大熵定理）的关系}
第 6 讲在\textbf{微分熵}框架下已证过这三条（支撑集固定 $\to$ 均匀；只给均值 $\to$ 指数；给均值与方差 $\to$ 高斯，$H_c=\frac12\lb(2\pi e\sigma^2)$）。第 8 章做的是把它们\emph{收进同一个框架}：一个统一解 $p(a_k)=e^{-1-\beta-\sum_m\lambda_mf_m(a_k)}$，代入不同 $f_m$ 就分别退化；再补上"为什么该选熵最大的"（不偏不倚 $+$ 熵集中）。所以\textbf{第 6 讲那三条是本讲结论的特例}——问"最大熵原理与前面最大熵定理的关系"，答这一句即可。
\end{idea}

\subsection{鉴别信息：两个分布"差多远"}

先看它从哪冒出来。真实分布是 $P_1$，你以为分布是 $P_2$，于是照 $P_2$ 设计码字（码长 $-\lb p_2(a_k)$）。真实数据服从 $P_1$，实际平均码长为 $\bar L=-\sum_k p_1(a_k)\lb p_2(a_k)$；若一开始就知道 $P_1$，最优也只要熵 $\Hx{X}_{P_1}=-\sum_{k}p_1\lb p_1$。差多少？

\begin{align*}
\bar L-\Hx{X}_{P_1}=-\sum_{k}p_1\lb p_2+\sum_{k}p_1\lb p_1=\sum_{k}p_1(a_k)\lb\frac{p_1(a_k)}{p_2(a_k)} .
\end{align*}

\begin{defbox}{鉴别信息（Kullback--Leibler 距离）}
设 $X$ 在假设 $H_1,H_2$ 下的分布为 $P_1,P_2$，定义
\begin{align*}
D(P_1\|P_2)=\sum_{k=1}^{K}p_1(a_k)\lb\frac{p_1(a_k)}{p_2(a_k)}
\qquad(\text{连续时求和换积分、}p\text{换密度}).
\end{align*}
别名：\textbf{相对熵、交叉熵、方向散度、Kullback--Leibler 距离}；教材用 Kullback 最早的名字"鉴别信息"。

\textbf{人话}：\textbf{用分布 $P_2$ 去近似真实分布 $P_1$，平均每个符号要多付多少 bit}。它衡量的是\emph{两个分布之间的差异}，不是某个分布自己的不确定性。注意方向：$P_1$ 是"真相"，$P_2$ 是"你以为的（先验）"。
\end{defbox}

\begin{calc}{手算 $D(P_1\|P_2)$：$P_1=(0.5,0.5)$，$P_2=(0.9,0.1)$}
\textbf{第 1 项}：$\frac{p_1}{p_2}=\frac{0.5}{0.9}=\frac59\approx0.5556$，$\lb 0.5556=\frac{\ln0.5556}{\ln2}=\frac{-0.5878}{0.6931}=-0.8480$，乘权重 $0.5\times(-0.8480)=-0.4240$。

\textbf{第 2 项}：$\frac{p_1}{p_2}=\frac{0.5}{0.1}=5$，$\lb 5=\frac{\ln5}{\ln2}=\frac{1.6094}{0.6931}=2.3219$，乘权重 $0.5\times2.3219=1.1610$。

\textbf{相加}：$D(P_1\|P_2)=-0.4240+1.1610=0.7370\approx0.737\bit$。

\textbf{验证这个数}：按 $P_2$ 设计的理想码，在真实分布 $P_1$ 下平均码长 $=0.5\lb\frac{1}{0.9}+0.5\lb\frac{1}{0.1}=0.0760+1.6610=1.737\bit$，而按真相设计只要 $\Hx{X}_{P_1}=1\bit$，多付 $0.737\bit$——和 $D$ 完全一致，不是巧合。

\textbf{反过来算一遍}（理解"不对称"的最好办法）：
$D(P_2\|P_1)=0.9\lb\frac{0.9}{0.5}+0.1\lb\frac{0.1}{0.5}=0.9\times0.8480+0.1\times(-2.3219)=0.7632-0.2322=0.5310\approx0.531\bit$。
同一对分布，一个方向 $0.737$、另一个方向 $0.531$，\textbf{不相等}。
\end{calc}

\begin{defbox}{鉴别信息的性质（简答题直接抄）}
\begin{enumerate}[leftmargin=1.6em,itemsep=2pt,topsep=2pt]
\item \textbf{非负性}：$D\ge0$，且 $=0$ 当且仅当 $P_1=P_2$。为什么：由 $\ln x\le x-1$ 可得——\emph{用错分布永远不可能少付码长}。
\item \textbf{不对称}：一般 $D(P_1\|P_2)\ne D(P_2\|P_1)$，故又叫\textbf{方向散度}（directed divergence）。例子就是上面 $0.737$ 对 $0.531$。
\item \textbf{不满足三角不等式}：$D(P_1\|P_3)\le D(P_1\|P_2)+D(P_2\|P_3)$ 一般\emph{不}成立，所以它\textbf{不是真正的距离函数}，只是"差异程度"的度量。
\item \textbf{凸性}：$D(P\|R)$ 对 $P$ 凸，$D(R\|P)$ 对 $P$ 也凸（教材性质 8.2.2）。这保证"求最小值"不会掉进局部最优——正是最小鉴别信息原理能用的根本原因。
\item \textbf{散度是对称的}：$J(P_1,P_2)=D(P_1\|P_2)+D(P_2\|P_1)\ge0$ 且 $J(P_1,P_2)=J(P_2,P_1)$（上面的例子 $J=0.737+0.531=1.268\bit$），但它同样不满足三角不等式。
\item 两条只记结论：随机变量\textbf{相互独立}时，联合鉴别信息等于各自的鉴别信息之和；但鉴别信息\textbf{没有分步可加性}，一般 $D(P_1\|P_2)\ne D(P_1\|P_3)+D(P_3\|P_2)$——这点与 Shannon 熵不同。
\end{enumerate}
\end{defbox}

\begin{qa}{它叫"距离"，为什么又不能说它是距离？}
真距离要满足非负、对称、三角不等式三条，$D$ 只满足第一条。标准答法：\emph{不是距离；它非负但不满足对称性与三角不等式，故称方向散度；两个方向相加得到的 $J$ 才对称。}
\end{qa}

\subsection{最小鉴别信息原理：手里有"老经验"时怎么办}

现实问题常是：你对 $X$ 有个\textbf{老经验}（先验分布 $p(x)$），现在又做了新测量、得到若干约束 $\int qf_m\,dx=C_m$，要在满足约束的 $q$ 里挑一个当估计。

\begin{defbox}{最小鉴别信息原理}
设真实分布 $q(x)$ 未知，已知先验 $p(x)$、约束 $\int qf_m\,dx=C_m\ (m=1..M)$ 与 $\int q\,dx=1$。则应在\textbf{满足全部约束}的 $q$ 中，取使鉴别信息 $D(q\|p)=\int q\lb\frac{q}{p}\,dx$ \textbf{最小}者为估计。
\end{defbox}

\textbf{人话}：约束说"必须满足这几条"，先验说"大概长什么样"；两者都得照顾，就选\emph{离先验最近}（改动最少、多付信息最少）的那个可行解。这是关于函数 $q$ 的变分问题，用拉格朗日乘子解得

\begin{align*}
q(x)=p(x)\exp\Bigl(-\lambda_0-\sum_{m=1}^{M}\lambda_m f_m(x)\Bigr),
\end{align*}

$M+1$ 个常数由 $M+1$ 个约束确定。人话：\textbf{先验乘一个"约束修正因子"}——没约束时 $\lambda_m=0$，$q$ 就等于 $p$（老经验直接用，一点不改）。

\begin{idea}{两条原理的关系（必考一句话）}
先验取\textbf{等概率}（离散 $p(a_k)=1/K$）时，上式变成 $q(a_k)\propto\exp\bigl(-\sum_m\lambda_m f_m(a_k)\bigr)$——正是上一节的最大熵解（常数 $1/K$ 被并进归一化因子）。

所以：\textbf{先验取均匀分布时，最小鉴别信息原理退化为最大熵原理}；反过来说，\textbf{最小鉴别信息原理是最大熵原理的推广}（教材 8.3 原话）。直觉也顺：均匀先验 $=$ "我没有任何偏爱"，此时"离先验最近"就等于"最不表态"，也就是"熵最大"。

教材习题 8.4 的结论再加固一句：若 $p_1$ 是约束下的最大熵分布、$p_2$ 是任一满足同样约束的分布，则 $\Hx{X}_{p_1}-\Hx{X}_{p_2}=D(p_2\|p_1)\ge0$——最大熵分布的熵确实不小于任何可行解，它就是"离所有可行解最近"的那个。
\end{idea}

\subsection{考试怎么考}

\begin{exam}{第 10 讲（第 8 章）的题型与答法}
\textbf{一、简答题（背下来就能拿分）}
\begin{enumerate}[leftmargin=1.6em,itemsep=1pt,topsep=2pt]
\item \emph{什么是最大熵原理}：只掌握部分信息（若干约束）时分布不唯一，应在满足全部约束的分布中选熵最大的那个（Jaynes 1957）；这是唯一不偏不倚的选择，别的选择等于额外加了假设。
\item \emph{为什么合理}：熵集中定理——熵略小于最大值的可行解个数少到可忽略（掷骰子例 $95\%$ 的可行解落在 $H_{\max}-0.005$ 以内），故最大熵解最保险。
\item \emph{三个典型结论}：只归一 $\to$ 均匀 $H=\lb K$；只给均值 $\to$ 指数（离散为几何）；给均值与方差 $\to$ 正态 $H_c=\frac12\lb(2\pi e\sigma^2)$。每条都要补一句"为什么"。
\item \emph{鉴别信息的定义与性质}：定义式 $+$ 非负、不对称（方向散度）、不满足三角不等式、凸性 $+$ 散度 $J$ 对称。
\item \emph{最小鉴别信息与最大熵的关系}：先验取均匀时前者退化为后者，故前者是推广。
\end{enumerate}

\textbf{二、计算题}
\begin{enumerate}[leftmargin=1.6em,itemsep=1pt,topsep=2pt]
\item \emph{给约束写最大熵分布}：用拉格朗日乘子写出 $p(a_k)=e^{-1-\beta-\sum_m\lambda_m f_m(a_k)}$，并说明"$\beta$ 与 $M$ 个 $\lambda_m$ 由 $M+1$ 个约束确定"；若再要求具体分布，就代 $f_m$ 判形状（一次式 $\to$ 指数/几何，二次式配方 $\to$ 高斯）。
\item \emph{算 $D(P_1\|P_2)$}：两步——逐项算 $\lb\frac{p_1}{p_2}$，再按 $p_1$ 加权求和，写清 bit；顺手算反方向，顺势答"不对称"。
\item \emph{用 $D$ 解释编码失配}：$\bar L-\Hx{X}_{P_1}=D(P_1\|P_2)$，这是第 11 讲的入口。
\end{enumerate}

\textbf{三、最容易丢分}：(1) 底数不声明；(2) $D$ 的方向写反——\textbf{权重永远挂在"真相"那一侧}；(3) 正态熵漏掉 $e$（是 $2\pi e\sigma^2$）；(4) 误以为"最大熵解一定是均匀分布"。
\end{exam}

\section{组合信息与算法信息}

\begin{idea}{开场：干脆不用概率}
前面一切（含上一讲）都建立在"有个概率分布"上。可现实里分布估不准；更极端的，\textbf{信号根本不是随机的}——一段程序代码、一个常数的小数展开，硬套随机模型，模型本身就错。

1965 年 Kolmogorov 在《关于信息度量定义的三种方法》中指出，除统计信息外还有两条路：\textbf{组合信息}（从计数、枚举出发——\emph{数一数有多少种可能的序列}）与\textbf{算法信息}（从计算机程序出发——\emph{产生这个串的最短程序有多长}）。对应的编码是 \textbf{Fitingof 通用编码}和 \textbf{Lempel--Ziv 通用编码}。
\end{idea}

\subsection{自适应统计编码：分布不知道时怎么编码}

第 4 讲的 Huffman 编码按"统计特性"造码书，故称统计编码。麻烦在于编码时\textbf{并不知道}真实分布：你拿猜的 $p(a_k)$ 造码书，而真实分布是 $g(a_k)$，这叫\textbf{统计失配}。失配要付多少钱？设码长取 $-\lb p(a_k)$，则

\begin{align*}
\bar L-\bar L^{*}=\sum_{k}g(a_k)\lb\frac{g(a_k)}{p(a_k)}=D\bigl(g\|p\bigr)\ \ge0 .
\end{align*}

\textbf{一句话：失配的代价正好是鉴别信息}——把分布估错，平均每个符号就多付 $D$ bit。这条把上一讲和第 4 章焊在一起，也是本章的出发点。

\begin{center}\small
\begin{tabular}{>{\raggedright\arraybackslash}p{0.15\textwidth}>{\raggedright\arraybackslash}p{0.77\textwidth}}
\toprule
类型 & 一句话 \\
\midrule
静态编码 & 用一套\textbf{固定}的统计特性造码书；猜得准就最优，猜错就惨（多付 $D$）。\\
半自适应 & 消息\textbf{扫两遍}：第一遍统计造码书，第二遍照码书编码。缺点是要把码书/统计特性一起传过去，还有两遍扫描的延时。\\
自适应 & 只扫\textbf{一遍}，边统计边编码：每步都用"以往消息"的统计特性编下一段；收发两端握有相同历史，可各自把码书同步更新（核心是"能逐步更新的数据结构"）。\\
\bottomrule
\end{tabular}
\end{center}

结论一句话：\textbf{自适应编码在任何情况下只略差于最佳静态编码，而静态编码在失配时会远劣于自适应编码。}

\subsection{组合信息：不数概率，只数"有多少种"}

\begin{idea}{核心思想}
\textbf{如果不管概率，只数一数有多少种可能的序列，也能定义信息量。}

道理很朴素：编码的本质是"给可能的情况\emph{编号}"，有 $N$ 种可能就得用 $\lb N$ 位二进制。于是\textbf{组合信息}定义为
\begin{align*}
H_c(X)=\lb N\qquad(\text{$N$ 为可能取值的个数}),
\end{align*}
下标 $c$ 是 Combination（组合），用以区别统计熵 $H(X)$；一般取底 2，单位仍是 bit。配套定义照搬 Shannon 的形式（但含义变了）：
\begin{align*}
H_c(X_1\cdots X_K)&=\textstyle\sum_{i}\lb N_i,\qquad
H_c(Y\mid X=a)=\lb N(Y\mid a),\\
I_c(X;Y)&=H_c(Y)-H_c(Y\mid X),
\end{align*}
$N(Y\mid a)$ 是"$X=a$ 时 $Y$ 可能取的值的个数"。\textbf{人话}：不给概率只数格子；知道 $X$ 后 $Y$ 的可能格子少了多少，就是 $X$ 的信息。
\end{idea}

教材表 9.1（$X,Y$ 各取 $1,2,3,4$，表里"$\surd$"表示该组合存在）：$\lb 4=2$，故 $I_c(X{=}1;Y)=2-\lb 4=0$，$I_c(X{=}2;Y)=2-\lb 3\approx0.415$，$I_c(X{=}3;Y)=2-\lb 2=1$，$I_c(X{=}4;Y)=2-\lb 1=2$。$X=1$ 时 $Y$ 仍有 $4$ 种可能，所以 $X=1$ 什么信息也没给；$X=4$ 时 $Y$ 只剩 $1$ 种，$X$ 把 $Y$ 全确定了。\textbf{提醒}：组合信息与统计信息性质\textbf{不一样}，\emph{不可}把 Shannon 的公式直接套过来（教材明确强调）。

\subsubsection{组合数与典型序列}

设字母表大小 $K$、序列长 $N$、第 $k$ 个字母出现 $N_k$ 次（频率 $f_k=N_k/N$），则"出现次数恰为这组 $N_k$"的序列个数是组合数

\begin{align*}
W=\frac{N!}{\prod_{k=1}^{K}N_k!},
\qquad
\lb W\ \approx\ N\Bigl(-\sum_{k=1}^{K}f_k\lb f_k\Bigr)=N\,H(f).
\end{align*}

"$\approx$"从哪来：由 Stirling 公式 $n!\approx\sqrt{2\pi n}\,(n/e)^n$ 代入，$\sqrt{\cdot}$ 那些因子变成 $N$ 的低阶项，主项就是 $N H(f)$。二元情形 $W=C(N,k)$，取 $k=\alpha N$ 得

\begin{align*}
\lb C(N,\alpha N)\ \approx\ N\,H_2(\alpha),\qquad
H_2(\alpha)=-\alpha\lb\alpha-(1-\alpha)\lb(1-\alpha).
\end{align*}

\textbf{人话}：长 $N$ 的串里，"$1$ 的占比约为 $\alpha$"的序列大约有 $2^{NH_2(\alpha)}$ 个——这一堆叫\textbf{典型序列}；在这堆里指定一个，需要的位数就是 $\lb C(N,\alpha N)\approx N H_2(\alpha)$。

\begin{calc}{数值体检：$N=6$、$\alpha=1/3$}
\textbf{(1) 精确算}：$C(6,2)=\frac{6!}{2!\,4!}=\frac{6\times5}{2}=15$ 个。而 $2^3=8<15<16=2^4$，故 $\lb 15\in(3,4)$，实际 $\lb 15=\frac{\ln15}{\ln2}=\frac{2.7081}{0.6931}=3.9069\approx3.91$ bit。

\textbf{(2) 公式算}：$H_2(1/3)=-\frac13\lb\frac13-\frac23\lb\frac23=\frac13\times1.5850+\frac23\times0.5850=0.5283+0.3900=0.9183$，于是 $N H_2(\alpha)=6\times0.9183=5.51$ bit。

\textbf{(3) 对照}：$3.91$ 对 $5.51$，差 $1.6$ bit，不小。这不是算错：$\lb C(N,\alpha N)\approx NH_2(\alpha)$ 是 $N\to\infty$ 的\textbf{渐近}式（严格说是 $\frac{2^{NH_2(\alpha)}}{N+1}\le C\le 2^{NH_2(\alpha)}$ 这种量级估计），$N=6$ 太小，轮不到它准。所以小 $N$ 必须用精确的 $C(6,2)=15$ 与 $\lb 15\approx3.91$。
\end{calc}

\subsubsection{课件那个例子：把 \texttt{001001} 压成 \texttt{0011}}

老师课件上原样的例子，考试很爱考这种"数一数、编个号"的题。本质只有一句：\textbf{编码 $=$ 在候选表里查位置，再把位置号写成二进制}。

\begin{calc}{手把手：序列 \texttt{001001} 的组合编码}
\textbf{第 1 步（数候选）}：$6$ 位二进制串总共 $2^6=64$ 个，但已知"恰好有 $2$ 个 $1$"，候选只有"$6$ 位里挑 $2$ 个位置放 $1$"的那些：
\begin{align*}
C(6,2)=\frac{6!}{2!\,4!}=\frac{6\times5}{2}=15\ \text{个}.
\end{align*}

\textbf{第 2 步（排序、编号，从 $0$ 开始）}：把这 $15$ 个串按字典序摆好，得下表。

\begin{center}\small
\begin{tabular}{clclcl}
\toprule
索引 & 序列 & 索引 & 序列 & 索引 & 序列 \\
\midrule
0 & \texttt{000011} & 5 & \texttt{001100} & 10 & \texttt{100001} \\
1 & \texttt{000101} & 6 & \texttt{010001} & 11 & \texttt{100010} \\
2 & \texttt{000110} & 7 & \texttt{010010} & 12 & \texttt{100100} \\
3 & \texttt{\textbf{001001}} & 8 & \texttt{010100} & 13 & \texttt{101000} \\
4 & \texttt{001010} & 9 & \texttt{011000} & 14 & \texttt{110000} \\
\bottomrule
\end{tabular}
\end{center}

查表：\texttt{001001} 在\textbf{索引 3}（从 $0$ 数：$0,1,2,3$，第四个）。

\textbf{第 3 步（编号要多宽）}：$15$ 个候选，编号范围 $0\sim14$。因为 $2^3=8<15\le16=2^4$，至少 $4$ 位，即 $\lceil\lb 15\rceil=\lceil3.9069\rceil=4$ 位（$3$ 位只能表示 $0\sim7$，装不下 $15$ 个）。

\textbf{第 4 步（号码转二进制）}：$3=0\times8+0\times4+1\times2+1\times1$，即四位二进制 \texttt{0011}。

\textbf{结论}：\texttt{001001}（$6$ 位）$\to$ \texttt{0011}（$4$ 位）。\textbf{解码}：收到 \texttt{0011} 读成 $3$，查同一张表的索引 $3$，还原 \texttt{001001}；表格两端约定即可，不必传输。

\textbf{一句提醒}：这里能省 \(2\) 位，是因为"恰好 $2$ 个 $1$"双方都知道。若没约定，还得额外传 $k$（$k$ 有 $7$ 种取法，约 $\lb 7\approx2.8$ 位），省下的就被吃掉了——这正是教材定理 9.3.4 中 $\log N+NH(k/N)$ 的来历：\textbf{先传 $k$，再传序号}。
\end{calc}

\subsection{Fitingof 通用编码：两步编号}

把上面的思路用到任意信源，就是 Fitingof 编码。一句话：\textbf{统计各符号出现次数（得频率分布 $F$），再用组合数编出"同类里的第几个"}。三步：

\textbf{(1)} 把序列切成若干长度 $N$ 的组（$N$ 足够大）。\textbf{(2)} 数出每组里各字母的出现次数 $N_k$、得频率分布 $F$；不同 $F$ 共 $W(K,N)=C_{N+K-1}^{\,N}$ 种，编"是哪种 $F$"需 $\lceil\lb W(K,N)\rceil$ bit。\textbf{(3)} 同一 $F$ 下序列仍有 $W(F)=\frac{N!}{\prod_kN_k!}$ 种，编"是其中第几个"需 $\lceil\lb W(F)\rceil$ bit。

于是 $L=\lceil\lb W(F)\rceil+\lceil\lb W(K,N)\rceil$。它\textbf{压根没用统计特性}（只数次数），但教材定理 9.2.1 证明：任给 $\varepsilon>0$，$N$ 足够大时平均每符号码长可小于 $H(P)+\varepsilon$，故称\textbf{通用}编码。缺点是那个 $+\varepsilon$ 收得\emph{很慢}，有限码长下反而不如统计编码，工程上因此用得不多。

\subsection{算法信息：最短程序有多长}

\begin{defbox}{Kolmogorov 算法熵（复杂度）}
\textbf{定义}：一个数据串的复杂性，就是\textbf{能产生这个数据串的最短二元程序的长度}。形式化地，设 $p$ 是"给定 $y$ 时让计算机 $C$ 输出 $x$"的二元程序、$l(p)$ 为其长度，则 $x$ 在条件 $y$ 下的条件算法熵为
\begin{align*}
K(x\mid y)=\min_{p:\,C(p,y)=x} l(p),
\end{align*}
无条件版记 $K(x)$；$y$ 提供的关于 $x$ 的信息记 $I(x;y)=K(x)-K(x\mid y)$。

\textbf{人话}：越有规律 $\to$ 程序越短 $\to$ 复杂度越低；越没规律 $\to$ 只能"照抄" $\to$ 程序和数据一样长 $\to$ 复杂度越高。
\end{defbox}

严格说 $K$ 依赖具体计算机，但换一台通用 Turing 机只差一个与 $x$ 无关的常数（相当于多一段"仿真程序"前缀），故"计算机"理解为任意一台通用机即可。

\begin{calc}{课件两个例子，把"复杂度"讲透}
\textbf{串 A $=0101010101010101$}（$16$ 位，"$01$"重复 $8$ 次）。本来要 $16$ 位才能照抄，但你可以写个短程序 \texttt{print("01"*8)}：只有十几个字符，$8$ bit ASCII 也就一百来 bit，而且\textbf{与串长无关}（串变成 $1000$ 位，写成 \texttt{"01"*1000} 即可）。所以 $K(A)$ 很小，\textbf{算法复杂度低}。

\textbf{串 B $=1001101001011010$}（$16$ 位，看着没什么规律）。想说清它只能照抄 \texttt{print("1001101001011010")}，程序里那段数据\textbf{和原串一样长}，$K(B)\approx$ 串长 $=16$ 位，\textbf{算法复杂度高}。

\textbf{结论（反直觉但必须记住）}：\textbf{随机串反而复杂度最高}——"有规律"正是可压缩的来源，完全随机 $=$ 完全没规律 $=$ 只能照抄，一个字都省不掉。
\end{calc}

这个结论有定理撑腰。教材定理 9.3.3：复杂度小于 $k$ 的串总数少于 $2^k$——因为长度小于 $k$ 的二元程序只有 $1+2+\cdots+2^{k-1}=2^k-1<2^k$ 个，一个程序至多产生一个输出，故\textbf{低复杂度串极少}。进一步（定理 9.3.5、9.3.6）：$N$ 位串有 $N\,H(X)\le\E[K(x\mid N)]\le N\,H(X)+\lb N+C$，且绝大多数串 $\lim_{N\to\infty}K/N=1$——一个字也压不动，称\textbf{不可压缩序列}（即"真正的随机序列"），故算法熵又叫\textbf{算法随机度}。

\begin{qa}{不可能有个通用程序，算出任意串的最短程序}
1970 年 Chaitin 证明：\textbf{数学上不存在通用办法对任意给定的串算出其算法熵}。

直观论证：若真有程序能算 $K(x)$，就可写个很短的程序"找出第一个 $K(x)\ge10^{10}$ 的二元串并打印它"。这个描述本身很短，打出来的串复杂度应该很小；可它偏偏是"复杂度 $\ge10^{10}$"的第一个，矛盾。这与 Perry 悖论（"用少于 $30$ 个英文单词不能表示的最小正整数"，而这句话本身只有 $20$ 个词）、Gödel 不完备定理同源。故 Kolmogorov 复杂度的价值不在求那段最短程序，而在给出研究随机性与可压缩性的理论框架。
\end{qa}

\textbf{Lempel--Ziv 压缩，一句话}：\textbf{找重复片段，用"偏移 + 长度"表示}。既然"理想最短程序"算不出来，就限制计算模型用"找重复"来逼近：LZ77 用滑动窗口，把当前内容在"已编码区"里找最长匹配，输出三元组（偏移量, 匹配长度, 下一个字符）。课件那个串 \texttt{ababbbabaabababaababbababa} 切成 \texttt{a.b.abb.baba.ababa.baababb.ababa}，编码为
\begin{align*}
(0,0,\texttt{a})\ (0,0,\texttt{b})\ (2,2,\texttt{b})\ (4,3,\texttt{a})\ (9,4,\texttt{a})\ (7,6,\texttt{b})\ (10,5,\varnothing).
\end{align*}
看 $(2,2,\texttt{b})$：往回退 $2$ 个字符（\texttt{ab}）、抄 $2$ 个（得 \texttt{ab}）、后面跟一个 \texttt{b}，拼出 \texttt{abb}——这就是"偏移+长度"的秘密（LZ78 改用"字典索引, 下一个字符"，本质一样）。它为何算逼近算法复杂度？"重复片段"就是"规律"，找到规律就能用短码代长串，正是 $K(x)$ 想度量的事；平稳遍历信源下压缩比当 $N\to\infty$ 趋于熵率 $H(X)$，且\emph{不需知统计特性}，故是通用编码，也是 gzip/zip 的核心。

\subsection{近似熵与样本熵（简答级）}

前面讲的都是"一个串"。工程里更常见的是\textbf{时间序列}（心率、血压、振动信号），要问的是：\emph{这串数据有多"有规律"？} 思路：切成很多长 $m$ 的小片段，看"两个片段像不像"；\textbf{片段越容易找到相似的伙伴，序列越有规律，熵越小}。

\begin{defbox}{近似熵 ApEn 与样本熵 SampEn}
\textbf{近似熵}（Pincus, 1991）：对每个 $i$，统计与第 $i$ 个 $m$ 长片段距离小于阈值 $r$ 的片段占比 $B^m(i)$，取对数平均
\begin{align*}
\phi^m(r)=\frac{1}{N-m+1}\sum_{i=1}^{N-m+1}\ln B^m(i),
\qquad
\mathrm{ApEn}(m,r,N)=\phi^{m}(r)-\phi^{m+1}(r).
\end{align*}
\textbf{样本熵}（Richman 等, 2000）：只在\emph{其余} $N-m$ 个片段中统计匹配（\textbf{不把自己算进去}），得平均 $B^m(r)$、$B^{m+1}(r)$，则
\begin{align*}
\mathrm{SampEn}(m,r,N)=-\ln\frac{B^{m+1}(r)}{B^{m}(r)}.
\end{align*}
\end{defbox}

\textbf{人话}：$\phi^m$ 是"用 $m$ 维模式看这串数据时，平均多容易找到相似片段"；维数由 $m$ 加到 $m+1$ 后相似度掉了多少，就是"新模式出现的概率"。掉得越少 $\Rightarrow$ 越有规律 $\Rightarrow$ 熵越小；周期信号的近似熵很小，随机信号较大。且它与幅值、相位无关，只与频率成分的复杂性有关。

\textbf{两者区别（必考一句）}：近似熵把数据和\textbf{自身}比较（自己当然和自己匹配，正好避免 $\ln0$），但引进了偏差、结果不恒定；样本熵\textbf{不计自身匹配}——"熵是新信息产生率"，拿数据和它自己比毫无意义——故精度更高、一致性更好。两者都依赖 $m,r$（$m$ 常取 $2\sim10$，$r$ 按数据定），没有公认最佳标准。

\subsection{考试怎么考}

\begin{exam}{第 11 讲（第 9 章）的题型与答法}
\textbf{一、组合信息计数题（最容易满分）}——四步写全就是满分：
\begin{enumerate}[leftmargin=1.6em,itemsep=1pt,topsep=2pt]
\item 数候选：确定有几个"$1$"（或几个某字母），算 $C(n,k)=\frac{n!}{k!(n-k)!}$（如 $C(6,2)=\frac{6\times5}{2}=15$）。
\item 排号：按字典序排出索引，\textbf{从 $0$ 开始}，找出目标串的索引（\texttt{001001}$\to3$）。
\item 定码长：最小的 $l$ 使 $2^l\ge$ 候选数，即 $\lceil\lb N\rceil$（$2^3=8<15\le2^4\Rightarrow4$ 位）。
\item 转二进制：索引写成 $l$ 位二进制（$3\to$\texttt{0011}），并说明解码端查同一张表即可还原。
\end{enumerate}
最后补一句"若 $k$ 未事先约定，还要额外传 $k$，约需 $\lb(n+1)$ 位"是加分点。

\textbf{二、算法复杂度判断题}——给两个串问哪个更随机/复杂度更高，答法固定：\emph{能写成短程序（如 \texttt{print("01"*8)}）的复杂度低；只能原样照抄的复杂度高，$K(x)\approx$ 串长}。再补反直觉结论：\textbf{随机串复杂度最高，因为没规律就无法压缩}，并引定理 9.3.3（复杂度小于 $k$ 的串不足 $2^k$ 个，低复杂度串极少，绝大多数串 $K/N\to1$ 即不可压缩）。

\textbf{三、简答题}
\begin{enumerate}[leftmargin=1.6em,itemsep=1pt,topsep=2pt]
\item \emph{Kolmogorov 复杂度定义}：产生该串所需最短二元程序的长度，条件版 $K(x\mid y)=\min l(p)$。
\item \emph{不可计算性}：Chaitin（1970）证明不存在通用算法算出任意串的算法熵；用"最小复杂度 $\ge10^{10}$ 的串"的自指矛盾说明，与 Gödel 不完备定理、Perry 悖论同源。
\item \emph{组合信息}：$H_c(X)=\lb N$（$N$ 为可能取值个数），与分布无关；典型序列 $\lb C(N,\alpha N)\approx NH_2(\alpha)$；\emph{不可}照搬 Shannon 公式。
\item \emph{自适应编码三类}：静态（固定特性）/ 半自适应（扫两遍、要传码书）/ 自适应（扫一遍、边统计边编码）；记住"失配代价 $=D(g\|p)$"与"自适应只略差于最佳静态"。
\item \emph{ApEn 与 SampEn 的区别}：样本熵不计自身匹配，偏差更小、一致性更好。
\end{enumerate}

\textbf{四、最容易丢分}：(1) 索引从 $0$ 还是 $1$ 开始没写清（课件与教材从 $0$ 开始）；(2) 码长写成 $\lb N$ 而不取整（必须 $\lceil\lb N\rceil$）；(3) 说"随机串信息量低"——恰恰相反，随机串\emph{算法}复杂度最高；(4) 把组合信息当熵去套 Shannon 公式。
\end{exam}


\end{document}
